\documentclass[11pt,openany]{book}

% ============================================================
% Packages
% ============================================================

\usepackage[a4paper,margin=1in]{geometry}

\usepackage[T1]{fontenc}
\usepackage[utf8]{inputenc}
\usepackage{lmodern}

\usepackage{amsmath,amssymb,amsthm,mathtools}
\usepackage{bm}
\usepackage{enumitem}
\usepackage{hyperref}
\usepackage{xcolor}
\usepackage[most]{tcolorbox}

\hypersetup{
  colorlinks=true,
  linkcolor=blue!50!black,
  citecolor=blue!50!black,
  urlcolor=blue!50!black
}

% ============================================================
% Theorem-like environments
% ============================================================

\theoremstyle{definition}
\newtheorem{definition}{Definition}[chapter]

\theoremstyle{remark}
\newtheorem{remark}{Remark}[chapter]

% ============================================================
% Pedagogical boxes
% ============================================================

\newtcolorbox{lessonbox}{
  colback=blue!3,
  colframe=blue!45!black,
  title=Lesson,
  fonttitle=\bfseries,
  breakable
}

\newtcolorbox{warningbox}{
  colback=red!3,
  colframe=red!55!black,
  title=Warning,
  fonttitle=\bfseries,
  breakable
}

\newtcolorbox{intuitionbox}{
  colback=green!3,
  colframe=green!40!black,
  title=Intuition,
  fonttitle=\bfseries,
  breakable
}

\newtcolorbox{checkpointbox}{
  colback=orange!4,
  colframe=orange!65!black,
  title=Checkpoint,
  fonttitle=\bfseries,
  breakable
}

\newtcolorbox{developmentnote}{
  colback=gray!6,
  colframe=gray!60!black,
  title=Development notes,
  fonttitle=\bfseries,
  breakable
}

% ============================================================
% Macros
% ============================================================

% Spaces
\newcommand{\M}{\mathcal M}
\newcommand{\D}{\mathcal D}
\newcommand{\Pspace}{\mathcal P}
\newcommand{\B}{\mathcal B}

% Operators
\newcommand{\G}{G}
\newcommand{\T}{T}
\newcommand{\A}{A}
\newcommand{\X}{X}
\newcommand{\C}{C}
\newcommand{\CD}{C_D}
\newcommand{\Czero}{C_0}

% Bold finite-dimensional objects
\newcommand{\bd}{\mathbf d}
\newcommand{\btilded}{\widetilde{\mathbf d}}
\newcommand{\bp}{\mathbf p}
\newcommand{\btildep}{\widetilde{\mathbf p}}
\newcommand{\bv}{\mathbf v}
\newcommand{\bw}{\mathbf w}
\newcommand{\bu}{\mathbf u}
\newcommand{\bK}{\mathbf K}
\newcommand{\bG}{\mathbf G}
\newcommand{\bX}{\mathbf X}
\newcommand{\bT}{\mathbf T}
\newcommand{\bA}{\mathbf A}
\newcommand{\bC}{\mathbf C}
\newcommand{\bCD}{\mathbf C_D}
\newcommand{\bM}{\mathbf M}
\newcommand{\bW}{\mathbf W}
\newcommand{\Tau}{\mathcal{T}}
% Measures and relations
\newcommand{\Lset}{L_{\mathrm{set}}}
\newcommand{\Lprob}{L_{\mathrm{prob}}}
\newcommand{\muprior}{\mu_{\mathrm{prior}}}
\newcommand{\nudobs}{\nu_{\mathbf d_{\mathrm{obs}}}}

% Misc
\newcommand{\dd}{\,\mathrm d}
\newcommand{\range}{\operatorname{ran}}
\newcommand{\kernel}{\operatorname{ker}}
\newcommand{\spanop}{\operatorname{span}}
\newcommand{\Tr}{\operatorname{Tr}}
\newcommand{\HS}{\mathrm{HS}}
\newcommand{\argmin}{\operatorname*{arg\,min}}
\newcommand{\pushforward}{_{\#}}

% ============================================================
% Title
% ============================================================

\title{
  My Take on SOLA:\\
  Resolving Kernels, Proxy Inversions, and the Problem with Shifted Noise
}
\author{Author Name}
\date{\today}

% ============================================================
% Document
% ============================================================

\begin{document}

\frontmatter

\maketitle
\tableofcontents

% ============================================================
% Preface
% ============================================================

\chapter*{Preface}
\addcontentsline{toc}{chapter}{Preface}

\section*{What this document is about}

\begin{developmentnote}
This preface should set the tone carefully.

The document is not a polemic against SOLA. The attitude should be:
SOLA is elegant, useful, and mathematically neat. The point is not to tear it down, but to
make its interpretation more precise.

The central slogan should be introduced early:
\[
  \text{Ask with targets. Interpret with resolving kernels. Add priors when you want truth.}
\]

The reader should feel that we are being fair to SOLA before we become critical.
\end{developmentnote}

SOLA is often introduced as a method for estimating localized averages, or more general linear
properties, of an unknown model without first reconstructing the whole model. This is a
beautiful idea. Instead of solving for everything, SOLA tries to build from the data exactly the
kind of quantity one wants to report.

The main purpose of these notes is to explain that idea from the ground up, and then to examine
carefully what the resulting numbers mean.

\begin{lessonbox}
SOLA is best understood as a method for reshuffling linear data information into property-like
information. The reshuffled information is real, but it must be labelled honestly.
\end{lessonbox}

\section*{The central caution}

\begin{developmentnote}
This section should introduce the main interpretive issue without yet proving it.

Important distinction:
\[
  \text{target kernel} \neq \text{resolving kernel}.
\]

The target kernel is what we ask for. The resolving kernel is what SOLA can actually build from
the data kernels.

Do not yet discuss pathological counterexamples in detail. Save that drama for Act VII.
\end{developmentnote}

A SOLA calculation begins with target kernels. These describe the linear properties we would
like to know. But the method generally returns resolving kernels. These are the linear
properties that the available data can actually synthesize.

If the resolving kernel equals the target kernel, then the interpretation is direct. If it does not,
then the SOLA number is not the exact target property of the true model. It is the resolving-kernel
property of the true model, up to observational noise.

\begin{warningbox}
The target kernel is a design request. The resolving kernel is the answer delivered by the data.
\end{warningbox}

\section*{How to read this document}

\begin{developmentnote}
Explain the narrative structure.

The document should follow the same style as the previous blog:
start with a simple concrete problem, let the natural method develop, then inspect what it really
means.

Avoid front-loading abstract Hilbert-space notation. Start with integrals and pictures.
\end{developmentnote}

The notes are organized as a guided story. We begin with a simple integral picture, build SOLA
in its cleanest form, add normalization constraints and noise, and only later ask the more delicate
questions: what did SOLA estimate, what does its uncertainty mean, and what kind of prior
information is needed if we want statements about the true target property?

\section*{Notation and conventions}

\begin{remark}[Notation for operators and components]\label{rem:notation}
  We adopt the following typographic conventions throughout these notes.

  \begin{itemize}
    \item \textbf{Finite-dimensional objects.} Boldface letters denote finite-dimensional vectors and matrices (coefficient/representation objects). For example \(\mathbf{v}\in\mathbb R^N\) and \(\mathbf{A}\in\mathbb R^{N\times N}\).
    \item \textbf{Infinite-dimensional objects.} Non-bold (lightface) letters denote infinite-dimensional operators or functions acting on the model Hilbert space \(\mathcal M\); e.g.\ \(G:\mathcal M\to\mathcal D\), \(C_0:\mathcal M\to\mathcal M\).
    \item \textbf{Components / entries.} The components of a vector or the entries of a matrix are written with square brackets and subscripts:
    \[
      [\mathbf v]_i = v_i,\qquad [\mathbf A]_{ij} = A_{ij},
      \qquad i,j=1,\dots,N.
    \]
    When an object is represented in a particular basis \(\{\phi_k\}_{k=1}^N\) we always mean the coordinates with respect to that basis; the matrix entries depend on the chosen basis.
    \item \textbf{Adjoints and transposes.} We reserve \(^{\top}\) for the Euclidean transpose of a finite matrix and \(^{*}\) for the Hilbert adjoint of an operator. Discrete adjoints that incorporate mass/weight matrices will be written explicitly (for example, the discrete adjoint of \(\mathbf G\) with mass matrix \(\mathbf M\) and data weight \(\mathbf W\) is \(\mathbf G^*=\mathbf M^{-1}\mathbf G^{\top}\mathbf W\)).
    \item \textbf{Infinite sequences.} When working with infinite coefficient sequences we use index notation \( (\alpha_i)_{i\ge1}\) or \(\{\alpha_i\}_{i=1}^\infty\); for finite coefficient vectors we write \(\mathbf\alpha=(\alpha_1,\dots,\alpha_N)^\top\) so that \([\mathbf\alpha]_i=\alpha_i\).
  \end{itemize}
\end{remark}

\mainmatter

```latex
% ============================================================
% Act I
% ============================================================

\part{Act I: The Kernel Game}

% ------------------------------------------------------------
\chapter{What the Data Can See}
% ------------------------------------------------------------

\section{The unknown model}

We begin with an unknown function
\[
  m:\Omega\to\mathbb R,
\]
defined on a physical domain \(\Omega\). In a geophysical problem, \(\Omega\) might be an
interval of depths, a radial coordinate inside a planet, a patch of the Earth's surface, or a
three-dimensional region. To keep the first picture simple, we imagine that
\[
  \Omega=[0,1],
\]
and that \(r\in\Omega\) is a spatial coordinate.

The function \(m(r)\) represents the model. It might be a velocity perturbation, a density
variation, a temperature anomaly, a sound-speed profile, or any other physical field we would
like to learn about. For now, the exact interpretation does not matter. What matters is that
\(m\) is the object hidden behind the data.

At this stage, we will not yet be too fussy about the exact mathematical home of \(m\). Later,
we may decide that \(m\) belongs to a Hilbert space \(\M\), perhaps something like
\(L^2(\Omega)\) or a smoother Sobolev space. That choice will matter. It determines what it
means for two models to be close, which functions are allowed, and which operations are
continuous.

But we do not need that machinery yet. For the moment, the essential picture is this:

\[
  \text{there is an unknown function }m(r),\text{ and the data only see it indirectly.}
\]

\begin{warningbox}
The phrase ``an unknown function'' sounds innocent, but it is already hiding a modelling choice.
Eventually we must say what kind of function \(m\) is allowed to be. If we do not make that
choice explicitly, some later computational step will make it for us.
\end{warningbox}

The important point is that we do not observe \(m(r)\) directly. We do not get to open the
Earth, the star, the material, or the experiment and read off the function value at every \(r\).
Instead, we observe numbers produced by an experiment. Those numbers are related to \(m\),
but they are not \(m\) itself.

The rest of this chapter is about understanding what kind of relationship the data have with
the model.

\section{Sensitivity kernels}

Suppose that our experiment gives us \(N_d\) data values. Associated with these data values are
functions
\[
  K_1(r),K_2(r),\dots,K_{N_d}(r).
\]
These are called sensitivity kernels.

Each kernel \(K_i\) describes how the \(i\)th datum responds to the model. If \(K_i(r)\) is large
near some location, then the \(i\)th datum is sensitive to the model near that location. If
\(K_i(r)\) is small or zero in a region, then the \(i\)th datum hardly sees that region at all.

\begin{intuitionbox}
A sensitivity kernel is the footprint of a datum on the model. Where the kernel is large, the
datum listens. Where the kernel is small, the datum is nearly deaf.
\end{intuitionbox}

One should not imagine the datum as measuring \(m(r)\) at a single point. Usually it does
something more blurred. It averages, smears, mixes, and filters. The kernel \(K_i\) tells us the
shape of that filtering.

For example, if \(K_i\) is sharply concentrated near \(r_0\), then the datum is mostly local to
\(r_0\). If \(K_i\) is broad, then the datum mixes information from a wide region. If \(K_i\)
oscillates between positive and negative values, then the datum compares different parts of the
model against each other.

This is already the beginning of the inverse problem. The data do not arrive as tiny labels
attached to the model. They arrive as integrals against kernels.

\begin{figure}[ht]
  \centering
  \fbox{
    \begin{minipage}{0.82\textwidth}
      \vspace{0.8cm}
      \centering
      Placeholder for a plot of many sensitivity kernels \(K_i(r)\), perhaps drawn as translucent
      curves over the same interval \(\Omega\).
      \vspace{0.8cm}
    \end{minipage}
  }
  \caption{A family of sensitivity kernels. Each datum sees the model through one of these
  kernels.}
  \label{fig:act1-sensitivity-kernels}
\end{figure}

The visual lesson of Figure~\ref{fig:act1-sensitivity-kernels} should be slightly unsettling.
Even if we have many data, each datum sees the model through its own distorted window. The
inverse problem is not a matter of simply reading off \(m\). It is a matter of understanding what
can be reconstructed, estimated, or summarized from these windows.

\section{Data as weighted integrals}

In the simplest noiseless setting that usually appears in SOLA applications, the \(i\)th datum is
\[
  d_i
  =
  \int_\Omega K_i(r)m(r)\dd r,
  \qquad i=1,\dots,N_d.
\]
This equation says that the datum \(d_i\) is obtained by multiplying the model \(m\) by the
kernel \(K_i\), then integrating over the domain.

Thus each datum is a weighted integral of the model. The model \(m\) supplies the unknown
field. The kernel \(K_i\) supplies the weights. The integral collapses the result to one number.

We collect the data into a finite-dimensional vector
\[
  \mathbf d
  =
  \begin{pmatrix}
    d_1\\
    \vdots\\
    d_{N_d}
  \end{pmatrix}
  \in\mathbb R^{N_d},
\]
so that
\[
  [\mathbf d]_i=d_i.
\]

This is the first place where the finite-dimensional and infinite-dimensional worlds meet. The
model \(m\) is a function. The data vector \(\mathbf d\) is just a list of numbers. The experiment
turns one into the other.

Later, once we choose a model Hilbert space \(\M\), the same equations can be written as
\[
  [G(m)]_i
  =
  \langle K_i,m\rangle_{\M},
  \qquad i=1,\dots,N_d,
\]
where
\[
  G:\M\to\D
\]
is the forward operator, and \(\D\simeq\mathbb R^{N_d}\) is the data space.

For now, this notation should be read as a compressed version of the integral equations:
\[
  G(m)
  =
  \begin{pmatrix}
    \int_\Omega K_1(r)m(r)\dd r\\
    \vdots\\
    \int_\Omega K_{N_d}(r)m(r)\dd r
  \end{pmatrix}.
\]

\begin{remark}[Why the data are linear]
Each datum depends linearly on \(m\). If \(m_1\) and \(m_2\) are two models and
\(\alpha,\beta\in\mathbb R\), then
\[
  \int_\Omega K_i(r)\bigl(\alpha m_1(r)+\beta m_2(r)\bigr)\dd r
  =
  \alpha\int_\Omega K_i(r)m_1(r)\dd r
  +
  \beta\int_\Omega K_i(r)m_2(r)\dd r.
\]
Therefore
\[
  G(\alpha m_1+\beta m_2)=\alpha G(m_1)+\beta G(m_2).
\]
This linearity is what makes SOLA possible in its usual form.
\end{remark}

It is useful to pause here. The data are not arbitrary numbers. They are responses of the model
against known kernels. If we know the kernels, then we know precisely what kind of information
the data contain.

But we also know what they do not contain.

\section{What finite data cannot see}

The map \(G\) sends a function to finitely many numbers:
\[
  G:\M\to\D,
  \qquad
  \D\simeq\mathbb R^{N_d}.
\]
This is a compression. Unless the model space is itself finite-dimensional in a very special way,
many different models will produce the same data.

A model perturbation \(h\in\M\) is invisible to the data if
\[
  G(h)=0.
\]
In integral form, this means
\[
  \int_\Omega K_i(r)h(r)\dd r=0,
  \qquad i=1,\dots,N_d.
\]
Such an \(h\) may be large, oscillatory, smooth, rough, physically plausible, or physically
implausible. The data do not care. If \(G(h)=0\), then adding \(h\) to a model does not change
the noiseless data:
\[
  G(m+h)=G(m)+G(h)=G(m).
\]

The collection of all such invisible perturbations is the nullspace
\[
  \kernel(G)
  =
  \{h\in\M:G(h)=0\}.
\]

\begin{intuitionbox}
If \(h\in\kernel(G)\), then \(h\) is written in invisible ink as far as these data are concerned.
It may be present in the model, but the experiment cannot read it.
\end{intuitionbox}

This is not a defect of SOLA. It is a structural fact about inverse problems with indirect,
finite data. The data only constrain the parts of the model that interact with the sensitivity
kernels. Everything orthogonal to those measurements is left unresolved.

This observation will matter later. For now, it tells us that asking for the whole model \(m\)
may be too ambitious. The data do not contain the whole model. They contain only finitely many
linear shadows of it.

\begin{checkpointbox}
The data do not show us the model. They show us finitely many shadows of the model.
\end{checkpointbox}

That is the first move in the SOLA story. If the data cannot reveal the whole model, perhaps we
should not ask them to. Perhaps we should ask for something smaller, sharper, and better matched
to what the data can actually say.

% ------------------------------------------------------------
\chapter{What We Would Like to Know}
% ------------------------------------------------------------

\section{The target property}

In many inverse problems, the goal is not really to reconstruct the whole model. The whole model
may be too much to ask for, and in many applications it is not even the thing we truly want.

Instead, we may want a particular property of the model.

For example, we may want the average value of \(m\) near a certain location. We may want the
contrast between two regions. We may want a smoothed profile. We may want a coefficient in
some basis. We may want a derivative-like quantity. The details vary, but the common structure
is that we want to ask the model a specific question.

The simplest such question has the form
\[
  p
  =
  \int_\Omega T(r)m(r)\dd r.
\]
Here \(T(r)\) is the target kernel, and \(p\in\mathbb R\) is the property we would like to know.

This looks very similar to the data equations
\[
  d_i
  =
  \int_\Omega K_i(r)m(r)\dd r.
\]
The difference is that the kernels \(K_i\) are given to us by the experiment, while the target
kernel \(T\) is chosen by us. The data kernels say what the experiment measured. The target
kernel says what we wish it had measured.

\begin{intuitionbox}
The sensitivity kernels are the questions the experiment actually asked. The target kernel is the
question we wish we could ask.
\end{intuitionbox}

At this point, the SOLA idea begins to appear. If the data are integrals against kernels, and the
desired property is also an integral against a kernel, then perhaps we can combine the data
integrals to imitate the desired one.

But before we combine anything, we should understand what kind of target we are asking for.

\section{Localized averages}

A common target is a localized average. We choose \(T\) to be concentrated near some location
\(r_0\). For instance, \(T\) might be a smooth bump centred at \(r_0\), small away from \(r_0\),
and normalized so that
\[
  \int_\Omega T(r)\dd r=1.
\]
Then
\[
  p
  =
  \int_\Omega T(r)m(r)\dd r
\]
is a weighted average of \(m\) near \(r_0\).

The normalization matters. If \(m\) is constant, say \(m(r)=c\), then
\[
  p
  =
  \int_\Omega T(r)c\dd r
  =
  c\int_\Omega T(r)\dd r
  =
  c.
\]
So the target kernel reproduces constants correctly. That is exactly what an averaging kernel
should do.

\begin{figure}[ht]
  \centering
  \fbox{
    \begin{minipage}{0.82\textwidth}
      \vspace{0.8cm}
      \centering
      Placeholder for a localized target kernel \(T(r)\), shown as a smooth bump with unit area.
      Optionally overlay a model \(m(r)\) in the background.
      \vspace{0.8cm}
    \end{minipage}
  }
  \caption{A localized averaging target. The target kernel defines the property we would like to
  estimate.}
  \label{fig:act1-target-kernel}
\end{figure}

The location and width of \(T\) encode the question. A narrow target asks for a highly localized
average. A broad target asks for a more smeared average. A target with oscillations asks for a
contrast or derivative-like feature rather than an ordinary average.

For this first act, however, we stay with localized averages. They are the cleanest way to see
what SOLA is trying to do.

\begin{lessonbox}
A target kernel is a question written as a function. It asks the model for a weighted answer.
\end{lessonbox}

\section{The property map}

For a single target kernel \(T\), the property is a scalar:
\[
  p=\int_\Omega T(r)m(r)\dd r.
\]
In practice, we often want several properties. For example, we might choose target kernels
\[
  T^{(1)},T^{(2)},\dots,T^{(N_p)},
\]
centred at different locations. Each one defines a property
\[
  p_k
  =
  \int_\Omega T^{(k)}(r)m(r)\dd r,
  \qquad k=1,\dots,N_p.
\]
Collecting these into a vector gives
\[
  \mathbf p
  =
  \begin{pmatrix}
    p_1\\
    \vdots\\
    p_{N_p}
  \end{pmatrix}
  \in\mathbb R^{N_p}.
\]

Once we have chosen a model Hilbert space \(\M\), this collection of properties can be written
as an operator
\[
  T:\M\to\Pspace,
\]
where \(\Pspace\simeq\mathbb R^{N_p}\), and
\[
  [T(m)]_k
  =
  \langle T^{(k)},m\rangle_{\M},
  \qquad k=1,\dots,N_p.
\]

The notation is compact, but the meaning is simple. The operator \(T\) asks the model
\(N_p\) linear questions. The \(k\)th question is encoded by the kernel \(T^{(k)}\).

We now have two maps:
\[
  G:\M\to\D,
  \qquad
  T:\M\to\Pspace.
\]
The first map \(G\) describes what the experiment gives us. The second map \(T\) describes what
we would like to know.

In an ideal world, the data would directly contain \(T(m)\). But the experiment gives us
\(G(m)\), not \(T(m)\). SOLA lives in the gap between these two maps.

\begin{checkpointbox}
The forward map \(G\) describes the questions the experiment actually answered. The property
map \(T\) describes the questions we wish had been answered.
\end{checkpointbox}

\section{The SOLA question}

We can now state the central SOLA question.

The data give us integrals against sensitivity kernels:
\[
  d_i
  =
  \int_\Omega K_i(r)m(r)\dd r.
\]
The desired property is an integral against a target kernel:
\[
  p
  =
  \int_\Omega T(r)m(r)\dd r.
\]
So SOLA asks:

\[
  \text{Can we combine the available data so that the combined kernel behaves like }T?
\]

More explicitly, can we find weights \(x_1,\dots,x_{N_d}\) such that
\[
  x_1K_1(r)+x_2K_2(r)+\cdots+x_{N_d}K_{N_d}(r)
\]
looks like the target kernel \(T(r)\)?

If we can, then the corresponding weighted data combination
\[
  x_1d_1+x_2d_2+\cdots+x_{N_d}d_{N_d}
\]
will behave like the desired property.

This is the whole spark of SOLA. We do not first reconstruct \(m\). We try to synthesize the
property directly from the data.

\begin{lessonbox}
SOLA begins by refusing to reconstruct everything. It tries to manufacture the property directly
from the data.
\end{lessonbox}

The next act turns this idea into algebra. The key observation will be almost embarrassingly
simple: because the data are linear integrals, combining data is the same as combining kernels.
That is where the kernel game truly begins.
```


```latex
% ============================================================
% Act II
% ============================================================

\part{Act II: Building SOLA from Scratch}

% ------------------------------------------------------------
\chapter{Combining Data Means Combining Kernels}
% ------------------------------------------------------------

\section{A weighted sum of data}

We now make the first SOLA move.

The data are
\[
  [\mathbf d]_i=d_i
  =
  \int_\Omega K_i(r)m(r)\dd r,
  \qquad i=1,\dots,N_d.
\]
Each datum is one weighted integral of the model. Suppose we choose scalar weights
\[
  x_1,\dots,x_{N_d}\in\mathbb R.
\]
Equivalently, we collect them into the finite-dimensional vector
\[
  \mathbf x
  =
  (x_1,\dots,x_{N_d})^\top
  \in\mathbb R^{N_d},
  \qquad
  [\mathbf x]_i=x_i.
\]

Now form the weighted sum of data
\[
  \sum_{i=1}^{N_d}x_i d_i.
\]
Substituting the integral expression for each datum gives
\[
  \sum_{i=1}^{N_d}x_i d_i
  =
  \sum_{i=1}^{N_d}
  x_i
  \int_\Omega K_i(r)m(r)\dd r.
\]
At first glance, this is only a weighted sum of numbers. But because each number came from
an integral, the weights can be moved inside the integral:
\[
  \sum_{i=1}^{N_d}x_i d_i
  =
  \int_\Omega
  \left(
    \sum_{i=1}^{N_d}x_iK_i(r)
  \right)
  m(r)\dd r.
\]

This is the small algebraic hinge on which the whole SOLA door swings.

A weighted sum of data is not just a weighted sum of numbers. It is itself another weighted
integral of the model. The new weight is obtained by taking the same weighted sum of the
sensitivity kernels.

\begin{intuitionbox}
Combining data combines the eyes of the experiment. The new datum sees the model through
the combined kernel
\[
  x_1K_1+\cdots+x_{N_d}K_{N_d}.
\]
\end{intuitionbox}

This explains why SOLA is a natural method in linear inverse problems. The data are linear
questions asked of the model. If we combine those questions linearly, we obtain another linear
question.

The hope is that this new question can be made to resemble the question we actually wanted
to ask.

\section{The resolving kernel}

The kernel
\[
  R_{\mathbf x}(r)
  :=
  \sum_{i=1}^{N_d}x_iK_i(r)
\]
is called the resolving kernel associated with the weights \(\mathbf x\).

With this notation, the weighted data sum becomes
\[
  \sum_{i=1}^{N_d}x_i d_i
  =
  \int_\Omega R_{\mathbf x}(r)m(r)\dd r.
\]
So the weighted data sum is the model \(m\) tested against the resolving kernel
\(R_{\mathbf x}\).

This is the first important change of perspective. We are no longer thinking only about data
values. We are thinking about what kernel those data values correspond to.

\begin{checkpointbox}
Choosing data weights is the same thing as choosing a kernel from the span of the sensitivity
kernels.
\end{checkpointbox}

In this scalar setting, the SOLA problem becomes almost visible by inspection. We have a target
kernel \(\Tau^{(1)}\), defining the desired property
\[
  [\mathbf p]_1
  =
  \int_\Omega \Tau^{(1)}(r)m(r)\dd r.
\]
We have a resolving kernel
\[
  R_{\mathbf x}
  =
  \sum_{i=1}^{N_d}x_iK_i,
\]
which defines the property-like quantity actually built from the data:
\[
  \sum_{i=1}^{N_d}x_i d_i
  =
  \int_\Omega R_{\mathbf x}(r)m(r)\dd r.
\]
SOLA asks us to choose \(\mathbf x\) so that
\[
  R_{\mathbf x}\approx \Tau^{(1)}.
\]

\begin{figure}[ht]
  \centering
  \fbox{
    \begin{minipage}{0.82\textwidth}
      \vspace{0.8cm}
      \centering
      Placeholder for a visual showing several sensitivity kernels \(K_i\), chosen weights
      \(x_i\), and the resulting resolving kernel \(R_{\mathbf x}\) overlaid with the target
      kernel \(\Tau^{(1)}\).
      \vspace{0.8cm}
    \end{minipage}
  }
  \caption{A weighted combination of sensitivity kernels produces a resolving kernel. SOLA
  chooses the weights so that this resolving kernel resembles the target kernel.}
  \label{fig:act2-resolving-kernel}
\end{figure}

The target kernel is the kernel we asked for. The resolving kernel is the kernel the data can
actually synthesize using the chosen weights.

For the moment, we do not judge this difference. We simply name it.

\section{Why the combination should be linear}

It is worth asking why SOLA searches for linear combinations of data in the first place. Why not
take a square, a product, a ratio, or some other nonlinear expression?

The answer is that the quantities we are trying to estimate are linear functionals of the model.
For example,
\[
  [\mathbf p]_1
  =
  \int_\Omega \Tau^{(1)}(r)m(r)\dd r
\]
is linear in \(m\). Likewise, each datum
\[
  d_i
  =
  \int_\Omega K_i(r)m(r)\dd r
\]
is linear in \(m\).

A linear combination of linear functionals is still a linear functional:
\[
  m
  \mapsto
  \sum_{i=1}^{N_d}x_i
  \int_\Omega K_i(r)m(r)\dd r
\]
is the same as
\[
  m
  \mapsto
  \int_\Omega R_{\mathbf x}(r)m(r)\dd r.
\]
Thus the weighted data sum remains the same kind of mathematical object as the target property.

A nonlinear combination would usually break this structure. For example,
\[
  d_1^2+d_2^2
\]
is generally not a linear functional of \(m\). It is quadratic in the data, and therefore usually
quadratic in the model. It cannot usually be represented as
\[
  \int_\Omega Q(r)m(r)\dd r
\]
for some kernel \(Q\).

This does not mean nonlinear methods are impossible or useless. It only means that they would
be solving a different kind of problem. SOLA is designed for the case where both the data and
the desired properties are linear questions about the model.

\begin{lessonbox}
SOLA uses linear combinations because it wants to turn linear data functionals into another
linear functional. The algebra is not an accident; it is the method.
\end{lessonbox}

There is another practical reason this matters. Once the data combination is linear, the effect of
data noise can be propagated cleanly. If the data covariance is \(\mathbf C_D\), and the SOLA
weights are collected into a matrix \(\mathbf X\), then the propagated covariance later takes the
form
\[
  \mathbf C_P=\mathbf X\mathbf C_D\mathbf X^*.
\]
This simple covariance propagation is one of the reasons SOLA remains computationally
attractive.

For now, however, we stay in the noiseless picture. The key object is still the resolving kernel.

\section{The span limitation}

The resolving kernel has the form
\[
  R_{\mathbf x}
  =
  \sum_{i=1}^{N_d}x_iK_i.
\]
Therefore it must belong to the finite-dimensional span
\[
  R_{\mathbf x}\in
  \spanop\{K_1,\dots,K_{N_d}\}.
\]

This is not a flaw. It is the geometry of the problem.

The data give us the kernels \(K_1,\dots,K_{N_d}\). By taking linear combinations of the data,
we can only create linear combinations of those kernels. If the target kernel \(\Tau^{(1)}\) lies in
this span, then there exist weights \(\mathbf x\) such that
\[
  R_{\mathbf x}=\Tau^{(1)}.
\]
In that special case, SOLA can reproduce the target kernel exactly.

But if
\[
  \Tau^{(1)}
  \notin
  \spanop\{K_1,\dots,K_{N_d}\},
\]
then no choice of weights can make the resolving kernel identical to the target kernel. The best
we can do is approximate it.

\begin{warningbox}
SOLA cannot synthesize a kernel that is not available through the span of the sensitivity kernels.
The method can ask for a target, but the data decide which resolving kernels are actually
reachable.
\end{warningbox}

This limitation should not be read too negatively. In practice, this is exactly why SOLA is useful:
it makes the mismatch visible. Instead of hiding the resolution of the estimate inside a black-box
model reconstruction, SOLA asks us to look directly at the resolving kernels.

The next chapter turns this geometric idea into an optimization problem.

% ------------------------------------------------------------
\chapter{The Noiseless, Unconstrained SOLA Problem}
% ------------------------------------------------------------

\section{Matching the target kernel}

We now formalize the simplest possible SOLA problem.

There is one target kernel \(\Tau^{(1)}\), and we look for weights
\[
  \mathbf x=(x_1,\dots,x_{N_d})^\top
\]
such that
\[
  R_{\mathbf x}
  =
  \sum_{i=1}^{N_d}x_iK_i
\]
is close to \(\Tau^{(1)}\).

To say ``close'', we need a way of measuring the size of the mismatch
\[
  \Tau^{(1)}-R_{\mathbf x}.
\]
For this first presentation, imagine that all kernels live in the same Hilbert space \(\M\), and
that closeness is measured using the norm \(\|\cdot\|_{\M}\). Then the simplest SOLA problem is
\[
  \mathbf x_0
  =
  \argmin_{\mathbf x\in\mathbb R^{N_d}}
  \left\|
    \Tau^{(1)}-R_{\mathbf x}
  \right\|_{\M}^2.
\]

In words, we choose the weights whose resolving kernel best approximates the target kernel.

\begin{remark}[A small functional-analytic caveat]
In a fully careful treatment, the sensitivity kernels and target kernels should be placed in the
correct dual or Riesz-represented Hilbert setting. For the present explanation, we identify linear
functionals with their representing kernels. This keeps the picture readable while preserving the
main idea.
\end{remark}

The solution is a projection problem. We are projecting the target kernel onto the space spanned
by the sensitivity kernels:
\[
  \spanop\{K_1,\dots,K_{N_d}\}.
\]
If \(\Tau^{(1)}\) already lies in that span, the projection is exact. Otherwise, the projection gives
the reachable kernel closest to the requested one.

\begin{intuitionbox}
The target kernel is a destination. The span of the sensitivity kernels is the road network. SOLA
chooses the reachable point closest to the destination.
\end{intuitionbox}

This is the cleanest SOLA idea. There is no noise, no covariance, no constraint, and no
regularization parameter. There is only the geometry of kernels.

\section{Operator form}

The scalar kernel picture is enough to understand the method, but it becomes cumbersome when
we want several properties at once. So we now introduce the operator notation.

The forward map is
\[
  G:\M\to\D,
\]
where \(\D\simeq\mathbb R^{N_d}\). Its components are
\[
  [G(m)]_i
  =
  \langle K_i,m\rangle_{\M},
  \qquad i=1,\dots,N_d.
\]

The property map is
\[
  \Tau:\M\to\Pspace,
\]
where \(\Pspace\simeq\mathbb R^{N_p}\). Its components are
\[
  [\Tau(m)]_k
  =
  \langle \Tau^{(k)},m\rangle_{\M},
  \qquad k=1,\dots,N_p.
\]
For a given model \(m\), the property vector is
\[
  \mathbf p=\Tau(m)\in\Pspace.
\]

SOLA now seeks a finite-dimensional matrix
\[
  \mathbf X:\D\to\Pspace
\]
such that
\[
  \mathbf XG\approx \Tau.
\]
The composite map
\[
  A:=\mathbf XG:\M\to\Pspace
\]
is the approximate property map built from the data.

Thus, for data
\[
  \mathbf d=G(m),
\]
the SOLA estimate is
\[
  \widetilde{\mathbf p}
  =
  \mathbf X\mathbf d.
\]
In the noiseless idealization, this is
\[
  \widetilde{\mathbf p}
  =
  \mathbf XG(m)
  =
  A(m).
\]

The target map is \(\Tau\). The map actually built from data is \(A=\mathbf XG\). The SOLA
objective is to make these two maps close.

The noiseless unconstrained SOLA problem is therefore
\[
  \mathbf X_0
  =
  \argmin_{\mathbf X}
  \left\|
    \Tau-\mathbf XG
  \right\|_{\HS}^2.
\]

Here \(\|\cdot\|_{\HS}\) denotes the Hilbert-Schmidt norm. In this finite-property setting, it can
be understood as summing the squared mismatch over all target components. If the \(k\)th
component of \(\mathbf XG\) has resolving kernel \(A^{(k)}\), then the objective is essentially
measuring
\[
  \sum_{k=1}^{N_p}
  \left\|
    \Tau^{(k)}-A^{(k)}
  \right\|_{\M}^2,
\]
up to the precise Hilbert-space identifications being used.

\begin{checkpointbox}
The scalar problem asks for one resolving kernel close to one target kernel. The operator problem
asks for a whole family of resolving kernels close to a whole family of target kernels.
\end{checkpointbox}

\section{The normal equation}

We now derive the solution of the unconstrained problem
\[
  \mathbf X_0
  =
  \argmin_{\mathbf X}
  \left\|
    \Tau-\mathbf XG
  \right\|_{\HS}^2.
\]
Define
\[
  f(\mathbf X)
  :=
  \left\|
    \Tau-\mathbf XG
  \right\|_{\HS}^2.
\]
Perturb \(\mathbf X\) by a finite-dimensional matrix \(\mathbf H:\D\to\Pspace\). Then
\[
  f(\mathbf X+\mathbf H)
  =
  \left\|
    \Tau-(\mathbf X+\mathbf H)G
  \right\|_{\HS}^2.
\]
Expanding,
\[
  f(\mathbf X+\mathbf H)
  =
  \left\|
    \Tau-\mathbf XG-\mathbf HG
  \right\|_{\HS}^2.
\]
Keeping only the terms linear in \(\mathbf H\), the first variation is
\[
  \delta f
  =
  -2
  \left\langle
    \Tau-\mathbf XG,
    \mathbf HG
  \right\rangle_{\HS}.
\]
Using the Hilbert-Schmidt adjoint relation, this can be written as
\[
  \delta f
  =
  -2
  \left\langle
    (\Tau-\mathbf XG)G^*,
    \mathbf H
  \right\rangle_{\HS}.
\]
Since this must vanish for every perturbation \(\mathbf H\), the stationarity condition is
\[
  (\Tau-\mathbf XG)G^*=0.
\]
Equivalently,
\[
  \mathbf XGG^*=\Tau G^*.
\]

If \(GG^*:\D\to\D\) is invertible, then the solution is
\[
  \mathbf X_0
  =
  \Tau G^*(GG^*)^{-1}.
\]

\begin{remark}[What is \(GG^*\)?]
The operator \(G^*:\D\to\M\) maps data-space vectors back into the model space. The
composition
\[
  GG^*:\D\to\D
\]
is therefore a data-space operator. Since \(\D\) is finite-dimensional, it is represented by a
finite matrix. In components, it contains the inner products between sensitivity kernels:
\[
  [\mathbf{GG^*}]_{ij}
  =
  \langle K_j,K_i\rangle_{\M},
  \qquad i,j=1,\dots,N_d,
\]
up to the chosen data-space inner product convention.
\end{remark}

The expression
\[
  \mathbf X_0
  =
  \Tau G^*(GG^*)^{-1}
\]
has a simple interpretation. The factor
\[
  G^*(GG^*)^{-1}
\]
back-projects data into a particular model-space representative. Then \(\Tau\) evaluates the
target properties of that representative.

But SOLA does not need to be introduced as model reconstruction followed by property
evaluation. It can be introduced more directly: \(\mathbf X_0\) is the finite-dimensional data
combination that makes \(\mathbf XG\) as close as possible to \(\Tau\).

\section{What this first solution means}

The operator
\[
  \mathbf X_0
  =
  \Tau G^*(GG^*)^{-1}
\]
is the best unconstrained linear reshuffling of the data, measured by the Hilbert-Schmidt
mismatch between the target property map \(\Tau\) and the approximate property map
\[
  A_0:=\mathbf X_0G.
\]

For each component \(k=1,\dots,N_p\), the \(k\)th row of \(\mathbf X_0\) gives weights
\[
  [\mathbf X_0]_{k1},\dots,[\mathbf X_0]_{kN_d}.
\]
These weights define the \(k\)th resolving kernel
\[
  A_0^{(k)}
  =
  \sum_{i=1}^{N_d}
  [\mathbf X_0]_{ki}K_i.
\]
The \(k\)th SOLA output is therefore
\[
  [\widetilde{\mathbf p}]_k
  =
  \sum_{i=1}^{N_d}
  [\mathbf X_0]_{ki}d_i
  =
  \int_\Omega A_0^{(k)}(r)m(r)\dd r.
\]

This is the object SOLA has actually built: not a direct measurement of the requested target
kernel \(\Tau^{(k)}\), but a direct measurement of the resolving kernel \(A_0^{(k)}\).

In the best case,
\[
  A_0^{(k)}=\Tau^{(k)}.
\]
Then the SOLA output equals the desired target property. More commonly,
\[
  A_0^{(k)}\approx \Tau^{(k)}.
\]
Then the SOLA output should be read through the resolving kernel itself.

For now, we simply note the distinction. The whole point of SOLA is that it gives us these
resolving kernels explicitly. They are not hidden inside the method. They are part of the
deliverable.

\begin{lessonbox}
At this stage, SOLA is a best-approximation problem: find the property-like map that can be
built from the available data map.
\end{lessonbox}

The next act adds an important calibration condition. If the target kernels are supposed to
represent averages, then the resolving kernels should behave like averages too. Shape alone is
not enough. The kernels also need the right mass.
```


```latex
% ============================================================
% Act III
% ============================================================

\part{Act III: Averages Need Mass}

% ------------------------------------------------------------
\chapter{Why Area One Matters}
% ------------------------------------------------------------

\section{Averages are calibrated linear functionals}

In the previous act, we saw the basic SOLA mechanism. A weighted sum of data produces a
resolving kernel:
\[
  R_{\mathbf x}(r)
  =
  \sum_{i=1}^{N_d}x_iK_i(r),
\]
and the corresponding data combination is
\[
  \sum_{i=1}^{N_d}x_id_i
  =
  \int_\Omega R_{\mathbf x}(r)m(r)\dd r.
\]
The natural goal was to choose \(\mathbf x\) so that \(R_{\mathbf x}\) resembles the target kernel
\(\mathcal T^{(1)}\).

But if the target is meant to represent an average, shape is not the whole story. An average must
also have the right scale.

Suppose \(\mathcal T^{(1)}\) is a localized averaging kernel. Then it should satisfy
\[
  \int_\Omega \mathcal T^{(1)}(r)\dd r=1.
\]
This condition makes the associated property
\[
  [\mathbf p]_1
  =
  \int_\Omega \mathcal T^{(1)}(r)m(r)\dd r
\]
a genuine weighted average.

To see why, test the target against a constant model. If
\[
  m(r)=c,
\]
then
\[
  [\mathbf p]_1
  =
  \int_\Omega \mathcal T^{(1)}(r)c\dd r
  =
  c\int_\Omega \mathcal T^{(1)}(r)\dd r
  =
  c.
\]
So a unit-mass target kernel reproduces constants correctly. That is what an average should do.

If instead
\[
  \int_\Omega \mathcal T^{(1)}(r)\dd r=2,
\]
then the same constant model would give
\[
  [\mathbf p]_1=2c.
\]
The kernel may still be localized, smooth, and visually reasonable, but it is no longer an
averaging kernel. It has the wrong calibration.

\begin{intuitionbox}
Averaging is not only about where the kernel is located. It is also about how much total weight
the kernel carries. The location tells us where we are looking. The mass tells us whether the
answer is on the right scale.
\end{intuitionbox}

For several localized averages, we choose target kernels
\[
  \mathcal T^{(1)},\dots,\mathcal T^{(N_p)}
\]
with
\[
  \int_\Omega \mathcal T^{(k)}(r)\dd r=1,
  \qquad k=1,\dots,N_p.
\]
The property vector is then
\[
  \mathbf p=\Tau(m)\in\Pspace,
\]
with components
\[
  [\mathbf p]_k
  =
  \int_\Omega \mathcal T^{(k)}(r)m(r)\dd r,
  \qquad k=1,\dots,N_p.
\]
Each component is intended to be a localized average of the same model, perhaps centred at a
different location.

\begin{lessonbox}
A localized averaging kernel must do two things: it must localize the model, and it must
reproduce constants correctly. The second requirement is the unit-mass condition.
\end{lessonbox}

\section{The resolving kernel may have the wrong mass}

Now return to the resolving kernel
\[
  R_{\mathbf x}(r)
  =
  \sum_{i=1}^{N_d}x_iK_i(r).
\]
Even if \(R_{\mathbf x}\) looks similar to the target kernel \(\mathcal T^{(1)}\), it need not have
the same total mass.

Its integral is
\[
  \int_\Omega R_{\mathbf x}(r)\dd r
  =
  \int_\Omega
  \sum_{i=1}^{N_d}x_iK_i(r)\dd r.
\]
By linearity,
\[
  \int_\Omega R_{\mathbf x}(r)\dd r
  =
  \sum_{i=1}^{N_d}
  x_i
  \int_\Omega K_i(r)\dd r.
\]
So the mass of the resolving kernel is controlled by the same data weights \(x_i\), but through
the integrals of the sensitivity kernels.

There is no automatic reason for this quantity to equal one:
\[
  \sum_{i=1}^{N_d}
  x_i
  \int_\Omega K_i(r)\dd r
  \neq 1
\]
in general.

This matters because the SOLA output associated with \(R_{\mathbf x}\) is
\[
  \sum_{i=1}^{N_d}x_id_i
  =
  \int_\Omega R_{\mathbf x}(r)m(r)\dd r.
\]
If \(R_{\mathbf x}\) is supposed to be interpreted as an averaging kernel, then it should act
correctly on constant models. For \(m(r)=c\),
\[
  \int_\Omega R_{\mathbf x}(r)c\dd r
  =
  c\int_\Omega R_{\mathbf x}(r)\dd r.
\]
Thus the result equals \(c\) only if
\[
  \int_\Omega R_{\mathbf x}(r)\dd r=1.
\]

\begin{warningbox}
A kernel can have a plausible shape but the wrong total mass. Then it is not the average we
meant to ask for.
\end{warningbox}

This is an easy issue to miss visually. Two kernels may appear to sit in the same region, have
similar widths, and even have similar oscillations, while still carrying different total mass. A
plot can make them look like siblings. The constant-model test reveals whether they are doing
the same averaging job.

\begin{figure}[ht]
  \centering
  \fbox{
    \begin{minipage}{0.82\textwidth}
      \vspace{0.8cm}
      \centering
      Placeholder for a figure showing two kernels with similar shape but different total area.
      The caption should emphasize that visual localization alone does not guarantee correct
      averaging.
      \vspace{0.8cm}
    \end{minipage}
  }
  \caption{Two kernels can look similarly localized while having different total mass. For
  averages, the mass is part of the meaning.}
  \label{fig:act3-wrong-mass}
\end{figure}

For multiple properties, the same issue appears component by component. If the approximate
property map is
\[
  A:=\mathbf XG:\M\to\Pspace,
\]
then the \(k\)th resolving kernel is
\[
  A^{(k)}(r)
  =
  \sum_{i=1}^{N_d}[\mathbf X]_{ki}K_i(r).
\]
For this resolving kernel to define an average, we need
\[
  \int_\Omega A^{(k)}(r)\dd r=1,
  \qquad k=1,\dots,N_p.
\]

\section{Unimodularity}

The condition
\[
  \int_\Omega R_{\mathbf x}(r)\dd r=1
\]
is often called a unimodularity condition. The word sounds slightly ornate, but the idea is
simple: the resolving kernel should have unit mass.

For one average, unimodularity says
\[
  \int_\Omega
  \sum_{i=1}^{N_d}x_iK_i(r)\dd r
  =
  1.
\]
Equivalently,
\[
  \sum_{i=1}^{N_d}
  x_i
  \int_\Omega K_i(r)\dd r
  =
  1.
\]

For many averages, each resolving kernel should have unit mass:
\[
  \int_\Omega A^{(k)}(r)\dd r=1,
  \qquad k=1,\dots,N_p.
\]

\begin{intuitionbox}
Unimodularity is the constant-model test. If the true model is flat, an averaging kernel should
return the flat value. Unit mass is what makes that happen.
\end{intuitionbox}

This condition does not guarantee that the resolving kernel equals the target kernel. It does not
guarantee perfect localization. It does not remove the limitations imposed by the span of the
sensitivity kernels.

But it does enforce a basic calibration. If the output is going to be called an average, then the
resolving kernel should at least behave like an average on constants.

\begin{checkpointbox}
In Act II we asked for shape: make the resolving kernel look like the target. In Act III we add
calibration: make the resolving kernel carry the right total mass.
\end{checkpointbox}

% ------------------------------------------------------------
\chapter{Constrained SOLA for Averages}
% ------------------------------------------------------------

\section{The average vector}

We now translate unimodularity into operator notation.

Let \(1_\Omega\in\M\) denote the constant function with value one:
\[
  1_\Omega(r)=1,
  \qquad r\in\Omega.
\]
Applying the forward map \(G\) to this constant function gives a finite-dimensional vector in the
data space:
\[
  \mathbf v:=G(1_\Omega)\in\D.
\]
Its components are
\[
  [\mathbf v]_i
  =
  \int_\Omega K_i(r)\dd r,
  \qquad i=1,\dots,N_d.
\]
So \(\mathbf v\) records the total mass of each sensitivity kernel.

This vector is useful because the mass of any resolving kernel can be written using it. For one
set of weights \(\mathbf x\),
\[
  R_{\mathbf x}(r)
  =
  \sum_{i=1}^{N_d}x_iK_i(r),
\]
and therefore
\[
  \int_\Omega R_{\mathbf x}(r)\dd r
  =
  \sum_{i=1}^{N_d}x_i[\mathbf v]_i.
\]
In Euclidean notation this is
\[
  \int_\Omega R_{\mathbf x}(r)\dd r
  =
  \mathbf x^\top\mathbf v.
\]

For multiple properties, the SOLA matrix \(\mathbf X:\D\to\Pspace\) has rows containing the
weights for each property. The \(k\)th resolving kernel is
\[
  A^{(k)}(r)
  =
  \sum_{i=1}^{N_d}[\mathbf X]_{ki}K_i(r),
\]
so its mass is
\[
  \int_\Omega A^{(k)}(r)\dd r
  =
  \sum_{i=1}^{N_d}[\mathbf X]_{ki}[\mathbf v]_i.
\]
But the right-hand side is exactly the \(k\)th component of \(\mathbf X\mathbf v\):
\[
  \int_\Omega A^{(k)}(r)\dd r
  =
  [\mathbf X\mathbf v]_k.
\]

\begin{lessonbox}
The vector \(\mathbf v=G(1_\Omega)\) lets us test the total mass of resolving kernels using only
data-space algebra.
\end{lessonbox}

This is a pleasant little piece of bookkeeping. Instead of writing one integral constraint for each
resolving kernel, we can write a single finite-dimensional vector constraint.

\section{The constraint \texorpdfstring{\(\mathbf X\mathbf v=\mathbf w\)}{Xv=w}}

For \(N_p\) averaging targets, we want each resolving kernel to have unit mass:
\[
  \int_\Omega A^{(k)}(r)\dd r=1,
  \qquad k=1,\dots,N_p.
\]
Using the average vector \(\mathbf v\), this becomes
\[
  [\mathbf X\mathbf v]_k=1,
  \qquad k=1,\dots,N_p.
\]

Define the desired mass vector
\[
  \mathbf w\in\Pspace
\]
by
\[
  [\mathbf w]_k=1,
  \qquad k=1,\dots,N_p.
\]
Then all unimodularity constraints can be written compactly as
\[
  \mathbf X\mathbf v=\mathbf w.
\]

This is the constrained SOLA condition for averages.

\begin{remark}[What the symbols mean]
The operator \(G:\M\to\D\) is lightface because it acts on the model function space. The vector
\(\mathbf v=G(1_\Omega)\) is bold because it is a finite-dimensional data-space object. The
matrix \(\mathbf X:\D\to\Pspace\) is bold because it is a finite-dimensional map between
Euclidean representation spaces. The vector \(\mathbf w\) is bold because it is a finite-dimensional
property-space object.
\end{remark}

The constraint
\[
  \mathbf X\mathbf v=\mathbf w
\]
has a clean interpretation:

\[
  \text{constant model}
  \quad\longmapsto\quad
  \text{constant average}.
\]

Indeed, if \(m=c1_\Omega\), then
\[
  G(m)=G(c1_\Omega)=cG(1_\Omega)=c\mathbf v.
\]
Applying the SOLA matrix gives
\[
  \mathbf XG(m)=\mathbf X(c\mathbf v)=c\mathbf X\mathbf v.
\]
If \(\mathbf X\mathbf v=\mathbf w\), and \([\mathbf w]_k=1\), then
\[
  [\mathbf XG(m)]_k=c,
  \qquad k=1,\dots,N_p.
\]
So each reported average reproduces constants correctly.

\begin{checkpointbox}
Unimodularity is not an abstract extra condition. It says that SOLA averages should return the
right answer when the model is constant.
\end{checkpointbox}

\section{The constrained minimization problem}

We now add this constraint to the noiseless SOLA problem.

The unconstrained problem was
\[
  \mathbf X_0
  =
  \argmin_{\mathbf X}
  \left\|
    \Tau-\mathbf XG
  \right\|_{\HS}^2.
\]
For averaging targets, we instead solve
\[
  \mathbf X_c
  =
  \argmin_{\mathbf X}
  \left\|
    \Tau-\mathbf XG
  \right\|_{\HS}^2
  \quad
  \text{subject to}
  \quad
  \mathbf X\mathbf v=\mathbf w.
\]

This says:

\[
  \text{match the target kernels as well as possible,}
\]
but only among data combinations whose resolving kernels have the correct total mass.

The constraint may force a slightly worse target fit than the unconstrained solution. That is the
price of calibration. If the unconstrained resolving kernels already have unit mass, then nothing
changes. But if they do not, the constrained solution adjusts the weights.

\begin{intuitionbox}
The unconstrained solution says: ``give me the closest shape.'' The constrained solution says:
``give me the closest shape among kernels that are actually averages.''
\end{intuitionbox}

The adjustment is surprisingly simple. It is a rank-one correction to the unconstrained solution.

\section{Closed-form correction}

Let
\[
  \mathbf X_0
  =
  \Tau G^*(GG^*)^{-1}
\]
be the noiseless unconstrained solution, assuming \(GG^*\) is invertible on \(\D\).

Define
\[
  \mathbf u
  :=
  (GG^*)^{-1}\mathbf v
  \in\D,
\]
and
\[
  \beta
  :=
  \langle \mathbf v,\mathbf u\rangle_{\D}.
\]
Equivalently,
\[
  \beta
  =
  \langle \mathbf v,(GG^*)^{-1}\mathbf v\rangle_{\D}.
\]
Assuming \(\mathbf v\neq 0\), we have \(\beta>0\).

The constrained solution is
\[
  \mathbf X_c
  =
  \mathbf X_0
  -
  \frac{\mathbf X_0\mathbf v-\mathbf w}{\beta}
  \mathbf u^*.
\]
Here
\[
  \frac{\mathbf X_0\mathbf v-\mathbf w}{\beta}
  \mathbf u^*
\]
is a rank-one map from \(\D\) to \(\Pspace\). It takes a data-space vector, tests it against
\(\mathbf u\), and returns a property-space correction in the direction
\[
  \mathbf X_0\mathbf v-\mathbf w.
\]

Let us check that the constraint is satisfied. Applying \(\mathbf X_c\) to \(\mathbf v\) gives
\[
  \mathbf X_c\mathbf v
  =
  \mathbf X_0\mathbf v
  -
  \frac{\mathbf X_0\mathbf v-\mathbf w}{\beta}
  \mathbf u^*\mathbf v.
\]
Since
\[
  \mathbf u^*\mathbf v
  =
  \langle \mathbf u,\mathbf v\rangle_{\D}
  =
  \beta,
\]
we obtain
\[
  \mathbf X_c\mathbf v
  =
  \mathbf X_0\mathbf v
  -
  (\mathbf X_0\mathbf v-\mathbf w)
  =
  \mathbf w.
\]
So the correction does exactly what it should.

\begin{remark}[Where the correction comes from]
One way to derive the formula is to introduce a Lagrange multiplier
\(\boldsymbol\lambda\in\Pspace\) and solve
\[
  \mathbf X_c
  =
  \argmin_{\mathbf X}
  \left\|
    \Tau-\mathbf XG
  \right\|_{\HS}^2
  \quad
  \text{subject to}
  \quad
  \mathbf X\mathbf v=\mathbf w.
\]
The stationarity equation has the form
\[
  (\Tau-\mathbf XG)G^*
  =
  \boldsymbol\lambda\mathbf v^*.
\]
Solving this equation gives the unconstrained SOLA operator plus a rank-one term. The multiplier
is then chosen so that \(\mathbf X_c\mathbf v=\mathbf w\), producing the formula above.
\end{remark}

The formula is useful because it shows how constrained SOLA differs from unconstrained SOLA.
It does not rebuild the whole solution from scratch. It corrects the unconstrained map by the
smallest adjustment needed to enforce the average condition.

At the kernel level, the correction modifies each resolving kernel just enough to give it unit
mass. In components,
\[
  A_c^{(k)}(r)
  =
  \sum_{i=1}^{N_d}[\mathbf X_c]_{ki}K_i(r),
\]
and the constraint ensures
\[
  \int_\Omega A_c^{(k)}(r)\dd r
  =
  [\mathbf X_c\mathbf v]_k
  =
  [\mathbf w]_k
  =
  1.
\]

\begin{figure}[ht]
  \centering
  \fbox{
    \begin{minipage}{0.82\textwidth}
      \vspace{0.8cm}
      \centering
      Placeholder for a figure comparing an unconstrained resolving kernel and a constrained
      resolving kernel. The constrained one should have unit area, even if its shape changes
      slightly.
      \vspace{0.8cm}
    \end{minipage}
  }
  \caption{The constrained SOLA solution adjusts the resolving kernel so that it has the correct
  mass for an average.}
  \label{fig:act3-constrained-kernel}
\end{figure}

\begin{lessonbox}
For averages, normalization is not decoration. It is part of the meaning of the estimate.
\end{lessonbox}

This closes the average-only story. We now know how to ask SOLA for localized averages, and
we know how to enforce the basic calibration that makes those averages meaningful.

But not every property is an average. Some targets represent gradients, contrasts, coefficients,
or more specialized linear diagnostics. Once we leave averages, the word ``unimodularity'' is no
longer broad enough. The next act replaces it with the more general idea of a resolving
constraint.



% ============================================================
% Act IV
% ============================================================

\part{Act IV: Beyond Averages}

% ------------------------------------------------------------
\chapter{Not Every Target Is an Average}
% ------------------------------------------------------------

\section{Derivative targets}

So far, the target kernels have represented localized averages. In that setting, unit mass was
essential. If
\[
  \int_\Omega \mathcal T^{(k)}(r)\dd r=1,
\]
then the component
\[
  [\mathbf p]_k
  =
  [\Tau(m)]_k
  =
  \int_\Omega \mathcal T^{(k)}(r)m(r)\dd r
\]
can be interpreted as a weighted average of \(m\).

But not every useful property is an average.

Sometimes we may want to estimate something more like a local derivative. For example, we may
want to know whether the model is increasing or decreasing near a location \(r_0\). A natural
target kernel for this kind of question would not be a positive bump. It would have positive and
negative lobes.

A cartoon derivative target might look like this:
\[
  \mathcal T^{(k)}(r)
  \approx
  \text{positive lobe to the right of }r_0
  -
  \text{positive lobe to the left of }r_0.
\]
Such a kernel compares nearby regions of the model. It is not trying to average the model. It is
trying to detect change.

For a derivative-like target, it is often natural to have
\[
  \int_\Omega \mathcal T^{(k)}(r)\dd r=0.
\]
This condition says that constants are ignored. Indeed, if \(m(r)=c\), then
\[
  [\mathbf p]_k
  =
  \int_\Omega \mathcal T^{(k)}(r)c\dd r
  =
  c\int_\Omega \mathcal T^{(k)}(r)\dd r
  =
  0.
\]
That is exactly what a derivative-like quantity should do. The derivative of a constant is zero.

\begin{intuitionbox}
An averaging kernel listens for level. A derivative kernel listens for change. One should
reproduce constants; the other should annihilate them.
\end{intuitionbox}

This already shows why ``unit mass'' cannot be the universal constraint. For averages, unit mass
is the correct calibration. For derivative-like targets, zero mass may be the correct calibration.

The target's meaning determines the constraint.

\begin{remark}[Point derivatives and regularity]
A true pointwise derivative, such as \(m'(r_0)\), is not defined for every model in a rough
function space such as \(L^2(\Omega)\). To make such a target rigorous, one either works in a
smoother model space or uses a regularized derivative target. In these notes, the phrase
``derivative target'' should be read as a smooth linear diagnostic designed to behave like a
localized derivative.
\end{remark}

In SOLA terms, the goal would be to choose \(\mathbf X\) so that the resolving kernel
\[
  A^{(k)}(r)
  =
  \sum_{i=1}^{N_d}[\mathbf X]_{ki}K_i(r)
\]
resembles the derivative-like target \(\mathcal T^{(k)}(r)\). But the relevant calibration would
not be
\[
  \int_\Omega A^{(k)}(r)\dd r=1.
\]
It might instead be
\[
  \int_\Omega A^{(k)}(r)\dd r=0,
\]
or some other condition that ensures the resolving kernel has the correct derivative-like
behavior.

\begin{checkpointbox}
The moment we leave averages, unimodularity becomes only one example of a broader idea:
the resolving kernel must satisfy the calibration needed for its intended interpretation.
\end{checkpointbox}

\section{Basis coefficient targets}

Another kind of target arises when we want to estimate coefficients of the model relative to a
chosen family of functions.

Suppose we have functions
\[
  \phi_1,\dots,\phi_{N_p}\in\M.
\]
If these functions are orthonormal, then the analysis coefficients of \(m\) are
\[
  [\mathbf p]_k
  =
  [\Tau(m)]_k
  =
  \langle \phi_k,m\rangle_{\M},
  \qquad k=1,\dots,N_p.
\]
In this case, the target kernels are simply
\[
  \mathcal T^{(k)}=\phi_k.
\]

This kind of target is not necessarily local. A basis function may be global, oscillatory, radial,
angular, compactly supported, or adapted to a numerical discretization. The property vector
\[
  \mathbf p=\Tau(m)
\]
then records how much of \(m\) is seen in each chosen direction.

\begin{intuitionbox}
An averaging target asks, ``what is the model doing near here?'' A coefficient target asks,
``how much of this chosen pattern is present in the model?''
\end{intuitionbox}

There is a small but important detail. If the family \(\{\phi_k\}_{k=1}^{N_p}\) is not orthonormal,
then the coefficient of \(m\) in that basis is not usually
\[
  \langle \phi_k,m\rangle_{\M}.
\]
Instead, one must use the appropriate dual family or the relevant mass matrix. For example, if
\(\{\phi^k\}_{k=1}^{N_p}\) is the dual family, then the coefficient-like target may be
\[
  [\mathbf p]_k
  =
  \langle \phi^k,m\rangle_{\M}.
\]
Equivalently, in finite-dimensional coordinates, one must keep track of the Gram matrix.

This is not a technical nuisance. It is the same lesson we have already met several times:
interpretation depends on geometry. A coefficient is not just a number. It is a number relative
to a chosen representation.

For SOLA, the target map is still linear:
\[
  \Tau:\M\to\Pspace,
  \qquad
  [\Tau(m)]_k
  =
  \langle \mathcal T^{(k)},m\rangle_{\M}.
\]
The only thing that changes is the meaning of the target kernels \(\mathcal T^{(k)}\). They are no
longer localized averages. They are coefficient-extraction kernels, or more generally
representation-dependent linear diagnostics.

The resolving kernels are again
\[
  A^{(k)}(r)
  =
  \sum_{i=1}^{N_d}[\mathbf X]_{ki}K_i(r).
\]
The question is now whether \(A^{(k)}\) reproduces the coefficient-like target
\(\mathcal T^{(k)}\) well enough for the intended use.

For such targets, unit mass may have no special meaning. The right constraint might instead be
orthogonality to some nuisance mode, correct response to a basis function, or exact reproduction
of a low-dimensional subspace.

\begin{warningbox}
A constraint that is natural for averages can be meaningless for coefficient targets. The target
decides the calibration, not the other way around.
\end{warningbox}

\section{Contrasts and ratios}

Some applications do not stop at individual SOLA outputs. They combine them further. For
example, one might estimate two model properties and then form a difference, contrast, or ratio.

This is common in geophysics. One may estimate two quantities, such as two wave-speed
parameters, and then form a ratio-like diagnostic. The danger is that each SOLA output comes
with its own resolving kernel.

Suppose we have two SOLA outputs,
\[
  [\widetilde{\mathbf p}]_1
  =
  \int_\Omega A^{(1)}(r)m_1(r)\dd r,
\]
and
\[
  [\widetilde{\mathbf p}]_2
  =
  \int_\Omega A^{(2)}(r)m_2(r)\dd r.
\]
If \(A^{(1)}\) and \(A^{(2)}\) are different, then these two numbers are not measuring the two
models in exactly the same way. They may correspond to different spatial averages, different
amounts of smoothing, or different mixtures of nearby structure.

A ratio such as
\[
  \frac{[\widetilde{\mathbf p}]_1}{[\widetilde{\mathbf p}]_2}
\]
is therefore not automatically the ratio of the intended target properties. It is the ratio of two
resolved quantities, each with its own resolving kernel.

For the intended ratio interpretation to be safe, one would like the relevant resolving kernels to
match in the way required by the physical question. In an idealized case, one might want
\[
  A^{(1)}=A^{(2)}
\]
or, more generally, one might want both resolving kernels to reproduce their targets in a jointly
compatible way.

\begin{warningbox}
Two SOLA numbers may look comparable while being averages with different kernels. A ratio of
such numbers is then a ratio of two different questions.
\end{warningbox}

This is not a reason to avoid contrasts or ratios altogether. It is a reason to label them
honestly. If the numerator and denominator have different resolving kernels, the ratio must be
interpreted through those kernels.

\begin{intuitionbox}
A ratio is safe only when the numerator and denominator have been measured with compatible
lenses. If the lenses differ, the ratio inherits the mismatch.
\end{intuitionbox}

This is another sign that the word ``unimodularity'' is too narrow for the general SOLA story.
For ratios, the relevant condition may not be unit mass. It may be matched resolution. For
derivatives, it may be zero response to constants. For coefficient targets, it may be exact response
to selected basis functions.

All of these are examples of the same broader idea: resolving constraints.

% ------------------------------------------------------------
\chapter{Resolving Constraints}
% ------------------------------------------------------------

\section{From unimodularity to calibration}

In the previous act, the constraint
\[
  \mathbf X\mathbf v=\mathbf w
\]
had a specific meaning. It enforced unit mass of the resolving kernels. For averaging targets,
this was exactly what we needed:
\[
  \int_\Omega A^{(k)}(r)\dd r=1.
\]

But now we have seen other kinds of targets. A derivative-like target might need zero response
to constants. A coefficient target might need correct response to a chosen basis function. A
ratio might require matched resolving behavior between two outputs.

So the same algebraic form can carry a broader meaning.

A resolving constraint is any condition imposed on the SOLA weights so that the resulting
resolving kernels have the structural property needed for interpretation.

For averages:
\[
  \text{resolving constraint}
  =
  \text{unit mass}.
\]
For derivative-like targets:
\[
  \text{resolving constraint}
  =
  \text{zero response to constants, or correct derivative calibration}.
\]
For coefficient targets:
\[
  \text{resolving constraint}
  =
  \text{correct response to selected basis directions}.
\]
For ratios:
\[
  \text{resolving constraint}
  =
  \text{compatible resolution between numerator and denominator}.
\]

\begin{lessonbox}
Unimodularity is not the principle. It is one example of the principle. The principle is
calibration of the resolving kernels.
\end{lessonbox}

This shift in language matters. If we keep saying ``unimodularity'', we may accidentally think
only about averages. But SOLA is broader than averages. It is a method for constructing linear
property-like quantities from linear data, and each property brings its own interpretation.

\section{General constraint form}

We now keep the same algebraic shape, but interpret it more generally.

Let
\[
  \mathbf v\in\D,
  \qquad
  \mathbf w\in\Pspace.
\]
The constrained SOLA problem has the form
\[
  \mathbf X_c
  =
  \argmin_{\mathbf X}
  \left\|
    \Tau-\mathbf XG
  \right\|_{\HS}^2
  \quad
  \text{subject to}
  \quad
  \mathbf X\mathbf v=\mathbf w.
\]

For averages, we used
\[
  \mathbf v=G(1_\Omega),
\]
and
\[
  [\mathbf w]_k=1.
\]
Then \(\mathbf X\mathbf v=\mathbf w\) forced each resolving kernel to have unit mass.

For derivative-like targets, one might still use
\[
  \mathbf v=G(1_\Omega),
\]
but choose
\[
  [\mathbf w]_k=0
\]
for components that should annihilate constants. This gives
\[
  \int_\Omega A^{(k)}(r)\dd r=0.
\]

For other targets, \(\mathbf v\) may represent a different test model. Suppose \(q\in\M\) is a
model pattern whose response should be calibrated. Define
\[
  \mathbf v=G(q).
\]
Then the constraint
\[
  \mathbf X\mathbf v=\mathbf w
\]
says
\[
  \mathbf XG(q)=\mathbf w.
\]
Since \(A=\mathbf XG\), this is
\[
  A(q)=\mathbf w.
\]
Thus the approximate property map \(A\) is forced to act on \(q\) in the prescribed way.

If the desired property map is \(\Tau\), a natural choice is often
\[
  \mathbf w=\Tau(q).
\]
Then the constraint becomes
\[
  A(q)=\Tau(q).
\]
In words: the SOLA approximation must reproduce the exact target property on the selected
test model \(q\).

\begin{intuitionbox}
A resolving constraint gives the resolving kernels a calibration exam. The test model \(q\) is the
exam question, and \(\mathbf w\) is the required answer.
\end{intuitionbox}

This interpretation makes the average case look less mysterious. There, the test model was the
constant function:
\[
  q=1_\Omega.
\]
The required answer was
\[
  \mathbf w=(1,\dots,1)^\top.
\]
So unimodularity was simply the demand that the SOLA averages pass the constant-model test.

\section{Several constraints}

Sometimes one calibration condition is not enough. We may want the resolving kernels to
respond correctly to several test models
\[
  q_1,\dots,q_L\in\M.
\]
Define the finite-dimensional data vectors
\[
  \mathbf v_\ell=G(q_\ell)\in\D,
  \qquad \ell=1,\dots,L.
\]
Collect these into the finite matrix
\[
  \mathbf V
  =
  \begin{pmatrix}
    \mathbf v_1 & \cdots & \mathbf v_L
  \end{pmatrix}.
\]
Likewise, choose the desired responses
\[
  \mathbf w_\ell\in\Pspace,
  \qquad \ell=1,\dots,L,
\]
and collect them into
\[
  \mathbf W
  =
  \begin{pmatrix}
    \mathbf w_1 & \cdots & \mathbf w_L
  \end{pmatrix}.
\]
Then the constraint becomes
\[
  \mathbf X\mathbf V=\mathbf W.
\]

A common choice is
\[
  \mathbf w_\ell=\Tau(q_\ell),
\]
so that
\[
  \mathbf XG(q_\ell)=\Tau(q_\ell),
  \qquad \ell=1,\dots,L.
\]
Equivalently,
\[
  A(q_\ell)=\Tau(q_\ell),
  \qquad \ell=1,\dots,L.
\]

This says that the SOLA approximate property map \(A=\mathbf XG\) must agree with the target
property map \(\Tau\) on a selected finite-dimensional set of test models.

\begin{checkpointbox}
A single constraint \(\mathbf X\mathbf v=\mathbf w\) calibrates one test direction. Multiple
constraints \(\mathbf X\mathbf V=\mathbf W\) calibrate several test directions at once.
\end{checkpointbox}

There is, of course, a trade-off. Every extra constraint reduces the freedom available to match
the target kernels and, later, to control noise. Constraints are useful only when they protect the
interpretation of the reported quantity.

\section{What the constraint should mean}

A resolving constraint should not be chosen because it is algebraically convenient. It should be
chosen because it protects the meaning of the SOLA output.

If the output is called an average, the resolving kernel should reproduce constants correctly.
That leads to unit mass:
\[
  \int_\Omega A^{(k)}(r)\dd r=1.
\]

If the output is called a derivative-like quantity, the resolving kernel should ignore constants:
\[
  \int_\Omega A^{(k)}(r)\dd r=0.
\]

If the output is meant to represent a coefficient, the resolving kernel should be calibrated
against the relevant basis or dual-basis functions.

If the output will be used in a ratio, the resolving kernels of the numerator and denominator
should be checked together, not separately.

The practical rule is:

\[
  \text{choose the constraint according to how the reported number will be interpreted.}
\]

\begin{warningbox}
A constraint can make a SOLA estimate more interpretable, but only if the constraint matches
the interpretation. The wrong constraint is just decorative algebra.
\end{warningbox}

This broader view also makes SOLA more flexible. The method is not restricted to localized
averages. It can be adapted to many linear diagnostics, provided we are honest about what the
resolving kernels actually represent.

\begin{lessonbox}
A resolving constraint says: preserve the feature of the resolving kernel that makes the reported
number meaningful.
\end{lessonbox}

We now have the clean deterministic skeleton of SOLA:

\[
  \text{choose target kernels},
  \qquad
  \text{build resolving kernels from data kernels},
  \qquad
  \text{enforce the constraints needed for interpretation}.
\]

The next ingredient is noise. Real data are not exact, and the weights used to sharpen a resolving
kernel can also amplify observational error. That is where the resolution-noise trade-off enters.

% ============================================================
% Act V
% ============================================================

\part{Act V: Noise Enters the Room}

% ------------------------------------------------------------
\chapter{Data Are Not Exact}
% ------------------------------------------------------------

\section{The noisy data model}

So far, the data have been treated as exact:
\[
  \mathbf d = G(\bar m).
\]
This was useful while building the basic SOLA idea. It allowed us to focus on the algebra of
kernels without worrying about measurement uncertainty.

But real data are not exact. The observed data are better written as
\[
  \widetilde{\mathbf d}
  =
  G(\bar m)+\boldsymbol\eta,
\]
where \(\bar m\in\M\) is the unknown true model and
\[
  \boldsymbol\eta\in\D
\]
is observational noise.

In the standard Gaussian version, we assume
\[
  \boldsymbol\eta\sim N(\mathbf 0,\mathbf C_D),
\]
where
\[
  \mathbf C_D\in\mathbb R^{N_d\times N_d}
\]
is the data-noise covariance matrix.

Thus the observed data vector is
\[
  \widetilde{\mathbf d}
  =
  \begin{pmatrix}
    \widetilde d_1\\
    \vdots\\
    \widetilde d_{N_d}
  \end{pmatrix},
  \qquad
  [\widetilde{\mathbf d}]_i=\widetilde d_i,
\]
with
\[
  \widetilde d_i
  =
  [G(\bar m)]_i+\eta_i.
\]

In integral form,
\[
  \widetilde d_i
  =
  \int_\Omega K_i(r)\bar m(r)\dd r+\eta_i,
  \qquad i=1,\dots,N_d.
\]

The covariance matrix \(\mathbf C_D\) tells us how uncertain the data are. If
\[
  \mathbf C_D=\sigma_D^2\mathbf I,
\]
then the data errors are independent and have common variance \(\sigma_D^2\). More generally,
the off-diagonal entries of \(\mathbf C_D\) encode correlations between different data errors.

\begin{intuitionbox}
The kernels tell us what the experiment tries to see. The covariance tells us how clearly it sees
it.
\end{intuitionbox}

This changes the SOLA problem. We no longer only care about constructing a resolving kernel
close to the target. We also care about how much noise the chosen data combination will amplify.

\section{Noise amplification}

Suppose we choose weights
\[
  \mathbf x=(x_1,\dots,x_{N_d})^\top.
\]
The corresponding SOLA-style data combination is
\[
  \sum_{i=1}^{N_d}x_i\widetilde d_i.
\]
Substituting the noisy data model gives
\[
  \sum_{i=1}^{N_d}x_i\widetilde d_i
  =
  \sum_{i=1}^{N_d}x_i[G(\bar m)]_i
  +
  \sum_{i=1}^{N_d}x_i\eta_i.
\]
The first term is the signal part:
\[
  \sum_{i=1}^{N_d}x_i[G(\bar m)]_i
  =
  \int_\Omega R_{\mathbf x}(r)\bar m(r)\dd r,
\]
where
\[
  R_{\mathbf x}(r)
  =
  \sum_{i=1}^{N_d}x_iK_i(r).
\]
The second term is the noise part:
\[
  \sum_{i=1}^{N_d}x_i\eta_i.
\]

So the same weights \(\mathbf x\) that build the resolving kernel also combine the noise.

This is the trade-off. If we choose aggressive weights, we may obtain a sharper resolving kernel,
but those same weights may amplify the observational errors.

To see this explicitly, consider the scalar noise contribution
\[
  \epsilon_{\mathbf x}
  =
  \sum_{i=1}^{N_d}x_i\eta_i.
\]
In vector notation,
\[
  \epsilon_{\mathbf x}
  =
  \mathbf x^\top\boldsymbol\eta.
\]
Since
\[
  \boldsymbol\eta\sim N(\mathbf 0,\mathbf C_D),
\]
we have
\[
  \epsilon_{\mathbf x}\sim N(0,\mathbf x^\top\mathbf C_D\mathbf x).
\]
Thus the variance of the weighted noise is
\[
  \operatorname{Var}(\epsilon_{\mathbf x})
  =
  \mathbf x^\top\mathbf C_D\mathbf x.
\]

A resolving kernel can be beautifully localized but numerically dangerous if producing it requires
large or unstable weights.

\begin{warningbox}
Sharper resolution is not free. The weights that sharpen the resolving kernel may also turn
quiet data noise into loud property noise.
\end{warningbox}

For multiple properties, the weights are collected into a finite matrix
\[
  \mathbf X:\D\to\Pspace.
\]
The SOLA estimate is
\[
  \widetilde{\mathbf p}
  =
  \mathbf X\widetilde{\mathbf d}.
\]
Using the noisy data model,
\[
  \widetilde{\mathbf p}
  =
  \mathbf XG(\bar m)+\mathbf X\boldsymbol\eta.
\]
The signal part is
\[
  \mathbf XG(\bar m)=A(\bar m),
\]
where
\[
  A:=\mathbf XG:\M\to\Pspace
\]
is the approximate property map. The noise part is
\[
  \mathbf X\boldsymbol\eta.
\]

Since \(\boldsymbol\eta\sim N(\mathbf 0,\mathbf C_D)\), the propagated property noise has
covariance
\[
  \mathbf C_P
  =
  \mathbf X\mathbf C_D\mathbf X^*.
\]

This formula is simple because the SOLA map is linear. That is one of the great conveniences of
the method.

\section{The resolution-noise trade-off}

We now have two competing desires.

First, we want the approximate property map
\[
  A=\mathbf XG
\]
to be close to the target property map
\[
  \Tau.
\]
This is the resolution desire:
\[
  \mathbf XG\approx \Tau.
\]

Second, we want the propagated noise
\[
  \mathbf X\boldsymbol\eta
\]
to be small. Since its covariance is
\[
  \mathbf C_P=\mathbf X\mathbf C_D\mathbf X^*,
\]
we need some scalar measure of the size of this covariance. A common choice is the trace:
\[
  \Tr(\mathbf X\mathbf C_D\mathbf X^*).
\]
This is the total propagated variance across the property components.

So noisy SOLA balances:
\[
  \text{small resolving-target mismatch}
\]
against
\[
  \text{small propagated data noise}.
\]

\begin{intuitionbox}
SOLA tunes two dials at once: one sharpens the station, the other suppresses static.
\end{intuitionbox}

If we care only about matching the target kernels, we may choose weights that are unstable.
If we care only about suppressing noise, we may choose weights that produce very smooth,
uninteresting resolving kernels. SOLA lives between these extremes.

The central compromise is:
\[
  \text{make }\mathbf XG\text{ close to }\Tau,
  \qquad
  \text{but do not let }\mathbf X\mathbf C_D\mathbf X^*\text{ become too large}.
\]

\begin{checkpointbox}
Noise changes the SOLA problem from pure kernel approximation into a resolution-noise
trade-off.
\end{checkpointbox}

% ------------------------------------------------------------
\chapter{Noise-Aware SOLA}
% ------------------------------------------------------------

\section{The noise-aware objective}

We now write the noisy SOLA problem as an optimization problem.

The noiseless unconstrained objective was
\[
  \left\|
    \Tau-\mathbf XG
  \right\|_{\HS}^2.
\]
This term measures how far the resolving kernels are from the target kernels.

To include data noise, we add the propagated variance term
\[
  \Tr(\mathbf X\mathbf C_D\mathbf X^*).
\]
Thus the noise-aware SOLA objective is
\[
  J(\mathbf X)
  =
  \left\|
    \Tau-\mathbf XG
  \right\|_{\HS}^2
  +
  \Tr(\mathbf X\mathbf C_D\mathbf X^*).
\]

The first term rewards resolution. The second term penalizes noise amplification.

\begin{remark}[About trade-off parameters]
Some presentations introduce an explicit trade-off parameter multiplying the noise term or the
resolution term. In the simplified form used here, the relative scale is encoded directly in
\(\mathbf C_D\) and in the chosen norms. One can add a tunable parameter if one wants an
explicit resolution-noise dial.
\end{remark}

The minimizer of
\[
  J(\mathbf X)
\]
is the unconstrained noise-aware SOLA operator.

Let us derive its normal equation. Perturb \(\mathbf X\) by \(\mathbf H\). The first variation of the
misfit term is
\[
  -2
  \left\langle
    (\Tau-\mathbf XG)G^*,
    \mathbf H
  \right\rangle_{\HS}.
\]
The first variation of the variance term is
\[
  2
  \left\langle
    \mathbf X\mathbf C_D,
    \mathbf H
  \right\rangle_{\HS}.
\]
Therefore stationarity gives
\[
  -(\Tau-\mathbf XG)G^*+\mathbf X\mathbf C_D=0.
\]
Equivalently,
\[
  (\Tau-\mathbf XG)G^*
  =
  \mathbf X\mathbf C_D.
\]
Rearranging,
\[
  \Tau G^*
  =
  \mathbf XGG^*+\mathbf X\mathbf C_D.
\]
Hence
\[
  \Tau G^*
  =
  \mathbf X(GG^*+\mathbf C_D).
\]

Define the finite-dimensional data-space matrix
\[
  \mathbf M
  :=
  GG^*+\mathbf C_D.
\]
If \(\mathbf C_D\) is positive definite, then \(\mathbf M\) is invertible. The unconstrained
noise-aware solution is then
\[
  \mathbf X_0
  =
  \Tau G^*\mathbf M^{-1}.
\]

Compare this with the noiseless unconstrained solution:
\[
  \mathbf X_0^{\mathrm{noiseless}}
  =
  \Tau G^*(GG^*)^{-1}.
\]
Noise has replaced \(GG^*\) by
\[
  \mathbf M=GG^*+\mathbf C_D.
\]

\begin{intuitionbox}
The covariance \(\mathbf C_D\) acts like a brake. Directions in data space with large noise become
less attractive to the SOLA weights.
\end{intuitionbox}

\section{The constrained noisy problem}

We now add the resolving constraint.

For averages, this constraint was unimodularity:
\[
  \mathbf X\mathbf v=\mathbf w.
\]
For more general targets, it may represent another calibration condition, but the algebraic form
is the same.

The constrained noisy SOLA problem is
\[
  \mathbf X_c
  =
  \argmin_{\mathbf X}
  \left\{
    \left\|
      \Tau-\mathbf XG
    \right\|_{\HS}^2
    +
    \Tr(\mathbf X\mathbf C_D\mathbf X^*)
  \right\}
  \quad
  \text{subject to}
  \quad
  \mathbf X\mathbf v=\mathbf w.
\]

This is the full basic SOLA problem in these notes. It contains:

\[
  \begin{array}{lll}
    \text{target matching} &:& \left\|\Tau-\mathbf XG\right\|_{\HS}^2,\\[0.3em]
    \text{noise control} &:& \Tr(\mathbf X\mathbf C_D\mathbf X^*),\\[0.3em]
    \text{interpretive calibration} &:& \mathbf X\mathbf v=\mathbf w.
  \end{array}
\]

\begin{checkpointbox}
Noise-aware constrained SOLA asks for the best resolving kernels that are calibrated correctly
and do not amplify data noise too much.
\end{checkpointbox}

The constrained problem can be solved by the same Lagrange multiplier logic as before. The
only change is that \(GG^*\) is replaced by
\[
  \mathbf M=GG^*+\mathbf C_D.
\]

\section{Closed-form solution}

Define
\[
  \mathbf M
  :=
  GG^*+\mathbf C_D.
\]
Then set
\[
  \mathbf u
  :=
  \mathbf M^{-1}\mathbf v,
\]
and
\[
  \beta
  :=
  \langle \mathbf v,\mathbf u\rangle_{\D}.
\]
Assuming \(\mathbf v\neq \mathbf 0\) and \(\mathbf M\) is positive definite, we have
\[
  \beta>0.
\]

The constrained noise-aware SOLA solution is
\[
  \mathbf X_c
  =
  \Tau G^*\mathbf M^{-1}
  -
  \frac{\Tau G^*(\mathbf u)-\mathbf w}{\beta}
  \mathbf u^*.
\]

This is the noisy analogue of the constrained noiseless correction. The first term
\[
  \Tau G^*\mathbf M^{-1}
\]
is the unconstrained noise-aware solution. The second term is the correction needed to enforce
\[
  \mathbf X_c\mathbf v=\mathbf w.
\]

Let us verify the constraint. Applying \(\mathbf X_c\) to \(\mathbf v\), we get
\[
  \mathbf X_c\mathbf v
  =
  \Tau G^*\mathbf M^{-1}\mathbf v
  -
  \frac{\Tau G^*(\mathbf u)-\mathbf w}{\beta}
  \mathbf u^*\mathbf v.
\]
Since
\[
  \mathbf u=\mathbf M^{-1}\mathbf v,
\]
the first term is
\[
  \Tau G^*\mathbf M^{-1}\mathbf v
  =
  \Tau G^*(\mathbf u).
\]
Also,
\[
  \mathbf u^*\mathbf v
  =
  \langle \mathbf u,\mathbf v\rangle_{\D}
  =
  \beta.
\]
Therefore,
\[
  \mathbf X_c\mathbf v
  =
  \Tau G^*(\mathbf u)
  -
  \bigl(\Tau G^*(\mathbf u)-\mathbf w\bigr)
  =
  \mathbf w.
\]
So the formula satisfies the resolving constraint exactly.

\begin{remark}[The role of \(\mathbf M\)]
The matrix
\[
  \mathbf M=GG^*+\mathbf C_D
\]
combines the geometry of the sensitivity kernels with the geometry of the data noise. The term
\(GG^*\) knows which data combinations are supported by the model-space kernels. The term
\(\mathbf C_D\) knows which data combinations are noisy.
\end{remark}

At the kernel level, the \(k\)th resolving kernel is still
\[
  A_c^{(k)}(r)
  =
  \sum_{i=1}^{N_d}
  [\mathbf X_c]_{ki}K_i(r).
\]
The difference is that the weights now know about noise. A kernel that could be obtained only by
using extremely noisy data combinations will be discouraged by the variance term.

\begin{lessonbox}
Noise-aware SOLA does not merely ask, ``which resolving kernel is closest to the target?'' It asks,
``which calibrated resolving kernel is closest to the target without paying too much in noise?''
\end{lessonbox}

\section{The usual propagated covariance}

Once the SOLA matrix \(\mathbf X\) has been chosen, the usual uncertainty calculation is
straightforward.

The observed data are
\[
  \widetilde{\mathbf d}
  =
  G(\bar m)+\boldsymbol\eta,
  \qquad
  \boldsymbol\eta\sim N(\mathbf 0,\mathbf C_D).
\]
The SOLA estimate is
\[
  \widetilde{\mathbf p}
  =
  \mathbf X\widetilde{\mathbf d}.
\]
Substituting the data model,
\[
  \widetilde{\mathbf p}
  =
  \mathbf XG(\bar m)
  +
  \mathbf X\boldsymbol\eta.
\]
The random part is
\[
  \mathbf X\boldsymbol\eta.
\]
Since \(\boldsymbol\eta\) has covariance \(\mathbf C_D\), the propagated covariance is
\[
  \mathbf C_P
  =
  \mathbf X\mathbf C_D\mathbf X^*.
\]

In components, the covariance between the \(k\)th and \(\ell\)th SOLA outputs is
\[
  [\mathbf C_P]_{k\ell}
  =
  \sum_{i=1}^{N_d}
  \sum_{j=1}^{N_d}
  [\mathbf X]_{ki}
  [\mathbf C_D]_{ij}
  [\mathbf X^*]_{j\ell}.
\]
If the data and property spaces use standard Euclidean inner products, then
\[
  \mathbf X^*=\mathbf X^\top,
\]
and this becomes
\[
  [\mathbf C_P]_{k\ell}
  =
  \sum_{i=1}^{N_d}
  \sum_{j=1}^{N_d}
  [\mathbf X]_{ki}
  [\mathbf C_D]_{ij}
  [\mathbf X]_{\ell j}.
\]

The diagonal entries
\[
  [\mathbf C_P]_{kk}
\]
are the propagated variances of the individual SOLA outputs. A common display is therefore
\[
  [\widetilde{\mathbf p}]_k
  \pm
  2\sqrt{[\mathbf C_P]_{kk}},
\]
which shows the estimate together with a nominal two-standard-deviation band.

\begin{figure}[ht]
  \centering
  \fbox{
    \begin{minipage}{0.82\textwidth}
      \vspace{0.8cm}
      \centering
      Placeholder for a figure showing SOLA property estimates
      \([\widetilde{\mathbf p}]_k\) with error bars
      \(2\sqrt{[\mathbf C_P]_{kk}}\).
      \vspace{0.8cm}
    \end{minipage}
  }
  \caption{The usual SOLA uncertainty display: propagated data-noise uncertainty after applying
  the SOLA map \(\mathbf X\).}
  \label{fig:act5-propagated-covariance}
\end{figure}

This covariance propagation is mathematically correct. If \(\boldsymbol\eta\) is Gaussian with
covariance \(\mathbf C_D\), then \(\mathbf X\boldsymbol\eta\) is Gaussian with covariance
\(\mathbf X\mathbf C_D\mathbf X^*\).

But we should be precise about what this means. It describes how observational noise moves
through the chosen SOLA map. It does not yet tell us everything one might want to know about
the true target property
\[
  \Tau(\bar m).
\]
That harder interpretive issue comes later.

\begin{checkpointbox}
The formula \(\mathbf C_P=\mathbf X\mathbf C_D\mathbf X^*\) is mathematically correct as linear
covariance propagation. The harder question is what inferential question it answers.
\end{checkpointbox}

We now have the classical SOLA ingredients in place: target kernels, resolving kernels,
calibration constraints, noise-aware optimization, and propagated covariance. The next act asks
what the resulting SOLA numbers actually estimate.

% ============================================================
% Act VI
% ============================================================

\part{Act VI: What Did SOLA Actually Estimate?}

% ------------------------------------------------------------
\chapter{The Target Is the Request, the Resolving Kernel Is the Answer}
% ------------------------------------------------------------

\section{The approximate property map}

We have now built the standard SOLA machinery. We have a forward map
\[
  G:\M\to\D,
\]
a property map
\[
  \Tau:\M\to\Pspace,
\]
and a finite-dimensional SOLA matrix
\[
  \mathbf X:\D\to\Pspace.
\]
The defining idea was to choose \(\mathbf X\) so that
\[
  \mathbf XG\approx \Tau.
\]

This composite map deserves its own name. We define
\[
  A:=\mathbf XG:\M\to\Pspace.
\]
The map \(A\) is the approximate property map synthesized from the data.

This notation is small, but the conceptual shift is large. The target map is \(\Tau\). The map
SOLA actually constructs is \(A\). These are not the same object unless the resolving kernels
exactly reproduce the target kernels.

For a model \(m\), the desired property vector is
\[
  \mathbf p=\Tau(m).
\]
The SOLA-accessible property vector, in the noiseless setting, is
\[
  A(m)=\mathbf XG(m).
\]

So there are two property maps in the room:
\[
  \Tau:\M\to\Pspace,
  \qquad
  A:\M\to\Pspace.
\]
The first one is what we asked for. The second one is what the data combination actually gives.

\begin{intuitionbox}
The target map \(\Tau\) is the question written on the wish list. The approximate map
\(A=\mathbf XG\) is the question the data can actually answer after SOLA has reshuffled them.
\end{intuitionbox}

At the kernel level, this means the following. The \(k\)th component of the target property is
represented by the target kernel \(\mathcal T^{(k)}\):
\[
  [\Tau(m)]_k
  =
  \langle \mathcal T^{(k)},m\rangle_{\M}.
\]
The \(k\)th component of the approximate property map is represented by the resolving kernel
\(A^{(k)}\):
\[
  [A(m)]_k
  =
  \langle A^{(k)},m\rangle_{\M}.
\]
The resolving kernel is built from the sensitivity kernels:
\[
  A^{(k)}
  =
  \sum_{i=1}^{N_d}[\mathbf X]_{ki}K_i.
\]

Thus SOLA produces the kernels \(A^{(k)}\), not merely the numbers
\([\widetilde{\mathbf p}]_k\). Those kernels are the actual lenses through which the model is
being viewed.

\begin{checkpointbox}
Once \(\mathbf X\) has been chosen, SOLA has chosen a new property map
\[
  A=\mathbf XG.
\]
The reported numbers should be interpreted through \(A\), not automatically through \(\Tau\).
\end{checkpointbox}

\section{The target property and the SOLA property}

Let \(\bar m\in\M\) denote the true model. The true target property vector is
\[
  \bar{\mathbf p}
  :=
  \Tau(\bar m).
\]
Its components are
\[
  [\bar{\mathbf p}]_k
  =
  [\Tau(\bar m)]_k
  =
  \langle \mathcal T^{(k)},\bar m\rangle_{\M}.
\]

Now ignore noise for a moment. If the data were exact, then
\[
  \mathbf d=G(\bar m).
\]
Applying the SOLA matrix gives
\[
  \widetilde{\mathbf p}
  =
  \mathbf X\mathbf d
  =
  \mathbf XG(\bar m)
  =
  A(\bar m).
\]
So the noiseless SOLA output is
\[
  \widetilde{\mathbf p}=A(\bar m).
\]

The desired target property is
\[
  \bar{\mathbf p}=\Tau(\bar m).
\]

Therefore the basic distinction is
\[
  A(\bar m)
  \quad\text{versus}\quad
  \Tau(\bar m).
\]

This is the interpretive fork in the road.

If
\[
  A=\Tau,
\]
then the SOLA property is exactly the target property:
\[
  A(\bar m)=\Tau(\bar m).
\]
In that case, there is no ambiguity.

But if
\[
  A\neq \Tau,
\]
then the SOLA output and the target property are generally different:
\[
  A(\bar m)\neq \Tau(\bar m).
\]

In components,
\[
  [\widetilde{\mathbf p}]_k
  =
  [A(\bar m)]_k
  =
  \langle A^{(k)},\bar m\rangle_{\M},
\]
whereas the desired target component is
\[
  [\bar{\mathbf p}]_k
  =
  [\Tau(\bar m)]_k
  =
  \langle \mathcal T^{(k)},\bar m\rangle_{\M}.
\]

So the difference is
\[
  [\bar{\mathbf p}]_k-[\widetilde{\mathbf p}]_k
  =
  \langle \mathcal T^{(k)}-A^{(k)},\bar m\rangle_{\M}.
\]

\begin{warningbox}
The SOLA output is exact for the resolving kernel, not automatically for the target kernel.
If \(A^{(k)}\neq\mathcal T^{(k)}\), then the two linear questions are different.
\end{warningbox}

This is not a failure of SOLA. It is the method telling us precisely what it managed to build.
The danger begins only when we forget the distinction and read
\[
  A(\bar m)
\]
as though it were
\[
  \Tau(\bar m).
\]

\section{Why kernel closeness is not enough}

A tempting thought now appears:

\[
  \text{if }A\text{ is close to }\Tau,\text{ then }A(\bar m)\text{ should be close to }\Tau(\bar m).
\]

This sounds reasonable, and in many practical situations it may be a useful heuristic. But as a
mathematical statement, it is missing a crucial ingredient.

The difference between the true target property and the SOLA-accessible property is
\[
  \Tau(\bar m)-A(\bar m)
  =
  (\Tau-A)(\bar m).
\]
Taking norms gives
\[
  \|\Tau(\bar m)-A(\bar m)\|_{\Pspace}
  =
  \|(\Tau-A)(\bar m)\|_{\Pspace}.
\]
Using the operator norm,
\[
  \|\Tau(\bar m)-A(\bar m)\|_{\Pspace}
  \le
  \|\Tau-A\|_{\M\to\Pspace}\,\|\bar m\|_{\M}.
\]

This inequality is useful, but it does not say what one might first hope. It says that the
property error is controlled by two things:

\[
  \text{resolving-target mismatch}
  \qquad\text{and}\qquad
  \text{size of the true model}.
\]

If \(\|\bar m\|_{\M}\) is controlled, then a small operator mismatch gives a small property error.
For example, if we know in advance that
\[
  \|\bar m\|_{\M}\le R,
\]
then
\[
  \|\Tau(\bar m)-A(\bar m)\|_{\Pspace}
  \le
  R\|\Tau-A\|_{\M\to\Pspace}.
\]
That is a genuine error bound.

But without such a model-side restriction, there is no uniform guarantee.

Indeed, suppose
\[
  \Tau-A\neq 0.
\]
Then there exists some direction \(h\in\M\) such that
\[
  (\Tau-A)(h)\neq \mathbf 0.
\]
For any scalar \(\alpha\), consider the model
\[
  m_\alpha=\alpha h.
\]
Then
\[
  \Tau(m_\alpha)-A(m_\alpha)
  =
  \alpha(\Tau-A)(h).
\]
Therefore
\[
  \|\Tau(m_\alpha)-A(m_\alpha)\|_{\Pspace}
  =
  |\alpha|\,\|(\Tau-A)(h)\|_{\Pspace}.
\]
By taking \(|\alpha|\) large enough, the property mismatch can be made arbitrarily large, even
though the operator mismatch \(\|\Tau-A\|\) has not changed.

\begin{warningbox}
A small mismatch between resolving and target kernels is not a model-independent error bound.
It becomes an error bound only after the model class has been restricted.
\end{warningbox}

This point is subtle because plots of kernels can be persuasive. If \(A^{(k)}\) visually resembles
\(\mathcal T^{(k)}\), it is natural to feel that the corresponding properties must be close. But
the model is what multiplies the difference:
\[
  [\Tau(\bar m)]_k-[A(\bar m)]_k
  =
  \langle \mathcal T^{(k)}-A^{(k)},\bar m\rangle_{\M}.
\]
A tiny kernel mismatch can matter a lot if the true model has large structure in exactly that
mismatch direction.

\begin{intuitionbox}
Kernel mismatch is the crack. The true model decides how much water flows through it.
\end{intuitionbox}

This is the first major limitation that must be stated carefully. SOLA gives us control over the
resolving kernels, and that control is valuable. But closeness of resolving kernels to target
kernels is not, by itself, a guarantee that the target properties of the true model have been
estimated accurately.

To get such a guarantee, one needs additional information about the true model. That information
may be a norm bound, a smoothness class, a deterministic admissible set, a probabilistic prior, or
some other model-side restriction. Without it, the data cannot rule out all directions that matter
to \(\Tau-A\).

\begin{checkpointbox}
The statement
\[
  A\approx\Tau
\]
is a statement about operators or kernels. To turn it into a statement about
\[
  A(\bar m)\approx\Tau(\bar m),
\]
we need some control of \(\bar m\).
\end{checkpointbox}

% ------------------------------------------------------------
\chapter{The Honest Sidestep}
% ------------------------------------------------------------

\section{Interpreting through resolving kernels}

There is a clean way to avoid overclaiming.

Do not pretend that the resolving kernel is the target kernel. Interpret the result through the
resolving kernel itself.

The SOLA output is
\[
  \widetilde{\mathbf p}
  =
  \mathbf X\widetilde{\mathbf d}.
\]
Ignoring observational noise for the moment,
\[
  \widetilde{\mathbf p}
  =
  A(\bar m).
\]
Therefore the \(k\)th component is
\[
  [\widetilde{\mathbf p}]_k
  =
  [A(\bar m)]_k
  =
  \langle A^{(k)},\bar m\rangle_{\M},
\]
where
\[
  A^{(k)}
  =
  \sum_{i=1}^{N_d}[\mathbf X]_{ki}K_i.
\]

This is not a second-rate interpretation. It is the direct interpretation of what SOLA has
constructed.

If \(A^{(k)}\) is a localized, unit-mass averaging kernel, then
\[
  [\widetilde{\mathbf p}]_k
\]
is a localized average of the true model using that resolving kernel. If \(A^{(k)}\) is broad, then
the number is a broad average. If \(A^{(k)}\) has side lobes, then the number includes those side
lobes. If \(A^{(k)}\) is derivative-like, then the number should be interpreted as that
derivative-like diagnostic.

\begin{lessonbox}
The safest SOLA interpretation is not ``this is the target property.'' It is ``this is the property
defined by the resolving kernel.''
\end{lessonbox}

This may sound like a retreat, but it is actually a strength of the method. SOLA exposes the
resolving kernels. It lets us inspect the actual linear functionals being reported. In that sense,
SOLA is unusually honest about resolution.

The mistake is not that SOLA gives resolving kernels. The mistake is to compute them, plot
them, admire them, and then interpret the numbers as if the target kernels had been achieved
exactly.

\section{Targets as auxiliary design objects}

The target kernels are still important. They guide the construction. Without them, the SOLA
optimization would not know what kind of property-like quantity we want.

But target kernels should be understood as auxiliary design objects.

A target kernel \(\mathcal T^{(k)}\) says:
\[
  \text{this is the linear question I would like to ask.}
\]
The resolving kernel \(A^{(k)}\) says:
\[
  \text{this is the linear question the data combination actually asks.}
\]

The SOLA optimization tries to make these two questions as similar as possible, subject to the
available data kernels, resolving constraints, and noise control. But the final interpretation
belongs to the question actually asked.

\begin{intuitionbox}
The target kernel is the sketch handed to the carpenter. The resolving kernel is the chair that
comes back from the workshop. Sit on the chair, not on the sketch.
\end{intuitionbox}

This view makes room for a mature use of SOLA. We can still design ambitious target kernels.
We can still use them to steer the calculation. We can still compare resolving kernels with
targets as a diagnostic.

But when reporting the result, we should say what was actually resolved.

For example, instead of saying
\[
  [\widetilde{\mathbf p}]_k
  \text{ estimates }
  \langle \mathcal T^{(k)},\bar m\rangle_{\M},
\]
we should say
\[
  [\widetilde{\mathbf p}]_k
  \text{ estimates }
  \langle A^{(k)},\bar m\rangle_{\M},
\]
and then show \(A^{(k)}\).

If \(A^{(k)}\) is close enough to \(\mathcal T^{(k)}\) for the scientific purpose at hand, excellent.
But that judgment depends on the model class and on the question being asked. It is not obtained
from the plot alone.

\begin{warningbox}
The target kernel is useful for design. The resolving kernel is necessary for interpretation.
Confusing those roles is where the trouble begins.
\end{warningbox}

\section{Each map point has its own kernel}

This becomes especially important when SOLA is used to build a map.

Suppose we choose target kernels
\[
  \mathcal T^{(1)},\dots,\mathcal T^{(N_p)}
\]
centred at different locations. The output is a finite-dimensional vector
\[
  \widetilde{\mathbf p}\in\Pspace,
\]
with components
\[
  [\widetilde{\mathbf p}]_k,
  \qquad k=1,\dots,N_p.
\]
It is tempting to plot these values as a profile or image and treat them as samples of the same
kind of quantity at different locations.

But each component has its own resolving kernel:
\[
  A^{(k)}
  =
  \sum_{i=1}^{N_d}[\mathbf X]_{ki}K_i.
\]
Thus each plotted value may correspond to a different average, a different width, a different
side-lobe structure, and a different response to unresolved model directions.

In a good SOLA calculation, the resolving kernels may be similar enough that the map has a
clean interpretation. But this has to be checked. It is not automatic.

\begin{figure}[ht]
  \centering
  \fbox{
    \begin{minipage}{0.82\textwidth}
      \vspace{0.8cm}
      \centering
      Placeholder for a figure showing a row of map values
      \([\widetilde{\mathbf p}]_k\), with the corresponding resolving kernels \(A^{(k)}(r)\)
      shown beneath each point. The visual message should be that every plotted number has its
      own kernel.
      \vspace{0.8cm}
    \end{minipage}
  }
  \caption{A SOLA map should be read together with its resolving kernels. Each reported value
  corresponds to its own resolved linear functional.}
  \label{fig:act6-map-with-resolving-kernels}
\end{figure}

This is particularly important near boundaries, gaps in coverage, or regions where the sensitivity
kernels change character. There, resolving kernels may become broader, more asymmetric, or
more oscillatory. The plotted SOLA number may still be perfectly valid, but its meaning has
changed.

\begin{checkpointbox}
A SOLA map is not merely a list of numbers. It is a list of numbers together with a list of
resolving kernels.
\end{checkpointbox}

This is the honest sidestep. We do not need to claim that SOLA has recovered the exact target
properties. We can say something more precise and more defensible:

\[
  \text{SOLA estimates the resolved properties defined by }A=\mathbf XG.
\]

If those resolved properties are scientifically useful, then SOLA has done something valuable.
If we want to go further and claim something about the exact target properties
\[
  \Tau(\bar m),
\]
then we need additional assumptions about \(\bar m\). That issue will return later, when we
discuss proxy models, prior information, and the difference between propagated noise and true
property uncertainty.

\begin{lessonbox}
The resolving kernel is not a diagnostic afterthought. It is part of the answer.
\end{lessonbox}

% ============================================================
% Act VII
% ============================================================

\part{Act VII: A Synthetic Cautionary Tale}

% ------------------------------------------------------------
\chapter{A Case That Looks Good}
% ------------------------------------------------------------

\section{Model, data, and properties}

We now build a small synthetic example. The point of the example is not to make SOLA look bad.
Quite the opposite: the example will first look like a perfectly respectable SOLA calculation.

The warning will come later.

Let
\[
  \M=L^2([0,1]),
\]
so the model is a square-integrable function
\[
  m:[0,1]\to\mathbb R.
\]
The data space is finite-dimensional:
\[
  \D\simeq\mathbb R^{N_d},
\]
and the property space is also finite-dimensional:
\[
  \Pspace\simeq\mathbb R^{N_p}.
\]

For the numerical example, one may take
\[
  N_d=100,
  \qquad
  N_p=20.
\]
Thus the data vector has one hundred components,
\[
  \mathbf d\in\mathbb R^{100},
\]
and the property vector has twenty components,
\[
  \mathbf p\in\mathbb R^{20}.
\]

The forward map is
\[
  G:\M\to\D,
\]
with components
\[
  [G(m)]_i
  =
  \langle K_i,m\rangle_{\M},
  \qquad i=1,\dots,N_d.
\]
Here
\[
  K_1,\dots,K_{N_d}\in\M
\]
are the sensitivity kernels.

The property map is
\[
  \Tau:\M\to\Pspace,
\]
with components
\[
  [\Tau(m)]_k
  =
  \langle \mathcal T^{(k)},m\rangle_{\M},
  \qquad k=1,\dots,N_p.
\]
The functions
\[
  \mathcal T^{(1)},\dots,\mathcal T^{(N_p)}
\]
are the target kernels.

For this example, the target kernels are localized bump averages. That is, each
\(\mathcal T^{(k)}\) is concentrated near a particular location and is normalized so that
\[
  \int_0^1 \mathcal T^{(k)}(r)\dd r=1.
\]
Thus each component
\[
  [\mathbf p]_k=[\Tau(m)]_k
\]
is intended to represent a localized average of the model.

\begin{figure}[ht]
  \centering
  \fbox{
    \begin{minipage}{0.82\textwidth}
      \vspace{0.8cm}
      \centering
      Placeholder for a plot of the synthetic sensitivity kernels
      \(K_1,\dots,K_{N_d}\) on \([0,1]\).
      \vspace{0.8cm}
    \end{minipage}
  }
  \caption{Synthetic sensitivity kernels \(K_i\) defining the forward map \(G\).}
  \label{fig:act7-sensitivity-kernels}
\end{figure}

\begin{figure}[ht]
  \centering
  \fbox{
    \begin{minipage}{0.82\textwidth}
      \vspace{0.8cm}
      \centering
      Placeholder for a plot of the localized averaging target kernels
      \(\mathcal T^{(1)},\dots,\mathcal T^{(N_p)}\).
      \vspace{0.8cm}
    \end{minipage}
  }
  \caption{Localized target kernels \(\mathcal T^{(k)}\) defining the property map \(\Tau\).}
  \label{fig:act7-target-kernels}
\end{figure}

The observed data are generated from a true model \(\bar m\), but we temporarily keep this model
hidden. We only show the data:
\[
  \widetilde{\mathbf d}
  =
  G(\bar m)+\boldsymbol\eta,
  \qquad
  \boldsymbol\eta\sim N(\mathbf 0,\mathbf C_D).
\]
This gives a standard-looking inverse problem. We have data, uncertainty estimates, sensitivity
kernels, and target kernels.

\begin{figure}[ht]
  \centering
  \fbox{
    \begin{minipage}{0.82\textwidth}
      \vspace{0.8cm}
      \centering
      Placeholder for observed data \(\widetilde{\mathbf d}\) with error bars, for example
      \(\pm 1\sigma\) from the data covariance \(\mathbf C_D\).
      \vspace{0.8cm}
    \end{minipage}
  }
  \caption{Observed data \(\widetilde{\mathbf d}\) with measurement uncertainty.}
  \label{fig:act7-observed-data}
\end{figure}

At this stage, nothing strange has happened. The setup is intentionally ordinary.

\begin{checkpointbox}
The example begins like a normal SOLA calculation: finite noisy data, known sensitivity kernels,
and localized averaging targets.
\end{checkpointbox}

\section{Sensitivity kernels and target kernels}

The sensitivity kernels \(K_i\) are generated from smooth synthetic functions, chosen to resemble
the kind of structured kernels one might see in a simplified physical problem. They are not
random spikes. They are smooth, correlated, and visually plausible.

The target kernels \(\mathcal T^{(k)}\) are localized bumps. Each asks for an average of the model
near a chosen location. If one could apply \(\Tau\) directly to the true model, the target property
vector would be
\[
  \bar{\mathbf p}
  :=
  \Tau(\bar m),
\]
with components
\[
  [\bar{\mathbf p}]_k
  =
  \langle \mathcal T^{(k)},\bar m\rangle_{\M}.
\]

But we cannot apply \(\Tau\) directly to \(\bar m\), because \(\bar m\) is unknown. We only have
the data. So we construct a SOLA matrix
\[
  \mathbf X:\D\to\Pspace,
\]
and compute
\[
  \widetilde{\mathbf p}
  =
  \mathbf X\widetilde{\mathbf d}.
\]

The approximate property map is
\[
  A:=\mathbf XG:\M\to\Pspace.
\]
Its components are represented by resolving kernels
\[
  A^{(k)}
  =
  \sum_{i=1}^{N_d}[\mathbf X]_{ki}K_i,
  \qquad k=1,\dots,N_p.
\]

The SOLA design goal is
\[
  A^{(k)}\approx \mathcal T^{(k)}
\]
for each \(k\). In the example, this goal appears to be achieved quite well.

\section{A standard SOLA diagnostic}

A standard SOLA diagnostic is to overlay the resolving kernels on the target kernels. If the
resolving kernels closely match the targets, the calculation looks encouraging.

In this example, the overlays look good. The resolving kernels are localized near the intended
targets. They have broadly similar shape. They pass the usual visual sniff test.

\begin{figure}[ht]
  \centering
  \fbox{
    \begin{minipage}{0.82\textwidth}
      \vspace{0.8cm}
      \centering
      Placeholder for resolving kernels \(A^{(k)}\) overlaid on target kernels
      \(\mathcal T^{(k)}\). The visual match should look good.
      \vspace{0.8cm}
    \end{minipage}
  }
  \caption{Resolving kernels \(A^{(k)}\) overlaid on target kernels
  \(\mathcal T^{(k)}\). The match appears visually encouraging.}
  \label{fig:act7-resolving-vs-target}
\end{figure}

The usual SOLA estimate is then plotted with propagated uncertainty:
\[
  \widetilde{\mathbf p}
  =
  \mathbf X\widetilde{\mathbf d},
\]
and
\[
  \mathbf C_P
  =
  \mathbf X\mathbf C_D\mathbf X^*.
\]
One may display
\[
  [\widetilde{\mathbf p}]_k
  \pm
  2\sqrt{[\mathbf C_P]_{kk}},
  \qquad k=1,\dots,N_p.
\]

\begin{figure}[ht]
  \centering
  \fbox{
    \begin{minipage}{0.82\textwidth}
      \vspace{0.8cm}
      \centering
      Placeholder for SOLA estimates \([\widetilde{\mathbf p}]_k\) with
      \(\pm 2\sigma\) uncertainty bands from \(\mathbf C_P\).
      \vspace{0.8cm}
    \end{minipage}
  }
  \caption{SOLA property estimates with propagated uncertainty bands.}
  \label{fig:act7-sola-estimates}
\end{figure}

Nothing about this diagnostic screams disaster. The data look reasonable. The targets are
reasonable. The resolving kernels look reasonably close to the targets. The uncertainty bands are
computed by the standard covariance propagation formula.

\begin{checkpointbox}
At this point, a practitioner might reasonably believe the SOLA calculation has behaved well.
\end{checkpointbox}

And in one sense, it has.

That last sentence is important. The point of the example is not that SOLA made an algebraic
mistake. The point is that a standard interpretation of the output is about to overstep what the
method actually estimated.

% ------------------------------------------------------------
\chapter{The Surprise}
% ------------------------------------------------------------

\section{Comparing SOLA estimates with true target properties}

Now we reveal the hidden truth.

The SOLA estimates are
\[
  \widetilde{\mathbf p}
  =
  \mathbf X\widetilde{\mathbf d}.
\]
The true target properties are
\[
  \bar{\mathbf p}
  :=
  \Tau(\bar m).
\]
Their components are
\[
  [\bar{\mathbf p}]_k
  =
  \langle \mathcal T^{(k)},\bar m\rangle_{\M}.
\]

A natural diagnostic is to compare
\[
  [\widetilde{\mathbf p}]_k
\]
with
\[
  [\bar{\mathbf p}]_k.
\]
In a synthetic example, we are allowed to do this because we know \(\bar m\). In a real inverse
problem, of course, \(\bar m\) is unknown, which is exactly why this diagnostic is unavailable.

When we make this comparison, the result can be shocking.

\begin{figure}[ht]
  \centering
  \fbox{
    \begin{minipage}{0.82\textwidth}
      \vspace{0.8cm}
      \centering
      Placeholder for a comparison between the SOLA estimates
      \([\widetilde{\mathbf p}]_k\) and the true target properties
      \([\bar{\mathbf p}]_k=[\Tau(\bar m)]_k\). The plot should show that the true values may
      lie far outside the propagated uncertainty bands.
      \vspace{0.8cm}
    \end{minipage}
  }
  \caption{SOLA estimates compared with true target properties. The disagreement appears
  catastrophic if the SOLA outputs are interpreted as estimates of \(\Tau(\bar m)\).}
  \label{fig:act7-sola-vs-truth}
\end{figure}

In the constructed example, the true target properties may lie far outside the nominal
\(\pm2\sigma\) bands. In fact, the SOLA estimates can even have the opposite sign from the true
target properties.

So we see a plot that appears to say:

\[
  \text{SOLA failed badly.}
\]

But that is not quite the right conclusion.

\section{The apparent failure}

The apparent failure has three visible features.

First, the signs can disagree:
\[
  [\widetilde{\mathbf p}]_k
  \text{ and }
  [\bar{\mathbf p}]_k
\]
may have opposite signs for many \(k\).

Second, the differences can be much larger than the reported propagated uncertainty:
\[
  \left|
    [\bar{\mathbf p}]_k-[\widetilde{\mathbf p}]_k
  \right|
  \gg
  2\sqrt{[\mathbf C_P]_{kk}}.
\]

Third, the true target properties may not behave like draws from the displayed uncertainty
bands. The usual visual story,
\[
  \text{estimate plus error bar should roughly cover the truth},
\]
breaks down.

This looks alarming because the resolving kernels looked good. The diagnostic in
Figure~\ref{fig:act7-resolving-vs-target} seemed to say that the SOLA calculation was resolving
the desired averages. Yet the comparison with truth says otherwise.

\begin{warningbox}
A good-looking resolving-kernel plot does not, by itself, guarantee that
\(\Tau(\bar m)\) is close to \(A(\bar m)\). The true model still matters.
\end{warningbox}

The key is to ask what the SOLA estimate was actually estimating.

The SOLA output is
\[
  \widetilde{\mathbf p}
  =
  \mathbf X\widetilde{\mathbf d}.
\]
Ignoring noise for clarity,
\[
  \widetilde{\mathbf p}
  =
  \mathbf XG(\bar m)
  =
  A(\bar m).
\]
But the plotted truth is
\[
  \bar{\mathbf p}
  =
  \Tau(\bar m).
\]

So the comparison is not really
\[
  \text{estimate}
  \quad\text{versus}\quad
  \text{the same quantity}.
\]
It is
\[
  A(\bar m)
  \quad\text{versus}\quad
  \Tau(\bar m).
\]

If \(A\neq\Tau\), these are different questions asked of the true model.

\section{What actually failed}

The dramatic plot appears when the SOLA numbers are interpreted as target-kernel properties.

But SOLA did not claim, at the algebraic level, to have computed
\[
  \Tau(\bar m)
\]
unless
\[
  A=\Tau.
\]
It computed, up to observational noise,
\[
  A(\bar m).
\]

In components,
\[
  [A(\bar m)]_k
  =
  \langle A^{(k)},\bar m\rangle_{\M},
\]
where
\[
  A^{(k)}
  =
  \sum_{i=1}^{N_d}[\mathbf X]_{ki}K_i.
\]
The true target component is
\[
  [\Tau(\bar m)]_k
  =
  \langle \mathcal T^{(k)},\bar m\rangle_{\M}.
\]
The difference is
\[
  [\Tau(\bar m)]_k-[A(\bar m)]_k
  =
  \langle \mathcal T^{(k)}-A^{(k)},\bar m\rangle_{\M}.
\]

Even if
\[
  \mathcal T^{(k)}-A^{(k)}
\]
looks small, this inner product can be large if \(\bar m\) has a large component in the mismatch
direction.

\begin{warningbox}
The method estimated the resolving-kernel properties. The dramatic plot appears when those
numbers are interpreted as target-kernel properties.
\end{warningbox}

So the right conclusion is not:

\[
  \text{SOLA failed.}
\]

The more careful conclusion is:

\[
  \text{The target-kernel interpretation failed.}
\]

That is a different diagnosis. It lets us keep what SOLA does well while being honest about what
it does not guarantee.

\begin{lessonbox}
SOLA can be working correctly as a resolving-kernel estimator even when it fails badly as an
estimator of the exact target-kernel property.
\end{lessonbox}

The next chapter explains how this kind of discrepancy can be deliberately manufactured. The
mechanism is not mysterious. It is the old nullspace problem wearing a very convincing hat.

% ------------------------------------------------------------
\chapter{How to Manufacture the Discrepancy}
% ------------------------------------------------------------

\section{A normal model}

We begin with an ordinary-looking model \(m_n\in\M\). For example, one may choose a smooth
moderate-norm model such as
\[
  m_n(r)
  =
  \exp\left(
    -\left(\frac{r-0.5}{0.5}\right)^2
  \right)
  \sin(5\pi r)+r.
\]
This model is not designed to be pathological. It is just a convenient synthetic truth from which
we generate data.

Define the corresponding noiseless data
\[
  \bar{\mathbf d}
  :=
  G(m_n),
\]
and the corresponding target properties
\[
  \mathbf p_n
  :=
  \Tau(m_n).
\]

If \(m_n\) were the true model, then the comparison between the SOLA estimate and
\(\mathbf p_n\) might look much less dramatic. In many ordinary examples, the SOLA proxy and
the true model may be close enough for the target interpretation to be practically useful.

But the data do not know that \(m_n\) is the truth. They only know
\[
  \bar{\mathbf d}=G(m_n).
\]
Any other model with the same data is equally compatible with the noiseless observation.

\section{A pathological data-equivalent model}

Now construct a different model \(\bar m\in\M\) satisfying
\[
  G(\bar m)=\bar{\mathbf d},
\]
but with target properties deliberately chosen to be very different. For instance, we can demand
\[
  \Tau(\bar m)=-\mathbf p_n.
\]

Thus \(\bar m\) has the same noiseless data as the normal model \(m_n\), but opposite target
properties:
\[
  G(\bar m)=G(m_n),
  \qquad
  \Tau(\bar m)=-\Tau(m_n).
\]

If such a model exists, then the data alone cannot distinguish between \(m_n\) and \(\bar m\).
They produce exactly the same noiseless data:
\[
  G(m_n)=G(\bar m).
\]
But the target property map distinguishes them sharply:
\[
  \Tau(m_n)\neq\Tau(\bar m).
\]

This is the whole trick. We build a model perturbation that is invisible to the data but visible to
the target properties.

Let
\[
  h:=\bar m-m_n.
\]
Then
\[
  G(h)=G(\bar m)-G(m_n)=\mathbf 0.
\]
So
\[
  h\in\kernel(G).
\]
But
\[
  \Tau(h)=\Tau(\bar m)-\Tau(m_n)
  =
  -\mathbf p_n-\mathbf p_n
  =
  -2\mathbf p_n.
\]
Thus \(h\) is invisible to \(G\), but visible to \(\Tau\).

\begin{intuitionbox}
The discrepancy is manufactured by hiding a property-changing component in the data nullspace.
The data cannot see it, but the target kernels can.
\end{intuitionbox}

This construction is artificial, but the lesson is not. Whenever the data have a nullspace, and
the target property sees some of that nullspace, there can be data-equivalent models with
different target properties.

\section{The stacked operator}

The construction can be written compactly using a stacked operator.

Define
\[
  \mathcal G
  :=
  \begin{bmatrix}
    G\\
    \Tau
  \end{bmatrix}
  :
  \M\to \D\oplus\Pspace.
\]
This operator sends a model to both its data and its target properties:
\[
  \mathcal G(m)
  =
  \begin{bmatrix}
    G(m)\\
    \Tau(m)
  \end{bmatrix}.
\]

To construct the pathological model \(\bar m\), solve
\[
  \mathcal G\bar m
  =
  \begin{bmatrix}
    \bar{\mathbf d}\\
    -\mathbf p_n
  \end{bmatrix}.
\]
Equivalently,
\[
  G(\bar m)=\bar{\mathbf d},
  \qquad
  \Tau(\bar m)=-\mathbf p_n.
\]

This is a finite set of linear constraints on the function \(\bar m\). Since \(\M\) is
infinite-dimensional and the constraints are finite-dimensional, there is often enough freedom
to satisfy them, although the resulting model may have a large norm or an unnatural shape.

That last point matters. The pathological model is often pathological precisely because it
contains a large component in a data-invisible direction. It is compatible with the data, but it
may not be a model one would have expected physically.

\begin{warningbox}
The construction does not say that every true model is pathological. It says that data alone do
not rule out such directions unless additional model information is supplied.
\end{warningbox}

This distinction is important for the tone of the argument. We are not saying that SOLA will
usually fail catastrophically. We are saying that SOLA, by itself, does not exclude certain
catastrophic interpretations of the target property.

If one believes the true model is not like \(\bar m\), that belief is model-side information. It may
be correct, but it should not be silently smuggled into the interpretation.

\section{Invisible to data, visible to properties}

The key mechanism can be summarized in one line:
\[
  h\in\kernel(G)
  \quad\text{but}\quad
  \Tau(h)\neq\mathbf 0.
\]

If such an \(h\) exists, then for any model \(m\),
\[
  G(m+h)=G(m),
\]
but
\[
  \Tau(m+h)=\Tau(m)+\Tau(h).
\]
So the data remain unchanged while the target properties change.

This is not a numerical accident. It is a structural feature of finite-data inverse problems.

The resolving map
\[
  A=\mathbf XG
\]
cannot see \(h\), because
\[
  A(h)=\mathbf XG(h)=\mathbf X\mathbf 0=\mathbf 0.
\]
But the target map \(\Tau\) may see \(h\):
\[
  \Tau(h)\neq\mathbf 0.
\]
Therefore
\[
  A(m+h)=A(m),
\]
while
\[
  \Tau(m+h)\neq\Tau(m).
\]

This proves the core point. SOLA can be perfectly consistent as an estimator of \(A(\bar m)\),
while being arbitrarily wrong as an estimator of \(\Tau(\bar m)\) over unrestricted model classes.

Indeed, if there exists \(h\in\kernel(G)\) with \(\Tau(h)\neq\mathbf 0\), then for
\[
  m_\alpha=m+\alpha h,
\]
we have
\[
  G(m_\alpha)=G(m)
\]
for every scalar \(\alpha\), but
\[
  \Tau(m_\alpha)=\Tau(m)+\alpha\Tau(h).
\]
Thus the target property can be moved arbitrarily far while the data remain fixed.

\begin{intuitionbox}
The data closed one eye. The target property looked exactly where that eye was blind.
\end{intuitionbox}

The honest conclusion is therefore:

\[
  \text{SOLA estimates resolved properties.}
\]
If we want exact target properties, we need additional model-side information that rules out, or
at least controls, the dangerous invisible directions.

\begin{lessonbox}
The problem is not that SOLA constructs resolving kernels. That is its strength. The problem is
interpreting resolved quantities as exact target quantities without controlling the model
directions that separate them.
\end{lessonbox}


% ============================================================
% Act VIII
% ============================================================

\part{Act VIII: SOLA as an Inversion in Disguise}

% ------------------------------------------------------------
\chapter{The Proxy Model}
% ------------------------------------------------------------

\section{From data map to proxy model}

One of the standard ways to describe SOLA is to say that it avoids reconstructing the whole
model. This is a good description of the practical workflow. SOLA does not ask us to first build
a model estimate \(\hat m\) and then evaluate its properties. It goes directly from data to
properties:
\[
  \widetilde{\mathbf p}
  =
  \mathbf X_c\widetilde{\mathbf d}.
\]

But algebra has a habit of opening side doors.

For the constrained SOLA operator, the estimate can also be written in the form
\[
  \widetilde{\mathbf p}
  =
  \mathbf X_c\widetilde{\mathbf d}
  =
  \Tau(\hat m),
\]
for a particular model \(\hat m\in\M\). This \(\hat m\) is not the true model. It is a proxy model,
constructed from the data and from the same ingredients that define the SOLA operator.

So there are two equivalent ways to read the same number:
\[
  \widetilde{\mathbf p}
  =
  \mathbf X_c\widetilde{\mathbf d}
\]
as a direct data-to-property estimate, or
\[
  \widetilde{\mathbf p}
  =
  \Tau(\hat m)
\]
as the exact target property of a particular proxy model.

\begin{intuitionbox}
SOLA walks past the front door of model reconstruction, but the algebra leaves a side entrance
open.
\end{intuitionbox}

This does not make SOLA bad. It makes SOLA more interpretable. The proxy model shows what
kind of model-side representative is hiding behind the data-to-property map.

Let
\[
  \mathbf M:=GG^*+\mathbf C_D,
  \qquad
  \mathbf u:=\mathbf M^{-1}\mathbf v,
  \qquad
  \beta:=\langle \mathbf v,\mathbf u\rangle_{\D}.
\]
Recall that the constrained noise-aware SOLA operator is
\[
  \mathbf X_c
  =
  \Tau G^*\mathbf M^{-1}
  -
  \frac{\Tau G^*(\mathbf u)-\mathbf w}{\beta}
  \mathbf u^*.
\]
Apply this operator to the observed data:
\[
  \mathbf X_c\widetilde{\mathbf d}
  =
  \Tau G^*\mathbf M^{-1}\widetilde{\mathbf d}
  -
  \frac{\langle \mathbf u,\widetilde{\mathbf d}\rangle_{\D}}{\beta}
  \bigl(\Tau G^*(\mathbf u)-\mathbf w\bigr).
\]
If we can find a model \(f\in\M\) satisfying
\[
  G(f)=\mathbf v,
  \qquad
  \Tau(f)=\mathbf w,
\]
then the second term can be written inside \(\Tau\):
\[
  \Tau G^*(\mathbf u)-\mathbf w
  =
  \Tau G^*(\mathbf u)-\Tau(f)
  =
  \Tau\bigl(G^*(\mathbf u)-f\bigr).
\]
Therefore
\[
  \mathbf X_c\widetilde{\mathbf d}
  =
  \Tau\left(
    G^*\mathbf M^{-1}\widetilde{\mathbf d}
    -
    \frac{\langle \mathbf u,\widetilde{\mathbf d}\rangle_{\D}}{\beta}
    \bigl(G^*(\mathbf u)-f\bigr)
  \right).
\]
This motivates the definition
\[
  \hat m
  :=
  G^*\mathbf M^{-1}\widetilde{\mathbf d}
  -
  \frac{\langle \mathbf u,\widetilde{\mathbf d}\rangle_{\D}}{\beta}
  \bigl(G^*(\mathbf u)-f\bigr).
\]
Then
\[
  \widetilde{\mathbf p}
  =
  \mathbf X_c\widetilde{\mathbf d}
  =
  \Tau(\hat m).
\]

\begin{checkpointbox}
The constrained SOLA estimate can be interpreted as the exact target property of a proxy model
\(\hat m\).
\end{checkpointbox}

This statement should be handled carefully. It does not say
\[
  \hat m=\bar m.
\]
It only says that the SOLA property estimate can be factored through one particular model-space
representative.

\section{The proxy formula}

Let us inspect the proxy model:
\[
  \hat m
  =
  G^*\mathbf M^{-1}\widetilde{\mathbf d}
  -
  \frac{\langle \mathbf u,\widetilde{\mathbf d}\rangle_{\D}}{\beta}
  \bigl(G^*(\mathbf u)-f\bigr).
\]

It has two pieces.

The first piece is
\[
  G^*\mathbf M^{-1}\widetilde{\mathbf d}.
\]
This resembles a regularized or noise-weighted back-projection of the data into the model space.
It is the model-space object produced by taking the data, weighting them by
\(\mathbf M^{-1}\), and applying the adjoint \(G^*\).

The second piece is
\[
  -
  \frac{\langle \mathbf u,\widetilde{\mathbf d}\rangle_{\D}}{\beta}
  \bigl(G^*(\mathbf u)-f\bigr).
\]
This is the correction associated with the resolving constraint
\[
  \mathbf X_c\mathbf v=\mathbf w.
\]
The scalar
\[
  \frac{\langle \mathbf u,\widetilde{\mathbf d}\rangle_{\D}}{\beta}
\]
depends on the data, while the model-space direction
\[
  G^*(\mathbf u)-f
\]
depends on the constraint.

The function \(f\in\M\) is chosen so that
\[
  G(f)=\mathbf v,
  \qquad
  \Tau(f)=\mathbf w.
\]
For averaging constraints, this says that \(f\) has the desired data response to the constant
test and the desired property response. In the simplest average-preserving case, \(f\) plays the
role of a model-space representative that makes the constraint visible inside \(\Tau\).

\begin{remark}[Non-uniqueness of \(f\)]
The function \(f\) need not be unique. If \(f_1\) and \(f_2\) both satisfy
\[
  G(f_j)=\mathbf v,
  \qquad
  \Tau(f_j)=\mathbf w,
  \qquad j=1,2,
\]
then their difference lies in the joint nullspace
\[
  f_1-f_2\in\kernel(G)\cap\kernel(\Tau).
\]
Changing \(f\) by such a direction changes the proxy model, but not its data response or target
property. This is another reminder that the proxy model is not a unique physical truth.
\end{remark}

In the noise-free case, where
\[
  \mathbf C_D=\mathbf 0
  \qquad\text{and hence}\qquad
  \mathbf M=GG^*,
\]
we have
\[
  G(G^*\mathbf u)
  =
  GG^*(GG^*)^{-1}\mathbf v
  =
  \mathbf v.
\]
Since \(G(f)=\mathbf v\), the correction direction satisfies
\[
  G\bigl(G^*(\mathbf u)-f\bigr)=\mathbf 0.
\]
Thus, in the noise-free constrained case, the correction lives in the nullspace of \(G\).

In the noise-aware case, the situation is slightly different. Since
\[
  \mathbf M=GG^*+\mathbf C_D,
\]
we have
\[
  GG^*\mathbf u
  =
  \mathbf v-\mathbf C_D\mathbf u.
\]
Therefore
\[
  G\bigl(G^*(\mathbf u)-f\bigr)
  =
  -\mathbf C_D\mathbf u.
\]
So the correction is not generally an exact nullspace correction when data noise is included.
Rather, it is a noise-weighted correction tied to the SOLA objective.

\begin{warningbox}
In the noise-free constrained case, the proxy correction is a true nullspace correction. In the
noise-aware case, it is better interpreted as a correction relative to the noise-weighted SOLA
geometry.
\end{warningbox}

This distinction matters because it prevents us from saying something too strong. The proxy
model is a useful algebraic representative, not a magical recovery of the true model.

\section{The meaning of the proxy}

The proxy model \(\hat m\) is not claimed to be the true model. It is one model-space object
whose target property equals the SOLA estimate:
\[
  \Tau(\hat m)=\widetilde{\mathbf p}.
\]

This gives a second interpretation of SOLA:
\[
  \text{direct view:}
  \qquad
  \widetilde{\mathbf p}
  =
  \mathbf X_c\widetilde{\mathbf d},
\]
and
\[
  \text{proxy-model view:}
  \qquad
  \widetilde{\mathbf p}
  =
  \Tau(\hat m).
\]

The direct view is the standard SOLA story. It emphasizes that the method avoids a full model
reconstruction. The proxy-model view emphasizes that the same data-to-property map can be
factored through a particular model-space representative.

Both views are useful. The first explains why SOLA is computationally attractive. The second
explains why SOLA is not entirely separate from classical inversion.

If the true model \(\bar m\) is close to the proxy model \(\hat m\), in the geometry relevant to
\(\Tau\), then
\[
  \Tau(\bar m)\approx \Tau(\hat m)=\widetilde{\mathbf p}.
\]
But if \(\bar m\) differs from \(\hat m\) by a direction that is hard for the data to see and easy
for \(\Tau\) to see, then
\[
  \Tau(\bar m)
\]
can be very different from
\[
  \Tau(\hat m).
\]

Thus the proxy interpretation does not remove the caution from the previous act. It sharpens it.

\begin{lessonbox}
SOLA may avoid explicitly reconstructing the model, but its algebra often carries an undercover
proxy reconstruction.
\end{lessonbox}

The proxy model is not the villain. It is a useful witness. It tells us what kind of model the SOLA
estimate is silently willing to be the property of.

% ------------------------------------------------------------
\chapter{Resolving Kernels and Resolution Operators}
% ------------------------------------------------------------

\section{Classical resolution operators}

In a classical linear inversion, one often builds a reconstruction operator
\[
  Y:\D\to\M.
\]
Given data \(\mathbf d\in\D\), the reconstructed model is
\[
  \hat m=Y\mathbf d.
\]
If the data are noise-free and generated by a true model \(m\), then
\[
  \mathbf d=G(m),
\]
so the reconstruction is
\[
  \hat m=YG(m).
\]
The composite map
\[
  \mathcal R:=YG:\M\to\M
\]
is called the resolution operator.

The resolution operator tells us what the inversion does to the true model:
\[
  m
  \longmapsto
  \mathcal R m.
\]
If
\[
  \mathcal R=I,
\]
then the inversion recovers every model exactly. In finite-data inverse problems, this almost
never happens. More commonly, \(\mathcal R\) is a smoothing, filtering, or projection-like
operator.

\begin{intuitionbox}
A resolution operator is the fingerprint left by the inverse method on the model. It tells us not
what model we wished to recover, but what version of the model the data and method can return.
\end{intuitionbox}

This is the classical resolution story:
\[
  \text{true model}
  \quad
  \xrightarrow{\;\mathcal R\;}
  \quad
  \text{resolved model}.
\]
The reconstructed model \(\hat m\) should therefore be interpreted as
\[
  \hat m=\mathcal R m
\]
in the ideal noise-free case, not automatically as \(m\).

\begin{warningbox}
Classical inversion does not usually recover the identity map. It recovers the model through a
resolution operator.
\end{warningbox}

This sounds familiar because SOLA has its own resolution object: the resolving kernel.

\section{SOLA resolving kernels}

SOLA defines the approximate property map
\[
  A:=\mathbf XG:\M\to\Pspace.
\]
For each component,
\[
  [A(m)]_k
  =
  \langle A^{(k)},m\rangle_{\M},
\]
where
\[
  A^{(k)}
  =
  \sum_{i=1}^{N_d}[\mathbf X]_{ki}K_i.
\]
The functions \(A^{(k)}\) are the resolving kernels.

Thus SOLA has the property-space resolution story
\[
  \text{true model}
  \quad
  \xrightarrow{\;A=\mathbf XG\;}
  \quad
  \text{resolved properties}.
\]

The target map is
\[
  \Tau:\M\to\Pspace.
\]
If SOLA achieved perfect resolution, then
\[
  A=\Tau.
\]
But in general,
\[
  A\neq\Tau.
\]
The resolving kernels describe the actual linear questions asked of the model.

Now connect this with the proxy-model view. Suppose there exists a linear reconstruction-like
operator
\[
  Y:\D\to\M
\]
such that
\[
  \mathbf X=\Tau Y.
\]
Then
\[
  A=\mathbf XG=\Tau YG.
\]
Using the classical resolution operator
\[
  \mathcal R:=YG,
\]
we obtain
\[
  A=\Tau\mathcal R.
\]

This identity is the bridge:
\[
  \boxed{
    \text{SOLA resolving map}
    =
    \text{target map after a model-space resolution operator}.
  }
\]

At the kernel level, if
\[
  [\Tau(m)]_k
  =
  \langle \mathcal T^{(k)},m\rangle_{\M},
\]
then
\[
  [A(m)]_k
  =
  [\Tau(\mathcal R m)]_k
  =
  \langle \mathcal T^{(k)},\mathcal R m\rangle_{\M}.
\]
Using the adjoint \(\mathcal R^*\),
\[
  [A(m)]_k
  =
  \langle \mathcal R^*\mathcal T^{(k)},m\rangle_{\M}.
\]
Therefore the resolving kernel is
\[
  A^{(k)}
  =
  \mathcal R^*\mathcal T^{(k)}.
\]

This is a revealing formula. It says that the SOLA resolving kernel can be seen as the target
kernel passed through the adjoint of a resolution operator.

\begin{checkpointbox}
Classical inversion modifies models through \(\mathcal R\). SOLA modifies target kernels through
\(\mathcal R^*\).
\end{checkpointbox}

This is not a criticism. It is a connection. SOLA gives direct access to resolving kernels, and that
is extremely useful. But those resolving kernels are not from another universe. They belong to
the same family of resolution ideas that appear in classical inversion.

\section{The shared resolution story}

We can now compare the two pictures.

Classical inversion says:
\[
  \mathbf d=G(m),
  \qquad
  \hat m=Y\mathbf d,
  \qquad
  \hat m=YG(m)=\mathcal R m.
\]
The method does not return \(m\). It returns \(\mathcal R m\).

SOLA says:
\[
  \mathbf d=G(m),
  \qquad
  \widetilde{\mathbf p}=\mathbf X\mathbf d,
  \qquad
  \widetilde{\mathbf p}=\mathbf XG(m)=A(m).
\]
The method does not necessarily return \(\Tau(m)\). It returns \(A(m)\).

If
\[
  \mathbf X=\Tau Y,
\]
then
\[
  A=\Tau\mathcal R.
\]
So the SOLA estimate is the target property of the resolved model:
\[
  \widetilde{\mathbf p}
  =
  \Tau(\mathcal R m).
\]
Equivalently, it is the true model tested against resolved target kernels:
\[
  [\widetilde{\mathbf p}]_k
  =
  \langle \mathcal R^*\mathcal T^{(k)},m\rangle_{\M}.
\]

These are two sides of the same coin:
\[
  \Tau(\mathcal R m)
  \quad\text{and}\quad
  (\Tau\mathcal R)(m).
\]

\begin{intuitionbox}
Classical inversion blurs the model and then asks the target question. SOLA can be read as
blurring the target question and then asking it of the true model.
\end{intuitionbox}

This helps explain both the power and the limitation of SOLA.

The power is that SOLA lets us design the target question and inspect the blurred version that
the data can actually support. Instead of hiding resolution inside a reconstructed model, SOLA
puts the resolving kernels on the table.

The limitation is that this is still resolution-limited information. The data do not grant perfect
access to \(\Tau(m)\) just because we avoided reconstructing \(m\). They grant access to
\[
  A(m)=\mathbf XG(m),
\]
which is a resolved version of the desired property.

\begin{warningbox}
Avoiding an explicit model reconstruction does not avoid resolution. It only changes where the
resolution appears.
\end{warningbox}

This is why the usual slogan that SOLA gives ``resolution control'' should be interpreted
carefully. SOLA does give meaningful control over the resolving kernels. That is a real advantage.
But it does not give full control over resolution in the sense of making the resolving kernels equal
to any targets we choose. The sensitivity kernels, noise covariance, constraints, and data geometry
still decide what is reachable.

\begin{lessonbox}
SOLA gives resolution control, not resolution sovereignty.
\end{lessonbox}

The next act turns from resolution to uncertainty. Once again, the theme will be the same:
the formula is mathematically correct, but its interpretation needs careful labels.

```latex
% ============================================================
% Act IX
% ============================================================

\part{Act IX: What About Uncertainty?}

% ------------------------------------------------------------
\chapter{Propagating Noise Is Not Posterior Uncertainty}
% ------------------------------------------------------------

\section{The propagated covariance}

We have reached one of the most delicate points in the interpretation of SOLA.

In the noisy setting, the observed data are
\[
  \widetilde{\mathbf d}
  =
  G(\bar m)+\boldsymbol\eta,
  \qquad
  \boldsymbol\eta\sim N(\mathbf 0,\mathbf C_D),
\]
where \(\bar m\in\M\) is the unknown true model and \(\mathbf C_D\) is the data-noise covariance.

Given a SOLA matrix
\[
  \mathbf X:\D\to\Pspace,
\]
the SOLA estimate is
\[
  \widetilde{\mathbf p}
  =
  \mathbf X\widetilde{\mathbf d}.
\]
Substituting the noisy data model gives
\[
  \widetilde{\mathbf p}
  =
  \mathbf XG(\bar m)
  +
  \mathbf X\boldsymbol\eta.
\]
The first term is the resolved signal:
\[
  \mathbf XG(\bar m)=A(\bar m),
  \qquad
  A:=\mathbf XG.
\]
The second term is the propagated observational noise:
\[
  \mathbf X\boldsymbol\eta.
\]

Since
\[
  \boldsymbol\eta\sim N(\mathbf 0,\mathbf C_D),
\]
and \(\mathbf X\) is linear, we have
\[
  \mathbf X\boldsymbol\eta
  \sim
  N(\mathbf 0,\mathbf C_P),
\]
where
\[
  \mathbf C_P
  =
  \mathbf X\mathbf C_D\mathbf X^*.
\]

This is the usual SOLA propagated covariance.

\begin{checkpointbox}
The covariance
\[
  \mathbf C_P=\mathbf X\mathbf C_D\mathbf X^*
\]
is the covariance of the noise term \(\mathbf X\boldsymbol\eta\) in the SOLA output.
\end{checkpointbox}

In components, the diagonal entries
\[
  [\mathbf C_P]_{kk}
\]
give the propagated variances of the individual SOLA components. One often displays
\[
  [\widetilde{\mathbf p}]_k
  \pm
  2\sqrt{[\mathbf C_P]_{kk}}.
\]
This is a perfectly valid calculation of the noise propagated through the chosen linear map
\(\mathbf X\).

The formula is not the problem. The label attached to the formula is the dangerous part.

\section{What this covariance measures}

The covariance
\[
  \mathbf C_P
  =
  \mathbf X\mathbf C_D\mathbf X^*
\]
answers a precise question:

\[
  \text{if the data noise has covariance }\mathbf C_D,
  \text{ what covariance does that noise induce in }\mathbf X\widetilde{\mathbf d}?
\]

This is a linear propagation question. It is the same kind of calculation one would do in any
linear measurement system.

Let
\[
  \boldsymbol\eta
  \sim
  N(\mathbf 0,\mathbf C_D).
\]
Then
\[
  \mathbf X\boldsymbol\eta
\]
has mean
\[
  \mathbb E[\mathbf X\boldsymbol\eta]
  =
  \mathbf X\mathbb E[\boldsymbol\eta]
  =
  \mathbf 0,
\]
and covariance
\[
  \mathbb E
  \left[
    (\mathbf X\boldsymbol\eta)
    (\mathbf X\boldsymbol\eta)^*
  \right]
  =
  \mathbf X
  \mathbb E[\boldsymbol\eta\boldsymbol\eta^*]
  \mathbf X^*
  =
  \mathbf X\mathbf C_D\mathbf X^*.
\]

So if the only randomness under discussion is the observational noise, and if \(\mathbf X\) is
fixed, then \(\mathbf C_P\) is exactly the right covariance.

\begin{intuitionbox}
The matrix \(\mathbf C_P\) tells us how much data static survives after SOLA has mixed the
channels.
\end{intuitionbox}

This covariance is useful. It tells us whether the chosen weights are stable or noisy. It lets us
compare two SOLA designs that have similar resolving kernels but different noise amplification.
It helps display the uncertainty induced by measurement error.

In that sense, \(\mathbf C_P\) is an important part of SOLA.

But it is not the whole uncertainty story.

\section{What this covariance does not measure}

The propagated covariance
\[
  \mathbf C_P=\mathbf X\mathbf C_D\mathbf X^*
\]
does not measure uncertainty about the exact target property
\[
  \bar{\mathbf p}
  =
  \Tau(\bar m).
\]

It measures uncertainty in the SOLA output caused by observational noise, around the resolved
quantity
\[
  A(\bar m)
  =
  \mathbf XG(\bar m).
\]
In other words, the noisy SOLA output has the form
\[
  \widetilde{\mathbf p}
  =
  A(\bar m)+\mathbf X\boldsymbol\eta.
\]
The covariance \(\mathbf C_P\) describes the second term:
\[
  \mathbf X\boldsymbol\eta.
\]
It does not describe the difference between
\[
  A(\bar m)
\]
and
\[
  \Tau(\bar m).
\]

That difference is
\[
  \Tau(\bar m)-A(\bar m)
  =
  (\Tau-A)(\bar m).
\]
This is a resolution or approximation error, not a propagated data-noise error.

\begin{warningbox}
The propagated covariance measures uncertainty in the resolved SOLA quantity \(A(\bar m)\)
caused by noisy data. It does not, by itself, measure uncertainty in the exact target property
\(\Tau(\bar m)\).
\end{warningbox}

This matters most when there are data-invisible model directions. Suppose
\[
  h\in\kernel(G),
\]
so that
\[
  G(h)=\mathbf 0.
\]
Then
\[
  A(h)=\mathbf XG(h)=\mathbf 0.
\]
The SOLA map cannot see \(h\). But the target map may see it:
\[
  \Tau(h)\neq\mathbf 0.
\]

For any scalar \(\alpha\), the models
\[
  \bar m
  \quad\text{and}\quad
  \bar m+\alpha h
\]
produce the same noiseless data:
\[
  G(\bar m+\alpha h)=G(\bar m).
\]
They also produce the same resolved SOLA property:
\[
  A(\bar m+\alpha h)=A(\bar m).
\]
But their target properties differ:
\[
  \Tau(\bar m+\alpha h)
  =
  \Tau(\bar m)+\alpha\Tau(h).
\]

Thus the exact target property can move while the SOLA output and its propagated covariance
do not notice.

This is the heart of the distinction:

\[
  \text{data-noise uncertainty}
  \neq
  \text{model/property uncertainty}.
\]

The first is described by
\[
  \mathbf X\mathbf C_D\mathbf X^*.
\]
The second requires some statement about which models are plausible among those compatible
with the data.

\begin{intuitionbox}
Noise covariance tells us how much the data wiggle. It does not tell us what might be hiding in
directions the data cannot see.
\end{intuitionbox}

If we want uncertainty about
\[
  \Tau(\bar m),
\]
then we need information about \(\bar m\), not only information about the data noise. That
information may be deterministic, such as a constraint set
\[
  \bar m\in B\subset\M,
\]
or probabilistic, such as a prior measure on \(\M\). Without such model-side information, the
target property may remain underdetermined even when the observational noise is small.

\begin{lessonbox}
Noise uncertainty is not the same thing as model/property uncertainty.
\end{lessonbox}

The usual SOLA covariance is therefore best read as a propagated noise covariance for the
resolved quantities. It is not, by itself, a posterior covariance for the true target properties.

% ------------------------------------------------------------
\chapter{The Shifted-Noise Issue}
% ------------------------------------------------------------

\section{The shifted noise measure}

There is a deeper way to phrase the same caution.

In the additive Gaussian setting, the data model is
\[
  \widetilde{\mathbf d}
  =
  G(\bar m)+\boldsymbol\eta,
  \qquad
  \boldsymbol\eta\sim N(\mathbf 0,\mathbf C_D).
\]
Once the observation \(\widetilde{\mathbf d}\) has been made, one might consider a new random
data-space object
\[
  \mathbf z_{\mathrm{shift}}
  =
  \widetilde{\mathbf d}+\boldsymbol\xi,
  \qquad
  \boldsymbol\xi\sim N(\mathbf 0,\mathbf C_D).
\]
This has distribution
\[
  \mathbf z_{\mathrm{shift}}
  \sim
  N(\widetilde{\mathbf d},\mathbf C_D).
\]
It is the noise distribution shifted so that its mean is the observed data.

Let us call this the shifted-noise measure.

There is nothing mathematically illegal about this object. It is a perfectly well-defined
probability measure on data space. We can push it through \(\mathbf X\), obtaining
\[
  \mathbf X\mathbf z_{\mathrm{shift}}
  =
  \mathbf X\widetilde{\mathbf d}
  +
  \mathbf X\boldsymbol\xi.
\]
This has mean
\[
  \mathbf X\widetilde{\mathbf d}
\]
and covariance
\[
  \mathbf X\mathbf C_D\mathbf X^*.
\]

So the usual SOLA display can be read as pushing forward the shifted-noise measure through
\(\mathbf X\):
\[
  N(\widetilde{\mathbf d},\mathbf C_D)
  \xrightarrow{\;\mathbf X\;}
  N(\mathbf X\widetilde{\mathbf d},\mathbf X\mathbf C_D\mathbf X^*).
\]

Again, this is mathematically valid.

The question is not whether the pushforward exists. The question is what it means.

\begin{checkpointbox}
The shifted-noise measure is a probability measure centred at the observed data. Pushing it
through \(\mathbf X\) gives a probability measure centred at the SOLA estimate.
\end{checkpointbox}

\section{Why this is observation-centred}

The shifted-noise construction starts from the realized observation
\[
  \widetilde{\mathbf d}
\]
and then spreads probability mass around it:
\[
  \widetilde{\mathbf d}
  \quad\leadsto\quad
  N(\widetilde{\mathbf d},\mathbf C_D).
\]
This makes it an observation-centred construction.

It answers something like:

\[
  \text{what happens if I jitter the observed data by another draw from the noise law?}
\]

That is a legitimate computational question. But it is not the fundamental generative question
behind the inverse problem.

The generative question is not centred on the observation. It is centred on a possible model:

\[
  \text{if }m\text{ were the true model, what data should I expect to observe?}
\]

In the additive Gaussian case, the answer is
\[
  \mathbf d_{\mathrm{obs}}
  \mid m
  \sim
  N(G(m),\mathbf C_D).
\]
That is, the distribution of possible observations is centred at \(G(m)\), not at the realized
\(\widetilde{\mathbf d}\).

\begin{intuitionbox}
The shifted-noise measure says, ``wiggle the observation.'' The generative model says,
``choose a model, then ask what observations it could have produced.''
\end{intuitionbox}

This difference may sound philosophical, but it matters. In inverse problems, the object we are
trying to learn about is not the observation. It is the model, or a property of the model. Therefore
the primitive uncertainty object should describe how models generate observations.

The shifted-noise measure does not do that. It is a data-side object centred after the fact on the
observed data.

\section{The generative question}

The generative formulation begins with a model-data relation.

In a probabilistic setting, this relation is a likelihood kernel
\[
  L_{\mathrm{prob}}(m,\cdot),
\]
where, for each \(m\in\M\),
\[
  L_{\mathrm{prob}}(m,\cdot)
\]
is a probability measure on the data space \(\D\).

In the additive Gaussian case,
\[
  L_{\mathrm{prob}}(m,\cdot)
  =
  N(G(m),\mathbf C_D).
\]
Equivalently, for a measurable data set \(B\subseteq\D\),
\[
  L_{\mathrm{prob}}(m,B)
  =
  \mathbb P\bigl(G(m)+\boldsymbol\eta\in B\bigr),
  \qquad
  \boldsymbol\eta\sim N(\mathbf 0,\mathbf C_D).
\]

This object answers the clean question:
\[
  \text{if }m\text{ were true, how likely would different observations be?}
\]

In a set-valued setting, the analogous object is a set likelihood
\[
  L_{\mathrm{set}}(m)\subseteq\D.
\]
This answers:
\[
  \text{if }m\text{ were true, which observations would be admissible?}
\]

For example, with additive bounded noise,
\[
  L_{\mathrm{set}}(m)
  =
  G(m)+S_{\mathrm{noise}},
\]
where \(S_{\mathrm{noise}}\subseteq\D\) is the set of allowed noise values.

Thus the probabilistic and set-valued versions have the same architecture:
\[
  m
  \quad\longmapsto\quad
  \text{possible data under }m.
\]

\begin{lessonbox}
The model-data relation is generative. It starts from a possible model and describes what that
model could produce in data space.
\end{lessonbox}

This is the object that belongs at the foundation. Once we have it, we can condition, restrict,
push forward, or compare models. Without it, we risk giving the observation itself a role it
cannot carry.

\section{The limitation of shifted-noise thinking}

The shifted-noise measure
\[
  N(\widetilde{\mathbf d},\mathbf C_D)
\]
may be useful in special additive settings. For example, it can be a convenient way to visualize
data-space uncertainty or to propagate observational perturbations through a fixed linear map.

But it should not be confused with a posterior distribution on models, or with a complete
description of uncertainty in target properties.

It is not a posterior on \(\M\). It lives on \(\D\).

It is not a posterior on \(\Pspace\) after applying \(\mathbf X\). Its pushforward lives on
\(\Pspace\), but it describes the distribution obtained by jittering the observation and applying
\(\mathbf X\). That is not the same as the distribution of
\[
  \Tau(\bar m)
\]
given the observed data.

It also does not account for model directions \(h\) satisfying
\[
  G(h)=\mathbf 0,
  \qquad
  \Tau(h)\neq\mathbf 0.
\]
Such directions can change the target property without changing the noiseless data. A
data-centred jitter around \(\widetilde{\mathbf d}\) does not discover them.

\begin{warningbox}
A shifted-noise measure is not a substitute for a model-data relation. It lives around the
observation, while the inverse problem asks what models and properties are compatible with that
observation.
\end{warningbox}

The correct reading is therefore modest:

\[
  N(\widetilde{\mathbf d},\mathbf C_D)
  \xrightarrow{\;\mathbf X\;}
  N(\mathbf X\widetilde{\mathbf d},\mathbf X\mathbf C_D\mathbf X^*)
\]
is a valid pushforward of a data-centred measure.

But the full inverse question asks something different. It asks how the observation constrains
\[
  \bar m
\]
or
\[
  \Tau(\bar m).
\]
For that, one needs a generative model-data relation and, if making posterior claims, a prior or
some other model-side information.

\begin{lessonbox}
A covariance pushed through \(\mathbf X\) answers a precise linear propagation question. It does
not, by itself, answer the full inverse question.
\end{lessonbox}

The next act proposes the corresponding reinterpretation of SOLA. Rather than pushing forward
an observation-centred shifted-noise object, we should ask how SOLA transforms the model-data
relation into a model-property relation.


% ============================================================
% Act X
% ============================================================

\part{Act X: Generative SOLA}

% ------------------------------------------------------------
\chapter{What Should SOLA Push Forward?}
% ------------------------------------------------------------

\section{The probabilistic model-data relation}

In the previous act, we separated two ideas that are often allowed to blur together.

On the one hand, we have a data-side shifted-noise object:
\[
  N(\widetilde{\mathbf d},\mathbf C_D).
\]
This is a probability measure centred at the observed data. It can be pushed through the SOLA
matrix \(\mathbf X\), giving
\[
  N(\mathbf X\widetilde{\mathbf d},\mathbf X\mathbf C_D\mathbf X^*).
\]
That calculation is mathematically valid.

On the other hand, the inverse problem itself is not born from a probability cloud placed around
the observation. It is born from a model-data relation. It begins with a possible model and asks
what data that model could have produced.

In the probabilistic setting, this model-data relation is encoded by a likelihood kernel
\[
  L_{\mathrm{prob}}:\M\times\Sigma_{\D}\to[0,1].
\]
For each model \(m\in\M\),
\[
  L_{\mathrm{prob}}(m,\cdot)
\]
is a probability measure on the data space \(\D\). It describes the possible observations under
the hypothesis that \(m\) is the true model.

Thus, for a measurable set \(B\subseteq\D\),
\[
  L_{\mathrm{prob}}(m,B)
\]
is the probability of observing data in \(B\) if \(m\) were true.

In the familiar additive Gaussian case,
\[
  \widetilde{\mathbf d}
  =
  G(m)+\boldsymbol\eta,
  \qquad
  \boldsymbol\eta\sim N(\mathbf 0,\mathbf C_D),
\]
the likelihood kernel is
\[
  L_{\mathrm{prob}}(m,\cdot)
  =
  N(G(m),\mathbf C_D).
\]
So for each model \(m\), the distribution of possible data is centred at \(G(m)\).

\begin{intuitionbox}
The likelihood kernel is a machine that takes a model as input and outputs a probability law on
data space. It answers: if this model were true, what observations should we expect?
\end{intuitionbox}

This is the direction of explanation that matters:
\[
  m
  \quad\longmapsto\quad
  \text{possible data}.
\]
The observation \(\widetilde{\mathbf d}\) is then compared against this family of model-indexed
data distributions.

This is different from starting with \(\widetilde{\mathbf d}\) and placing a cloud around it. The
shifted-noise cloud is centred on what happened. The likelihood kernel describes what could
happen under each possible model.

\begin{checkpointbox}
The shifted-noise measure is observation-centred. The likelihood kernel is model-centred, or
generative.
\end{checkpointbox}

This distinction is not merely aesthetic. Once the likelihood is written as
\[
  m\mapsto L_{\mathrm{prob}}(m,\cdot),
\]
it can handle much more than additive Gaussian noise. For example, one may have
model-dependent noise,
\[
  L_{\mathrm{prob}}(m,\cdot)
  =
  N(G(m),\mathbf C_D(m)),
\]
or non-Gaussian observations, or data whose distribution is not naturally written as
``prediction plus additive error''.

The generative relation is the flexible object. The shifted-noise measure is only a special
data-centred construction.

\section{Pushing the relation through SOLA}

Now bring SOLA back into the story.

A SOLA matrix
\[
  \mathbf X:\D\to\Pspace
\]
transforms data-space vectors into property-space vectors:
\[
  \mathbf d
  \quad\longmapsto\quad
  \mathbf X\mathbf d.
\]

If each model \(m\) generates a probability measure on data space,
\[
  L_{\mathrm{prob}}(m,\cdot),
\]
then \(\mathbf X\) pushes that measure forward to a probability measure on property space:
\[
  \mathbf X_{\#}L_{\mathrm{prob}}(m,\cdot).
\]

Thus SOLA induces a new model-property relation:
\[
  m
  \quad\longmapsto\quad
  \mathbf X_{\#}L_{\mathrm{prob}}(m,\cdot).
\]

This is the probabilistic version of generative SOLA.

For each model \(m\), the measure
\[
  \mathbf X_{\#}L_{\mathrm{prob}}(m,\cdot)
\]
describes the possible SOLA outputs that could be observed if \(m\) were the true model.

In the additive Gaussian case,
\[
  L_{\mathrm{prob}}(m,\cdot)
  =
  N(G(m),\mathbf C_D).
\]
Pushing this through \(\mathbf X\) gives
\[
  \mathbf X_{\#}L_{\mathrm{prob}}(m,\cdot)
  =
  N(\mathbf XG(m),\mathbf X\mathbf C_D\mathbf X^*).
\]
Since
\[
  A:=\mathbf XG,
\]
this becomes
\[
  \mathbf X_{\#}L_{\mathrm{prob}}(m,\cdot)
  =
  N(A(m),\mathbf X\mathbf C_D\mathbf X^*).
\]

This is subtly but importantly different from the shifted-noise pushforward
\[
  N(\mathbf X\widetilde{\mathbf d},\mathbf X\mathbf C_D\mathbf X^*).
\]
The generative pushforward is indexed by the model:
\[
  m\mapsto N(A(m),\mathbf X\mathbf C_D\mathbf X^*).
\]
The shifted-noise pushforward is centred at the observed SOLA output:
\[
  N(\mathbf X\widetilde{\mathbf d},\mathbf X\mathbf C_D\mathbf X^*).
\]

\begin{warningbox}
The same covariance \(\mathbf X\mathbf C_D\mathbf X^*\) can appear in two different stories.
One story is generative and model-indexed. The other is observation-centred. The formulas may
look similar, but the interpretation is not the same.
\end{warningbox}

The generative SOLA relation answers:
\[
  \text{if }m\text{ were true, what SOLA outputs should we expect?}
\]
That is a clean question. It is the property-space analogue of the original observation model.

\begin{lessonbox}
SOLA pushes the likelihood kernel from data space to property space. It turns a model-data
relation into a model-resolved-property relation.
\end{lessonbox}

Notice the phrase ``resolved-property''. The mean of the pushed-forward Gaussian is
\[
  A(m)=\mathbf XG(m),
\]
not generally
\[
  \Tau(m).
\]
So even in the generative interpretation, SOLA naturally produces a relation between models and
resolved properties. The target property \(\Tau(m)\) enters only when
\[
  A\approx\Tau,
\]
or when additional assumptions justify replacing one by the other.

\section{The set-valued analogue}

There is an exactly parallel set-valued version.

Instead of assigning to each model a probability measure on data space, suppose the observation
model assigns to each model a set of admissible data:
\[
  L_{\mathrm{set}}(m)\subseteq\D.
\]
This means:
\[
  \text{if }m\text{ were true, the observation should lie in }L_{\mathrm{set}}(m).
\]

For example, in an additive bounded-noise model, one might write
\[
  L_{\mathrm{set}}(m)
  =
  G(m)+S_{\mathrm{noise}},
\]
where
\[
  S_{\mathrm{noise}}\subseteq\D
\]
is the set of allowed noise values.

Now apply SOLA. The finite-dimensional map
\[
  \mathbf X:\D\to\Pspace
\]
sends admissible data sets to admissible property-space sets:
\[
  \mathbf X(L_{\mathrm{set}}(m))
  :=
  \{\mathbf X\mathbf d:\mathbf d\in L_{\mathrm{set}}(m)\}
  \subseteq\Pspace.
\]
Thus SOLA induces the set-valued model-property relation
\[
  m
  \quad\longmapsto\quad
  \mathbf X(L_{\mathrm{set}}(m)).
\]

In the additive bounded-noise case,
\[
  L_{\mathrm{set}}(m)
  =
  G(m)+S_{\mathrm{noise}},
\]
so
\[
  \mathbf X(L_{\mathrm{set}}(m))
  =
  \mathbf XG(m)+\mathbf X(S_{\mathrm{noise}}).
\]
Using
\[
  A=\mathbf XG,
\]
this becomes
\[
  \mathbf X(L_{\mathrm{set}}(m))
  =
  A(m)+\mathbf X(S_{\mathrm{noise}}).
\]

Again, the centre is the resolved property
\[
  A(m),
\]
not the exact target property
\[
  \Tau(m).
\]

\begin{intuitionbox}
The probabilistic version pushes probability clouds through \(\mathbf X\). The set-valued version
pushes admissible data sets through \(\mathbf X\). In both cases, SOLA transports the
model-data relation into property space.
\end{intuitionbox}

This set-valued viewpoint is useful because it makes the architecture visible without requiring
probability. The fundamental object is still the relation
\[
  m
  \quad\longmapsto\quad
  \text{possible observations}.
\]
SOLA then transforms the possible observations into possible resolved-property outputs.

\begin{checkpointbox}
Whether uncertainty is probabilistic or set-valued, the generative direction is the same:
\[
  m
  \longmapsto
  \text{data possibilities}
  \longmapsto
  \text{SOLA-output possibilities}.
\]
\end{checkpointbox}

\section{The new interpretation}

We can now state the conceptual replacement.

The observation-centred picture says:
\[
  \text{shifted noise on data}
  \quad
  \xrightarrow{\;\mathbf X\;}
  \quad
  \text{property uncertainty}.
\]

The generative picture says:
\[
  \text{model-data relation}
  \quad
  \xrightarrow{\;\mathbf X\;}
  \quad
  \text{model-property relation}.
\]

The second picture is the one that belongs at the foundation.

In the probabilistic case:
\[
  m\mapsto L_{\mathrm{prob}}(m,\cdot)
\]
becomes
\[
  m\mapsto \mathbf X_{\#}L_{\mathrm{prob}}(m,\cdot).
\]

In the set-valued case:
\[
  m\mapsto L_{\mathrm{set}}(m)
\]
becomes
\[
  m\mapsto \mathbf X(L_{\mathrm{set}}(m)).
\]

This is generative SOLA.

It does not say that SOLA solves the whole inverse problem. It does not magically turn resolved
properties into exact target properties. It does not remove the need for model-side assumptions
when those assumptions are needed.

What it does is give a cleaner answer to the question:

\[
  \text{what kind of uncertainty object should SOLA act on?}
\]

It should act on the model-data relation.

\begin{lessonbox}
Generative SOLA does not push forward an observation-centred shifted noise object. It pushes
forward the model-data relation.
\end{lessonbox}

This interpretation also clarifies what SOLA is good at. SOLA is a map from data space to
property space. Therefore it is naturally suited to transform data-side consequences of a model
into property-side consequences of that same model.

In short:
\[
  \boxed{
    \text{SOLA reshuffles generative information.}
  }
\]

But if we want to infer the true target property
\[
  \Tau(\bar m),
\]
then we need one more ingredient: information about which models are plausible after seeing the
data.

That is where prior information belongs.

% ------------------------------------------------------------
\chapter{Where Priors Belong}
% ------------------------------------------------------------

\section{When resolving-kernel interpretation is enough}

Sometimes the resolving-kernel interpretation is exactly what we need.

Suppose the scientific goal is to report the resolved quantity
\[
  A(\bar m)=\mathbf XG(\bar m).
\]
Then SOLA is directly matched to the task. The resolving kernels \(A^{(k)}\) define the reported
linear functionals:
\[
  [A(\bar m)]_k
  =
  \langle A^{(k)},\bar m\rangle_{\M}.
\]
If the data are noisy, the SOLA output is
\[
  \widetilde{\mathbf p}
  =
  A(\bar m)+\mathbf X\boldsymbol\eta,
\]
and the propagated covariance
\[
  \mathbf C_P=\mathbf X\mathbf C_D\mathbf X^*
\]
describes the observational noise in those resolved quantities.

This is a perfectly coherent use of SOLA.

The result should simply be labelled correctly:
\[
  \text{this is a resolved average, contrast, coefficient, or diagnostic,}
\]
with resolving kernels shown and interpreted.

\begin{checkpointbox}
If the reported object is \(A(\bar m)\), then SOLA supplies both the estimate and the resolving
kernels that define its meaning.
\end{checkpointbox}

In this mode, target kernels are design tools. They guide the construction of \(A\), but the final
scientific object is the resolved property itself.

\section{When true target properties are desired}

A different goal is to make statements about the exact target property
\[
  \bar{\mathbf p}
  =
  \Tau(\bar m).
\]
This is a stronger claim.

The difference between the exact target property and the resolved property is
\[
  \Tau(\bar m)-A(\bar m)
  =
  (\Tau-A)(\bar m).
\]
As discussed earlier, kernel closeness alone does not give a model-independent bound on this
quantity. We need some control of \(\bar m\).

That control is prior information, in the broad sense. It may be deterministic or probabilistic.
It may be explicit or implicit. It may come from physics, previous experiments, smoothness
assumptions, boundedness assumptions, geological plausibility, or mathematical regularity.

The important thing is that it lives on the model side.

\begin{warningbox}
If the claim is about \(\Tau(\bar m)\), then assumptions about \(\bar m\) are unavoidable. The
only choice is whether those assumptions are stated explicitly or smuggled into the
interpretation.
\end{warningbox}

This is the point where SOLA, deterministic inversion, and Bayesian inversion should stop
arguing past each other. They answer different questions unless we carefully align their
assumptions.

SOLA gives resolved data-supported properties. To turn those into statements about exact target
properties, we need to say which unresolved model directions are plausible.

\section{Deterministic prior information}

In deterministic linear inversion, model-side information can be encoded as a constraint set
\[
  B\subset\M.
\]
For example, \(B\) might be a norm ball,
\[
  B=\{m\in\M:\|m-m_0\|_{\M}\le R\},
\]
or an ellipsoid,
\[
  B=\{m\in\M:\|C_0^{-1/2}(m-m_0)\|_{\M}\le R\},
\]
or a set defined by bounds, smoothness, monotonicity, or other physical restrictions.

The data define a feasible set. In the noiseless case,
\[
  F_{\mathbf d}
  :=
  \{m\in\M:G(m)=\mathbf d\}.
\]
With bounded noise, one might instead use
\[
  F_{\widetilde{\mathbf d}}
  :=
  \{m\in\M:\widetilde{\mathbf d}\in L_{\mathrm{set}}(m)\}.
\]

If we believe
\[
  \bar m\in B,
\]
then the admissible models after observing the data are
\[
  F_{\widetilde{\mathbf d}}\cap B.
\]
The corresponding attainable target properties are
\[
  \Tau(F_{\widetilde{\mathbf d}}\cap B)
  :=
  \{\Tau(m):m\in F_{\widetilde{\mathbf d}}\cap B\}
  \subseteq\Pspace.
\]

This set is a deterministic uncertainty statement about the exact target property. It says:
given the data and the model constraint \(B\), these are the target-property vectors still
possible.

\begin{intuitionbox}
Deterministic prior information fences off the model space. The data choose the compatible
region inside the fence. The property map projects that region into property space.
\end{intuitionbox}

This is different from the SOLA propagated covariance. The SOLA covariance describes noise
propagation through \(\mathbf X\). The deterministic property set describes the range of exact
target properties over admissible models.

Both can be useful, but they answer different questions.

\section{Probabilistic prior information}

In probabilistic linear inversion, model-side information is encoded by a prior measure
\[
  \mu_{\mathrm{prior}}
\]
on \(\M\). This prior describes which models are considered plausible before seeing the data.

Together with the likelihood kernel
\[
  L_{\mathrm{prob}}(m,\cdot),
\]
the prior defines a joint measure on model-data space. Informally,
\[
  m\sim\mu_{\mathrm{prior}},
  \qquad
  \widetilde{\mathbf d}\mid m\sim L_{\mathrm{prob}}(m,\cdot).
\]
After observing data, one obtains a posterior distribution on models:
\[
  \mu_{\mathrm{post}}(\cdot\mid \widetilde{\mathbf d}).
\]
The posterior distribution of the exact target property is then the pushforward
\[
  \Tau_{\#}\mu_{\mathrm{post}}(\cdot\mid \widetilde{\mathbf d}).
\]

This is the probabilistic answer to the question:
\[
  \text{given the data and the prior, what do we believe about }\Tau(\bar m)?
\]

In a linear Gaussian setting, this posterior property distribution may be Gaussian, with a mean
and covariance that can be computed explicitly. But the interpretation is very different from
the SOLA covariance
\[
  \mathbf X\mathbf C_D\mathbf X^*.
\]
The Bayesian covariance reflects both the data and the prior model uncertainty. The SOLA
covariance reflects propagated observational noise through a chosen data-to-property map.

\begin{warningbox}
A propagated noise covariance is not a posterior covariance unless additional assumptions make
that identification valid. The prior is not a decoration; it is what turns likelihood into posterior
belief.
\end{warningbox}

This also explains why ``no prior'' can be both attractive and dangerous. SOLA can avoid a
probabilistic prior if the goal is to report resolved quantities. But if the goal is a probabilistic
statement about exact target properties, then a prior, or something playing the role of a prior,
is needed.

\section{The principle}

The lesson is not that every SOLA calculation must become Bayesian. Nor is it that deterministic
constraints are always better than probabilistic priors. The lesson is about labels.

If we report
\[
  A(\bar m),
\]
we should show and interpret the resolving kernels.

If we report uncertainty from
\[
  \mathbf X\mathbf C_D\mathbf X^*,
\]
we should call it propagated data-noise uncertainty.

If we report statements about
\[
  \Tau(\bar m),
\]
we should state the model-side assumptions that make those statements possible.

Those assumptions may be deterministic:
\[
  \bar m\in B.
\]
They may be probabilistic:
\[
  m\sim\mu_{\mathrm{prior}}.
\]
They may be physical, geological, smoothness-based, or structural. But they should be visible.

\begin{lessonbox}
Put priors where beliefs lie. Do not hide them inside the interpretation of a resolving kernel.
\end{lessonbox}

This gives a fair role to each framework.

SOLA is excellent for designing and reporting resolved linear properties. Deterministic inversion
is excellent for asking what exact properties remain possible under explicit constraints. Bayesian
inversion is excellent for assigning posterior probabilities once a prior and likelihood are
specified.

The mistake is to ask one framework to secretly do the work of another without paying the
assumptions bill.

\begin{checkpointbox}
Resolved-property claims need resolving kernels. Exact-property claims need model-side
assumptions. Posterior claims need priors.
\end{checkpointbox}

% ============================================================
% Act XI
% ============================================================

\part{Act XI: What SOLA Is Good For}

% ------------------------------------------------------------
\chapter{A Fair Summary}
% ------------------------------------------------------------

\section{What SOLA does well}

We have taken the long route. We began with weighted integrals, built resolving kernels, added
calibration constraints, introduced noise, inspected proxy models, and separated propagated
noise from model uncertainty.

After all that, it is worth returning to a simple question:

\[
  \text{what is SOLA good for?}
\]

The answer is: quite a lot.

SOLA is excellent at constructing controlled linear combinations of data. Given data
\[
  \widetilde{\mathbf d}\in\D,
\]
it builds a finite-dimensional map
\[
  \mathbf X:\D\to\Pspace
\]
and reports
\[
  \widetilde{\mathbf p}
  =
  \mathbf X\widetilde{\mathbf d}.
\]
Because the map is linear, the output is easy to compute, easy to inspect, and easy to propagate
data noise through.

But the deeper strength of SOLA is not merely that it gives a formula. Its strength is that it
gives resolving kernels.

For each property component,
\[
  [\widetilde{\mathbf p}]_k
  =
  \sum_{i=1}^{N_d}[\mathbf X]_{ki}\widetilde d_i,
\]
and, in the noiseless idealization,
\[
  [\widetilde{\mathbf p}]_k
  =
  \langle A^{(k)},\bar m\rangle_{\M},
\]
where
\[
  A^{(k)}
  =
  \sum_{i=1}^{N_d}[\mathbf X]_{ki}K_i.
\]
Thus the method tells us not only the number it reports, but also the linear functional of the
model that produced that number.

\begin{lessonbox}
SOLA does not merely output estimates. It outputs estimates together with the kernels that
define what those estimates mean.
\end{lessonbox}

This is a real advantage. In many inverse methods, resolution appears only after one derives a
resolution operator or performs additional diagnostics. In SOLA, resolution is part of the main
construction. The resolving kernels are there from the beginning, waving little flags that say:
\[
  \text{this is what the data combination actually measured.}
\]

SOLA also gives a clear way to balance resolution against noise. The noise-aware objective
\[
  \left\|
    \Tau-\mathbf XG
  \right\|_{\HS}^2
  +
  \Tr(\mathbf X\mathbf C_D\mathbf X^*)
\]
contains two competing desires:

\[
  \text{make resolving kernels look like target kernels,}
\]
and
\[
  \text{avoid amplifying data noise too much.}
\]

This trade-off is not hidden. It can be inspected, tuned, plotted, and discussed.

SOLA is also property-focused. It does not demand that we reconstruct the whole model before
asking for the quantity we care about. If the scientific question is already a linear property of
the model, then SOLA tries to build that property directly from the data.

That is elegant. It is the mathematical equivalent of taking the stairs directly to the room you
need instead of touring the entire building.

\begin{checkpointbox}
SOLA is strongest when the desired output is a resolved linear property and when the resolving
kernels are treated as part of the reported result.
\end{checkpointbox}

A fair list of SOLA's strengths is therefore:

\begin{itemize}[leftmargin=*]
  \item it constructs explicit linear combinations of data;
  \item it gives direct access to resolving kernels;
  \item it makes the resolution-noise trade-off visible;
  \item it supports property-focused estimation without requiring a full model reconstruction;
  \item it provides useful diagnostics for interpreting what the data can and cannot resolve;
  \item it gives a clean propagated data-noise covariance
  \[
    \mathbf C_P=\mathbf X\mathbf C_D\mathbf X^*.
  \]
\end{itemize}

These are not small virtues. They are exactly why SOLA is useful.

\section{What SOLA should not be asked to do alone}

The caution is not that SOLA is weak. The caution is that SOLA is sometimes asked to carry
more meaning than its construction supplies.

SOLA constructs
\[
  A:=\mathbf XG.
\]
The target property map is
\[
  \Tau:\M\to\Pspace.
\]
The desired exact target property of the true model is
\[
  \bar{\mathbf p}
  =
  \Tau(\bar m).
\]
The noiseless resolved property produced by SOLA is
\[
  A(\bar m).
\]

These are the same only if
\[
  A=\Tau
\]
on the relevant model class. If not, then the difference is
\[
  \Tau(\bar m)-A(\bar m)
  =
  (\Tau-A)(\bar m).
\]
Thus a small resolving-target mismatch is not, by itself, a model-independent guarantee that
\[
  A(\bar m)\approx\Tau(\bar m).
\]
To make that implication, we need some control of \(\bar m\).

\begin{warningbox}
SOLA can estimate resolved properties without a prior. It cannot turn resolved properties into
exact target properties without some assumption about the model.
\end{warningbox}

Similarly, the propagated covariance
\[
  \mathbf C_P=\mathbf X\mathbf C_D\mathbf X^*
\]
is a covariance for propagated data noise. It tells us how the observational noise
\(\boldsymbol\eta\) moves through \(\mathbf X\). It does not describe all uncertainty about
\[
  \Tau(\bar m).
\]
It does not account for directions \(h\in\M\) such that
\[
  G(h)=\mathbf 0,
  \qquad
  \Tau(h)\neq\mathbf 0.
\]
Those directions can change the exact target property without changing the noiseless data.

So SOLA alone should not be asked to provide:

\begin{itemize}[leftmargin=*]
  \item exact target properties \(\Tau(\bar m)\) when \(A\neq\Tau\);
  \item full uncertainty about \(\Tau(\bar m)\);
  \item posterior uncertainty without a prior or an equivalent model-side assumption;
  \item safe ratios or contrasts without checking the relevant resolving kernels together;
  \item model-independent error bounds from kernel plots alone.
\end{itemize}

This is not a demotion. It is just bookkeeping. SOLA answers a precise question, and it answers
it beautifully:

\[
  \text{which resolved linear properties can I build from the available data?}
\]

If we want a stronger answer, such as
\[
  \text{what is the exact target property of the true model?}
\]
then we need stronger assumptions.

\begin{intuitionbox}
SOLA can tell us what the data can say after we choose a way of listening. It cannot tell us what
the data never heard.
\end{intuitionbox}

The safest language is therefore:

\[
  \widetilde{\mathbf p}
  \text{ estimates }
  A(\bar m)
\]
with propagated data-noise covariance
\[
  \mathbf C_P.
\]
If we also want to say that
\[
  A(\bar m)\approx\Tau(\bar m),
\]
then we should state what model-side assumptions make that approximation reliable.

\section{Ratios and multi-property diagnostics}

The same caution becomes sharper when SOLA outputs are combined after the fact.

Suppose we estimate two quantities, perhaps associated with two model components
\[
  m_1,m_2\in\M,
\]
or with two different physical fields. Let
\[
  \Tau_1:\M\to\mathbb R,
  \qquad
  \Tau_2:\M\to\mathbb R
\]
be the corresponding target property maps. Since their outputs are finite-dimensional, we write
their one-dimensional property values as vectors:
\[
  \mathbf p_1=\Tau_1(m_1)\in\mathbb R,
  \qquad
  \mathbf p_2=\Tau_2(m_2)\in\mathbb R.
\]
Their scalar components are
\[
  [\mathbf p_1]_1,
  \qquad
  [\mathbf p_2]_1.
\]

SOLA does not generally return these exact target properties. It returns resolved quantities
defined by approximate maps
\[
  A_1:=\mathbf X_1G_1,
  \qquad
  A_2:=\mathbf X_2G_2.
\]
Thus the reported quantities are, in the noiseless idealization,
\[
  A_1(m_1),
  \qquad
  A_2(m_2).
\]

If we form a ratio of the SOLA outputs, we are forming something like
\[
  \frac{[A_1(m_1)]_1}{[A_2(m_2)]_1}.
\]
But the intended target ratio may be
\[
  \frac{[\Tau_1(m_1)]_1}{[\Tau_2(m_2)]_1}.
\]
In general,
\[
  \frac{[A_1(m_1)]_1}{[A_2(m_2)]_1}
  \neq
  \frac{[\Tau_1(m_1)]_1}{[\Tau_2(m_2)]_1}.
\]

The issue is not only that each resolving kernel might differ from its target. The issue is that
the two resolving kernels may differ from each other in ways that matter to the ratio.

For example, suppose the two target kernels are meant to localize the two quantities in the same
region. If the resolving kernel for the numerator is narrow, but the resolving kernel for the
denominator is broad, then the ratio is mixing two different spatial averages. If one resolving
kernel has side lobes and the other does not, the ratio inherits that asymmetry.

\begin{warningbox}
A ratio of SOLA outputs is a ratio of resolved quantities. It is not automatically the ratio of the
corresponding target quantities.
\end{warningbox}

This does not mean ratios are forbidden. It means they require matched interpretation.

A safer ratio statement would be:

\[
  \frac{[A_1(m_1)]_1}{[A_2(m_2)]_1}
\]
is the ratio of two resolved quantities, where the numerator is resolved through \(A_1\) and the
denominator is resolved through \(A_2\). If \(A_1\) and \(A_2\) have sufficiently compatible
resolving kernels for the intended scientific question, then the ratio may be meaningful.

The key word is compatible.

\begin{intuitionbox}
A ratio is a duet. If the two instruments are playing in different rooms, the harmony is mostly
wishful thinking.
\end{intuitionbox}

For multi-property diagnostics, the resolving kernels should therefore be inspected jointly. It is
not enough to say that each individual kernel looks acceptable. We must ask whether the family
of resolving kernels supports the operation we perform on the outputs.

For differences, this means checking whether the two resolved quantities are comparable.

For ratios, this means checking whether numerator and denominator share the required
resolution.

For gradients or contrasts, this means checking whether the resolving kernels reproduce the
intended signed structure.

For maps, this means remembering that every point may have its own kernel.

\begin{lessonbox}
Multi-property diagnostics require multi-property resolution checks. Looking at each resolving
kernel alone may miss the mismatch that matters.
\end{lessonbox}

\section{The final slogan}

We can now compress the whole story into three instructions:

\[
  \boxed{
    \text{Ask with targets. Interpret with resolving kernels. Add priors when you want truth.}
  }
\]

The first phrase,
\[
  \text{Ask with targets,}
\]
keeps the constructive heart of SOLA. We choose target kernels because they express the
scientific questions we would like to ask.

The second phrase,
\[
  \text{Interpret with resolving kernels,}
\]
keeps the answer honest. The data combination actually measures the resolving-kernel property,
not necessarily the exact target-kernel property.

The third phrase,
\[
  \text{Add priors when you want truth,}
\]
marks the boundary between resolved data-supported quantities and claims about the exact
target property of the true model. If we want to say something about
\[
  \Tau(\bar m),
\]
rather than merely
\[
  A(\bar m),
\]
then we need assumptions about \(\bar m\).

Those assumptions may be deterministic:
\[
  \bar m\in B\subset\M.
\]
They may be probabilistic:
\[
  m\sim\mu_{\mathrm{prior}}.
\]
They may come from physics, geology, smoothness, boundedness, or earlier experiments. But
they should be visible.

\begin{lessonbox}
SOLA is not a bad inverse method. It is a good reshuffling method that deserves honest labels.
\end{lessonbox}

So my take on SOLA is not that it should be discarded. Far from it. SOLA is a neat and useful
method precisely because it focuses on properties, exposes resolving kernels, and makes the
resolution-noise trade-off visible.

But its outputs should be read with the right grammar.

A SOLA number is not just a number. It comes with a kernel.

A SOLA error bar is not all uncertainty. It is propagated data-noise uncertainty.

A target kernel is not a promise. It is a design request.

A resolving kernel is not decoration. It is the measurement.

And if we want to turn resolved measurements into exact claims about the true model, we must
say what we believe about the model.

\begin{checkpointbox}
The mature use of SOLA is neither blind trust nor rejection. It is careful labelling:
resolved quantity, resolving kernel, propagated noise, and explicit model assumptions when
exact target properties are claimed.
\end{checkpointbox}

% ============================================================
% Back matter
% ============================================================

\backmatter

% ============================================================
% Summary
% ============================================================

\chapter{Summary of the Argument}

\section{The construction}

The mathematical skeleton of SOLA has three main ingredients.

First, there is the forward map
\[
  G:\M\to\D,
\]
which sends a model \(m\in\M\) to finite-dimensional data
\[
  \mathbf d=G(m)\in\D\simeq\mathbb R^{N_d}.
\]
In components,
\[
  [G(m)]_i
  =
  \langle K_i,m\rangle_{\M},
  \qquad i=1,\dots,N_d.
\]
The functions \(K_i\) are the sensitivity kernels.

Second, there is the property map
\[
  \Tau:\M\to\Pspace,
\]
which sends a model to the finite-dimensional property vector
\[
  \mathbf p=\Tau(m)\in\Pspace\simeq\mathbb R^{N_p}.
\]
In components,
\[
  [\mathbf p]_k
  =
  [\Tau(m)]_k
  =
  \langle \mathcal T^{(k)},m\rangle_{\M},
  \qquad k=1,\dots,N_p.
\]
The functions \(\mathcal T^{(k)}\) are the target kernels.

Third, SOLA chooses a finite-dimensional map
\[
  \mathbf X:\D\to\Pspace
\]
so that
\[
  \mathbf XG\approx \Tau.
\]
This is the core construction.

The simplest noiseless version solves
\[
  \mathbf X_0
  =
  \argmin_{\mathbf X}
  \left\|
    \Tau-\mathbf XG
  \right\|_{\HS}^2.
\]
With noise and a resolving constraint, the basic problem becomes
\[
  \mathbf X_c
  =
  \argmin_{\mathbf X}
  \left\{
    \left\|
      \Tau-\mathbf XG
    \right\|_{\HS}^2
    +
    \Tr(\mathbf X\mathbf C_D\mathbf X^*)
  \right\}
  \quad
  \text{subject to}
  \quad
  \mathbf X\mathbf v=\mathbf w.
\]

Here
\[
  \mathbf C_D
\]
is the data-noise covariance, and
\[
  \mathbf X\mathbf v=\mathbf w
\]
is a resolving constraint. For averages, this constraint enforces unit mass of the resolving
kernels. For other properties, it encodes the calibration needed for the intended interpretation.

\begin{lessonbox}
The SOLA construction is simple at heart:
\[
  \text{choose }\mathbf X\text{ so that }\mathbf XG\text{ behaves like }\Tau.
\]
The details decide what ``behaves like'' means.
\end{lessonbox}

\section{The interpretation}

Once \(\mathbf X\) has been chosen, the approximate property map is
\[
  A:=\mathbf XG:\M\to\Pspace.
\]
This is the map that the data combination actually synthesizes.

The target map is
\[
  \Tau:\M\to\Pspace.
\]
The SOLA map is
\[
  A=\mathbf XG:\M\to\Pspace.
\]

For the true model \(\bar m\), the exact target property is
\[
  \bar{\mathbf p}
  =
  \Tau(\bar m).
\]
The noiseless SOLA-accessible property is
\[
  A(\bar m)
  =
  \mathbf XG(\bar m).
\]
Thus the central interpretive distinction is
\[
  \Tau(\bar m)
  \quad\text{versus}\quad
  A(\bar m).
\]

In components,
\[
  [\Tau(\bar m)]_k
  =
  \langle \mathcal T^{(k)},\bar m\rangle_{\M},
\]
whereas
\[
  [A(\bar m)]_k
  =
  \langle A^{(k)},\bar m\rangle_{\M},
\]
with resolving kernels
\[
  A^{(k)}
  =
  \sum_{i=1}^{N_d}
  [\mathbf X]_{ki}K_i.
\]

Therefore the SOLA output should be interpreted through the resolving kernels \(A^{(k)}\). The
target kernels \(\mathcal T^{(k)}\) guide the construction, but the resolving kernels define the
reported quantities.

\begin{checkpointbox}
The target kernel is the request. The resolving kernel is the answer.
\end{checkpointbox}

\section{The limitation}

The difference between the exact target property and the resolved SOLA property is
\[
  \Tau(\bar m)-A(\bar m)
  =
  (\Tau-A)(\bar m).
\]
Hence
\[
  \|\Tau(\bar m)-A(\bar m)\|_{\Pspace}
  \le
  \|\Tau-A\|_{\M\to\Pspace}\,\|\bar m\|_{\M}.
\]

This inequality is useful, but it also exposes the limitation. A small resolving-target mismatch
does not give a uniform property-error bound over all possible models. It becomes an error
bound only when the true model is controlled.

If we know, for example, that
\[
  \|\bar m\|_{\M}\le R,
\]
then
\[
  \|\Tau(\bar m)-A(\bar m)\|_{\Pspace}
  \le
  R\|\Tau-A\|_{\M\to\Pspace}.
\]
But without such a model-side restriction, the difference can be arbitrarily large whenever
\(\Tau-A\neq 0\).

The nullspace mechanism makes this especially clear. If
\[
  h\in\kernel(G)
  \qquad\text{but}\qquad
  \Tau(h)\neq\mathbf 0,
\]
then
\[
  A(h)=\mathbf XG(h)=\mathbf 0,
\]
while
\[
  \Tau(h)\neq\mathbf 0.
\]
Thus the exact target property can change in directions that the SOLA map cannot see.

\begin{warningbox}
Resolving kernels can be close to target kernels and still fail to control target-property error
over unrestricted model classes. The missing ingredient is control of the model.
\end{warningbox}

\section{The uncertainty distinction}

The usual SOLA propagated covariance is
\[
  \mathbf C_P
  =
  \mathbf X\mathbf C_D\mathbf X^*.
\]
This is the covariance of the propagated data-noise term
\[
  \mathbf X\boldsymbol\eta
\]
when
\[
  \boldsymbol\eta\sim N(\mathbf 0,\mathbf C_D).
\]

It is therefore a covariance for observational noise in the resolved SOLA quantity
\[
  A(\bar m).
\]
It is not, by itself, a posterior covariance for the exact target property
\[
  \Tau(\bar m).
\]

To obtain uncertainty about \(\Tau(\bar m)\), one needs model-side information. This may be a
deterministic constraint set
\[
  \bar m\in B\subset\M,
\]
or a probabilistic prior
\[
  m\sim\mu_{\mathrm{prior}}.
\]

\begin{lessonbox}
Propagated data-noise uncertainty is not the same thing as uncertainty about the exact target
property of the true model.
\end{lessonbox}

\section{The generative correction}

The shifted-noise interpretation begins with an observation-centred object such as
\[
  N(\widetilde{\mathbf d},\mathbf C_D),
\]
then pushes it through \(\mathbf X\):
\[
  N(\widetilde{\mathbf d},\mathbf C_D)
  \xrightarrow{\;\mathbf X\;}
  N(\mathbf X\widetilde{\mathbf d},\mathbf X\mathbf C_D\mathbf X^*).
\]
This pushforward is mathematically valid, but it is not the fundamental inverse-problem object.

The generative object is the model-data relation.

In the probabilistic case, this is a likelihood kernel
\[
  m\mapsto L_{\mathrm{prob}}(m,\cdot).
\]
SOLA pushes it forward to
\[
  m
  \mapsto
  \mathbf X_{\#}L_{\mathrm{prob}}(m,\cdot).
\]

In the set-valued case, the model-data relation is
\[
  m\mapsto L_{\mathrm{set}}(m)\subseteq\D.
\]
SOLA pushes it forward to
\[
  m
  \mapsto
  \mathbf X(L_{\mathrm{set}}(m))\subseteq\Pspace.
\]

Thus the generative interpretation is
\[
  \text{model-data relation}
  \xrightarrow{\;\mathbf X\;}
  \text{model-property relation}.
\]

\begin{lessonbox}
Generative SOLA reshuffles model-indexed information. It does not need to place the observed
data at the centre of the uncertainty story.
\end{lessonbox}

\section{The final principle}

The whole argument can be compressed into one line:
\[
  \boxed{
    \text{Ask with targets. Interpret with resolving kernels. Add priors when you want truth.}
  }
\]

Targets are design objects. Resolving kernels are interpretation objects. Priors or deterministic
constraints are what allow claims about exact target properties of the true model.

% ============================================================
% Appendix A
% ============================================================

\chapter{Appendix A: Algebra of SOLA Minimizers}

\section{Preliminaries}

This appendix collects the basic algebra for the SOLA minimizers used in the main text.

Let
\[
  G:\M\to\D,
  \qquad
  \Tau:\M\to\Pspace,
\]
with finite-dimensional spaces
\[
  \D\simeq\mathbb R^{N_d},
  \qquad
  \Pspace\simeq\mathbb R^{N_p}.
\]
Let
\[
  \mathbf X:\D\to\Pspace
\]
be the SOLA matrix.

The objective in the noiseless case is
\[
  J_0(\mathbf X)
  =
  \left\|
    \Tau-\mathbf XG
  \right\|_{\HS}^2.
\]

For compactness, define the data-space Gram operator
\[
  \mathbf H:=GG^*:\D\to\D.
\]
With the conventions of these notes, \(\mathbf H\) is represented by a finite-dimensional matrix.
In components, under standard choices of inner products,
\[
  [\mathbf H]_{ij}
  =
  \langle K_j,K_i\rangle_{\M}.
\]

For the noisy case, define
\[
  \mathbf M:=\mathbf H+\mathbf C_D=GG^*+\mathbf C_D.
\]

We assume the inverses appearing below exist. In practice, positive data covariance, regularized
systems, or pseudoinverses may be used when exact invertibility fails.

\section{Unconstrained noiseless derivation}

The unconstrained noiseless SOLA problem is
\[
  \mathbf X_0
  =
  \argmin_{\mathbf X}
  \left\|
    \Tau-\mathbf XG
  \right\|_{\HS}^2.
\]

Let
\[
  E(\mathbf X):=\Tau-\mathbf XG.
\]
Then
\[
  J_0(\mathbf X)
  =
  \langle E(\mathbf X),E(\mathbf X)\rangle_{\HS}.
\]
Perturb \(\mathbf X\) by \(\varepsilon\mathbf H_1\), where
\[
  \mathbf H_1:\D\to\Pspace
\]
is an arbitrary finite-dimensional perturbation. Then
\[
  E(\mathbf X+\varepsilon\mathbf H_1)
  =
  \Tau-(\mathbf X+\varepsilon\mathbf H_1)G
  =
  E(\mathbf X)-\varepsilon\mathbf H_1G.
\]
Therefore
\[
  J_0(\mathbf X+\varepsilon\mathbf H_1)
  =
  \|E(\mathbf X)-\varepsilon\mathbf H_1G\|_{\HS}^2.
\]
Differentiating at \(\varepsilon=0\),
\[
  \delta J_0(\mathbf X)[\mathbf H_1]
  =
  -2
  \left\langle
    E(\mathbf X),
    \mathbf H_1G
  \right\rangle_{\HS}.
\]
Using the adjoint relation,
\[
  \left\langle
    E(\mathbf X),
    \mathbf H_1G
  \right\rangle_{\HS}
  =
  \left\langle
    E(\mathbf X)G^*,
    \mathbf H_1
  \right\rangle_{\HS}.
\]
Thus
\[
  \delta J_0(\mathbf X)[\mathbf H_1]
  =
  -2
  \left\langle
    (\Tau-\mathbf XG)G^*,
    \mathbf H_1
  \right\rangle_{\HS}.
\]
Stationarity for every perturbation \(\mathbf H_1\) gives
\[
  (\Tau-\mathbf XG)G^*=0.
\]
Hence
\[
  \Tau G^*-\mathbf XGG^*=0,
\]
or
\[
  \mathbf X\mathbf H=\Tau G^*.
\]
If \(\mathbf H\) is invertible, then
\[
  \boxed{
    \mathbf X_0=\Tau G^*\mathbf H^{-1}.
  }
\]
Since
\[
  \mathbf H=GG^*,
\]
this is often written as
\[
  \mathbf X_0
  =
  \Tau G^*(GG^*)^{-1}.
\]

\begin{checkpointbox}
The unconstrained noiseless minimizer is obtained by solving the normal equation
\[
  \mathbf XGG^*=\Tau G^*.
\]
\end{checkpointbox}

\section{Constrained noiseless derivation}

Now impose a resolving constraint
\[
  \mathbf X\mathbf v=\mathbf w,
\]
where
\[
  \mathbf v\in\D,
  \qquad
  \mathbf w\in\Pspace.
\]

The constrained noiseless problem is
\[
  \mathbf X_c
  =
  \argmin_{\mathbf X}
  \left\|
    \Tau-\mathbf XG
  \right\|_{\HS}^2
  \quad
  \text{subject to}
  \quad
  \mathbf X\mathbf v=\mathbf w.
\]

Introduce a Lagrange multiplier
\[
  \boldsymbol\lambda\in\Pspace.
\]
Consider the Lagrangian
\[
  \mathcal L(\mathbf X,\boldsymbol\lambda)
  =
  \left\|
    \Tau-\mathbf XG
  \right\|_{\HS}^2
  +
  2
  \left\langle
    \boldsymbol\lambda,
    \mathbf X\mathbf v-\mathbf w
  \right\rangle_{\Pspace}.
\]

The first variation with respect to \(\mathbf X\) is
\[
  \delta_{\mathbf X}\mathcal L
  =
  -2
  \left\langle
    (\Tau-\mathbf XG)G^*,
    \mathbf H_1
  \right\rangle_{\HS}
  +
  2
  \left\langle
    \boldsymbol\lambda,
    \mathbf H_1\mathbf v
  \right\rangle_{\Pspace}.
\]
The second term can be written as
\[
  \left\langle
    \boldsymbol\lambda,
    \mathbf H_1\mathbf v
  \right\rangle_{\Pspace}
  =
  \left\langle
    \boldsymbol\lambda\mathbf v^*,
    \mathbf H_1
  \right\rangle_{\HS}.
\]
Therefore stationarity gives
\[
  -(\Tau-\mathbf XG)G^*
  +
  \boldsymbol\lambda\mathbf v^*
  =
  0.
\]
Equivalently,
\[
  (\Tau-\mathbf XG)G^*
  =
  \boldsymbol\lambda\mathbf v^*.
\]
Rearranging,
\[
  \mathbf XGG^*
  =
  \Tau G^*-\boldsymbol\lambda\mathbf v^*.
\]
Using \(\mathbf H=GG^*\),
\[
  \mathbf X\mathbf H
  =
  \Tau G^*-\boldsymbol\lambda\mathbf v^*.
\]
Hence
\[
  \mathbf X
  =
  \Tau G^*\mathbf H^{-1}
  -
  \boldsymbol\lambda\mathbf v^*\mathbf H^{-1}.
\]

Define
\[
  \mathbf X_0:=\Tau G^*\mathbf H^{-1},
\]
and
\[
  \mathbf u:=\mathbf H^{-1}\mathbf v.
\]
Since \(\mathbf H\) is self-adjoint,
\[
  \mathbf v^*\mathbf H^{-1}=\mathbf u^*.
\]
Therefore
\[
  \mathbf X
  =
  \mathbf X_0-\boldsymbol\lambda\mathbf u^*.
\]

Now impose the constraint:
\[
  \mathbf X\mathbf v=\mathbf w.
\]
Substituting,
\[
  (\mathbf X_0-\boldsymbol\lambda\mathbf u^*)\mathbf v
  =
  \mathbf w.
\]
Thus
\[
  \mathbf X_0\mathbf v
  -
  \boldsymbol\lambda\,\mathbf u^*\mathbf v
  =
  \mathbf w.
\]
Define
\[
  \beta
  :=
  \mathbf u^*\mathbf v
  =
  \langle \mathbf u,\mathbf v\rangle_{\D}
  =
  \langle \mathbf H^{-1}\mathbf v,\mathbf v\rangle_{\D}.
\]
Then
\[
  \boldsymbol\lambda
  =
  \frac{\mathbf X_0\mathbf v-\mathbf w}{\beta}.
\]

Therefore
\[
  \boxed{
  \mathbf X_c
  =
  \mathbf X_0
  -
  \frac{\mathbf X_0\mathbf v-\mathbf w}{\beta}
  \mathbf u^*
  }
\]
with
\[
  \mathbf X_0=\Tau G^*\mathbf H^{-1},
  \qquad
  \mathbf u=\mathbf H^{-1}\mathbf v,
  \qquad
  \beta=\langle \mathbf v,\mathbf u\rangle_{\D}.
\]

\begin{lessonbox}
The constrained noiseless solution is the unconstrained solution plus a rank-one correction that
enforces the resolving constraint.
\end{lessonbox}

\section{Unconstrained noisy derivation}

The noise-aware objective is
\[
  J(\mathbf X)
  =
  \left\|
    \Tau-\mathbf XG
  \right\|_{\HS}^2
  +
  \Tr(\mathbf X\mathbf C_D\mathbf X^*).
\]

The first variation of the first term is the same as before:
\[
  -2
  \left\langle
    (\Tau-\mathbf XG)G^*,
    \mathbf H_1
  \right\rangle_{\HS}.
\]

For the noise term, using the standard trace derivative,
\[
  \delta
  \Tr(\mathbf X\mathbf C_D\mathbf X^*)[\mathbf H_1]
  =
  2
  \left\langle
    \mathbf X\mathbf C_D,
    \mathbf H_1
  \right\rangle_{\HS},
\]
assuming \(\mathbf C_D\) is self-adjoint.

Thus stationarity gives
\[
  -(\Tau-\mathbf XG)G^*+\mathbf X\mathbf C_D=0.
\]
Equivalently,
\[
  \Tau G^*-\mathbf XGG^*
  =
  \mathbf X\mathbf C_D.
\]
Hence
\[
  \Tau G^*
  =
  \mathbf X(GG^*+\mathbf C_D).
\]

Define
\[
  \mathbf M:=GG^*+\mathbf C_D.
\]
Then
\[
  \mathbf X\mathbf M=\Tau G^*.
\]
If \(\mathbf M\) is invertible,
\[
  \boxed{
    \mathbf X_0=\Tau G^*\mathbf M^{-1}.
  }
\]

\begin{checkpointbox}
The noisy unconstrained solution is obtained from the noiseless solution by replacing
\(GG^*\) with \(GG^*+\mathbf C_D\).
\end{checkpointbox}

\section{Constrained noisy derivation}

Finally, impose the resolving constraint
\[
  \mathbf X\mathbf v=\mathbf w
\]
on the noisy objective.

The constrained noisy problem is
\[
  \mathbf X_c
  =
  \argmin_{\mathbf X}
  \left\{
    \left\|
      \Tau-\mathbf XG
    \right\|_{\HS}^2
    +
    \Tr(\mathbf X\mathbf C_D\mathbf X^*)
  \right\}
  \quad
  \text{subject to}
  \quad
  \mathbf X\mathbf v=\mathbf w.
\]

Use the Lagrangian
\[
  \mathcal L(\mathbf X,\boldsymbol\lambda)
  =
  \left\|
    \Tau-\mathbf XG
  \right\|_{\HS}^2
  +
  \Tr(\mathbf X\mathbf C_D\mathbf X^*)
  +
  2
  \left\langle
    \boldsymbol\lambda,
    \mathbf X\mathbf v-\mathbf w
  \right\rangle_{\Pspace}.
\]

The stationarity condition is
\[
  -(\Tau-\mathbf XG)G^*
  +
  \mathbf X\mathbf C_D
  +
  \boldsymbol\lambda\mathbf v^*
  =
  0.
\]
Rearranging,
\[
  \mathbf X(GG^*+\mathbf C_D)
  =
  \Tau G^*-\boldsymbol\lambda\mathbf v^*.
\]
Using
\[
  \mathbf M:=GG^*+\mathbf C_D,
\]
we obtain
\[
  \mathbf X\mathbf M
  =
  \Tau G^*-\boldsymbol\lambda\mathbf v^*.
\]
Thus
\[
  \mathbf X
  =
  \Tau G^*\mathbf M^{-1}
  -
  \boldsymbol\lambda\mathbf v^*\mathbf M^{-1}.
\]

Define
\[
  \mathbf X_0:=\Tau G^*\mathbf M^{-1},
  \qquad
  \mathbf u:=\mathbf M^{-1}\mathbf v.
\]
Since \(\mathbf M\) is self-adjoint,
\[
  \mathbf v^*\mathbf M^{-1}=\mathbf u^*.
\]
Therefore
\[
  \mathbf X
  =
  \mathbf X_0-\boldsymbol\lambda\mathbf u^*.
\]

Apply the constraint:
\[
  (\mathbf X_0-\boldsymbol\lambda\mathbf u^*)\mathbf v=\mathbf w.
\]
Hence
\[
  \mathbf X_0\mathbf v-\boldsymbol\lambda\,\mathbf u^*\mathbf v=\mathbf w.
\]
Define
\[
  \beta
  :=
  \mathbf u^*\mathbf v
  =
  \langle \mathbf u,\mathbf v\rangle_{\D}.
\]
Then
\[
  \boldsymbol\lambda
  =
  \frac{\mathbf X_0\mathbf v-\mathbf w}{\beta}.
\]

Since
\[
  \mathbf X_0\mathbf v
  =
  \Tau G^*\mathbf M^{-1}\mathbf v
  =
  \Tau G^*(\mathbf u),
\]
we can write
\[
  \boldsymbol\lambda
  =
  \frac{\Tau G^*(\mathbf u)-\mathbf w}{\beta}.
\]

Thus
\[
  \boxed{
  \mathbf X_c
  =
  \Tau G^*\mathbf M^{-1}
  -
  \frac{\Tau G^*(\mathbf u)-\mathbf w}{\beta}
  \mathbf u^*
  }
\]
where
\[
  \mathbf M=GG^*+\mathbf C_D,
  \qquad
  \mathbf u=\mathbf M^{-1}\mathbf v,
  \qquad
  \beta=\langle \mathbf v,\mathbf u\rangle_{\D}.
\]

\begin{lessonbox}
The constrained noisy solution is the noise-aware SOLA operator plus the correction needed to
enforce the resolving constraint.
\end{lessonbox}

% ============================================================
% Appendix B
% ============================================================

\chapter{Appendix B: Proxy Model Derivation}

\section{Factorization of the constrained operator}

This appendix explains why the constrained SOLA estimate can sometimes be written as the
target property of a proxy model.

Start from the constrained noisy SOLA operator
\[
  \mathbf X_c
  =
  \Tau G^*\mathbf M^{-1}
  -
  \frac{\Tau G^*(\mathbf u)-\mathbf w}{\beta}
  \mathbf u^*,
\]
with
\[
  \mathbf M=GG^*+\mathbf C_D,
  \qquad
  \mathbf u=\mathbf M^{-1}\mathbf v,
  \qquad
  \beta=\langle \mathbf v,\mathbf u\rangle_{\D}.
\]

Apply \(\mathbf X_c\) to the observed data \(\widetilde{\mathbf d}\):
\[
  \mathbf X_c\widetilde{\mathbf d}
  =
  \Tau G^*\mathbf M^{-1}\widetilde{\mathbf d}
  -
  \frac{\Tau G^*(\mathbf u)-\mathbf w}{\beta}
  \mathbf u^*\widetilde{\mathbf d}.
\]
Since
\[
  \mathbf u^*\widetilde{\mathbf d}
  =
  \langle \mathbf u,\widetilde{\mathbf d}\rangle_{\D},
\]
this is
\[
  \mathbf X_c\widetilde{\mathbf d}
  =
  \Tau G^*\mathbf M^{-1}\widetilde{\mathbf d}
  -
  \frac{\langle \mathbf u,\widetilde{\mathbf d}\rangle_{\D}}{\beta}
  \bigl(\Tau G^*(\mathbf u)-\mathbf w\bigr).
\]

Now suppose there exists \(f\in\M\) such that
\[
  G(f)=\mathbf v,
  \qquad
  \Tau(f)=\mathbf w.
\]
Then
\[
  \Tau G^*(\mathbf u)-\mathbf w
  =
  \Tau G^*(\mathbf u)-\Tau(f)
  =
  \Tau\bigl(G^*(\mathbf u)-f\bigr).
\]
Therefore
\[
  \mathbf X_c\widetilde{\mathbf d}
  =
  \Tau G^*\mathbf M^{-1}\widetilde{\mathbf d}
  -
  \frac{\langle \mathbf u,\widetilde{\mathbf d}\rangle_{\D}}{\beta}
  \Tau\bigl(G^*(\mathbf u)-f\bigr).
\]
By linearity of \(\Tau\),
\[
  \mathbf X_c\widetilde{\mathbf d}
  =
  \Tau\left(
    G^*\mathbf M^{-1}\widetilde{\mathbf d}
    -
    \frac{\langle \mathbf u,\widetilde{\mathbf d}\rangle_{\D}}{\beta}
    \bigl(G^*(\mathbf u)-f\bigr)
  \right).
\]

This gives the proxy model formula.

\section{Construction of the proxy model}

Define
\[
  \hat m
  :=
  G^*\mathbf M^{-1}\widetilde{\mathbf d}
  -
  \frac{\langle \mathbf u,\widetilde{\mathbf d}\rangle_{\D}}{\beta}
  \bigl(G^*(\mathbf u)-f\bigr).
\]
Then
\[
  \boxed{
    \mathbf X_c\widetilde{\mathbf d}=\Tau(\hat m).
  }
\]

Thus the constrained SOLA estimate can be interpreted as the exact target property of the proxy
model \(\hat m\).

The first part,
\[
  G^*\mathbf M^{-1}\widetilde{\mathbf d},
\]
is a noise-weighted back-projection of the data. The second part,
\[
  -
  \frac{\langle \mathbf u,\widetilde{\mathbf d}\rangle_{\D}}{\beta}
  \bigl(G^*(\mathbf u)-f\bigr),
\]
is the correction associated with the resolving constraint.

\begin{warningbox}
The proxy model \(\hat m\) is an algebraic representative. It is not automatically the true
model.
\end{warningbox}

The condition
\[
  G(f)=\mathbf v,
  \qquad
  \Tau(f)=\mathbf w
\]
means that \(f\) realizes both the data-side calibration vector and the property-side desired
response. Such an \(f\) exists if
\[
  (\mathbf v,\mathbf w)
  \in
  \range
  \begin{bmatrix}
    G\\
    \Tau
  \end{bmatrix}.
\]
That is, there must be at least one model whose forward data are \(\mathbf v\) and whose target
properties are \(\mathbf w\).

If \(f\) is not unique, the proxy model is not unique either. If \(f_1\) and \(f_2\) both satisfy
\[
  G(f_j)=\mathbf v,
  \qquad
  \Tau(f_j)=\mathbf w,
  \qquad j=1,2,
\]
then
\[
  f_1-f_2\in\kernel(G)\cap\kernel(\Tau).
\]
Changing \(f\) by such a direction changes \(\hat m\), but not the SOLA estimate
\[
  \Tau(\hat m).
\]

\begin{lessonbox}
The proxy model is useful for interpretation, but it should not be mistaken for a uniquely
recovered physical model.
\end{lessonbox}

\section{Connection with null-space components}

The proxy formula becomes especially transparent in the noiseless constrained case.

Set
\[
  \mathbf C_D=\mathbf 0,
  \qquad
  \mathbf M=GG^*.
\]
Then
\[
  \mathbf u=(GG^*)^{-1}\mathbf v.
\]
Since
\[
  \mathbf v=G(f),
\]
we have
\[
  G(G^*\mathbf u)
  =
  GG^*(GG^*)^{-1}\mathbf v
  =
  \mathbf v.
\]
Therefore
\[
  G\bigl(G^*\mathbf u-f\bigr)
  =
  G(G^*\mathbf u)-G(f)
  =
  \mathbf v-\mathbf v
  =
  \mathbf 0.
\]
So in the noiseless constrained case,
\[
  G^*\mathbf u-f\in\kernel(G).
\]

The proxy model is then
\[
  \hat m
  =
  G^*(GG^*)^{-1}\widetilde{\mathbf d}
  -
  \frac{\langle \mathbf u,\widetilde{\mathbf d}\rangle_{\D}}{\beta}
  \bigl(G^*\mathbf u-f\bigr).
\]
It is the minimum-norm style back-projection
\[
  G^*(GG^*)^{-1}\widetilde{\mathbf d}
\]
plus a nullspace correction chosen to enforce the resolving constraint at the property level.

This is why SOLA can be viewed as an inversion in disguise. It may not explicitly reconstruct a
model as its primary output, but its constrained estimate can often be represented as the target
property of a particular proxy model.

In the noisy case, the correction is not generally an exact \(G\)-nullspace direction, because
\[
  \mathbf M=GG^*+\mathbf C_D.
\]
Then
\[
  GG^*\mathbf u
  =
  \mathbf v-\mathbf C_D\mathbf u,
\]
and therefore
\[
  G(G^*\mathbf u-f)
  =
  -\mathbf C_D\mathbf u.
\]
So the correction belongs to the noise-weighted geometry of the problem rather than the exact
nullspace of \(G\).

\begin{checkpointbox}
Noise-free constrained SOLA has a clean nullspace-proxy interpretation. Noise-aware SOLA has
a related proxy interpretation, but the geometry is modified by \(\mathbf C_D\).
\end{checkpointbox}

% ============================================================
% Appendix C
% ============================================================

\chapter{Appendix C: Suggested Figures}

\section{Figures for Acts I--III}

\subsection*{Figure 1: The unknown model behind a curtain}

\begin{itemize}[leftmargin=*]
  \item Show a smooth unknown function \(m(r)\) on \([0,1]\).
  \item Partially obscure it with a translucent curtain or dashed style.
  \item Message: the model exists, but it is not observed directly.
\end{itemize}

\subsection*{Figure 2: Sensitivity kernels as listening devices}

\begin{itemize}[leftmargin=*]
  \item Plot several sensitivity kernels \(K_i(r)\) on the same axis.
  \item Use light transparent curves so their family structure is visible.
  \item Optional annotation:
  \[
    d_i=\int_\Omega K_i(r)m(r)\dd r.
  \]
  \item Message: each datum listens to the model through a different kernel.
\end{itemize}

\subsection*{Figure 3: Data as integrals}

\begin{itemize}[leftmargin=*]
  \item Show one sensitivity kernel \(K_i(r)\), the model \(m(r)\), and their product
  \(K_i(r)m(r)\).
  \item Shade the area under the product curve.
  \item Message: the datum is one number obtained by integrating the product.
\end{itemize}

\subsection*{Figure 4: The target kernel as the desired question}

\begin{itemize}[leftmargin=*]
  \item Plot a localized target kernel \(\mathcal T^{(1)}(r)\).
  \item Show it centred at some location \(r_0\).
  \item Annotate:
  \[
    [\mathbf p]_1=\int_\Omega \mathcal T^{(1)}(r)m(r)\dd r.
  \]
  \item Message: the target kernel encodes the property we wish the data could measure.
\end{itemize}

\subsection*{Figure 5: Weighted kernels becoming a resolving kernel}

\begin{itemize}[leftmargin=*]
  \item Show several kernels \(K_i\) with weights \(x_i\).
  \item Then show the weighted sum
  \[
    R_{\mathbf x}=\sum_i x_iK_i.
  \]
  \item Overlay the target kernel \(\mathcal T^{(1)}\).
  \item Message: choosing data weights is choosing a resolving kernel.
\end{itemize}

\subsection*{Figure 6: Same shape, wrong mass}

\begin{itemize}[leftmargin=*]
  \item Plot two similar-looking kernels with different total area.
  \item Shade their areas differently.
  \item Annotate:
  \[
    \int R_{\mathbf x}(r)\dd r\neq 1.
  \]
  \item Message: visual localization is not enough for an average.
\end{itemize}

\subsection*{Figure 7: Unimodularity correction}

\begin{itemize}[leftmargin=*]
  \item Show an unconstrained resolving kernel and the constrained resolving kernel.
  \item Indicate that the constrained version satisfies
  \[
    \int_\Omega A^{(k)}(r)\dd r=1.
  \]
  \item Message: the constraint calibrates the resolving kernel.
\end{itemize}

\section{Figures for Acts IV--VI}

\subsection*{Figure 8: Different target types}

\begin{itemize}[leftmargin=*]
  \item Show three target kernels:
  \begin{enumerate}[label=(\alph*)]
    \item a positive localized average;
    \item a derivative-like target with positive and negative lobes;
    \item an oscillatory basis coefficient target.
  \end{enumerate}
  \item Message: not all targets are averages, so not all constraints should be unit-mass
  constraints.
\end{itemize}

\subsection*{Figure 9: Resolving constraints as calibration tests}

\begin{itemize}[leftmargin=*]
  \item Show a test model \(q\), the data response \(G(q)=\mathbf v\), and the desired property
  response \(\Tau(q)=\mathbf w\).
  \item Then show the constraint
  \[
    \mathbf X\mathbf v=\mathbf w.
  \]
  \item Message: resolving constraints enforce correct behavior on selected test directions.
\end{itemize}

\subsection*{Figure 10: Resolution-noise trade-off slider}

\begin{itemize}[leftmargin=*]
  \item Make a schematic with a slider.
  \item Left side: broad stable resolving kernel with small error bars.
  \item Right side: sharp resolving kernel with large error bars.
  \item Message: sharper resolution can amplify noise.
\end{itemize}

\subsection*{Figure 11: Propagated covariance}

\begin{itemize}[leftmargin=*]
  \item Show data error bars being pushed through \(\mathbf X\) to property error bars.
  \item Annotate:
  \[
    \mathbf C_P=\mathbf X\mathbf C_D\mathbf X^*.
  \]
  \item Message: this is propagated data-noise uncertainty.
\end{itemize}

\subsection*{Figure 12: Target versus resolving kernel}

\begin{itemize}[leftmargin=*]
  \item Overlay \(\mathcal T^{(k)}\) and \(A^{(k)}\).
  \item Highlight their difference
  \[
    \mathcal T^{(k)}-A^{(k)}.
  \]
  \item Message: the mismatch direction is where interpretation can leak.
\end{itemize}

\subsection*{Figure 13: Every map point has its own kernel}

\begin{itemize}[leftmargin=*]
  \item Show a row of SOLA map values.
  \item Under each value, show its associated resolving kernel.
  \item Message: a SOLA map is a map of values plus a map of kernels.
\end{itemize}

\section{Figures for Acts VII--XI}

\subsection*{Figure 14: Synthetic cautionary example setup}

\begin{itemize}[leftmargin=*]
  \item Panel 1: sensitivity kernels \(K_i\).
  \item Panel 2: target kernels \(\mathcal T^{(k)}\).
  \item Panel 3: observed data \(\widetilde{\mathbf d}\) with \(\pm1\sigma\) bars.
  \item Message: the example begins as a normal SOLA calculation.
\end{itemize}

\subsection*{Figure 15: Good-looking resolving kernels}

\begin{itemize}[leftmargin=*]
  \item Overlay resolving kernels \(A^{(k)}\) on target kernels \(\mathcal T^{(k)}\).
  \item Make the match look visually persuasive.
  \item Message: standard diagnostics may look good.
\end{itemize}

\subsection*{Figure 16: The apparent failure}

\begin{itemize}[leftmargin=*]
  \item Plot SOLA estimates
  \[
    [\widetilde{\mathbf p}]_k
  \]
  with \(\pm2\sigma\) bands.
  \item Overlay true target properties
  \[
    [\bar{\mathbf p}]_k=[\Tau(\bar m)]_k.
  \]
  \item Message: the target-property interpretation can fail dramatically.
\end{itemize}

\subsection*{Figure 17: Nullspace cartoon}

\begin{itemize}[leftmargin=*]
  \item Show model space as a plane or room.
  \item Show data-visible directions and data-invisible directions.
  \item Mark a direction \(h\) with
  \[
    G(h)=\mathbf 0,
    \qquad
    \Tau(h)\neq\mathbf 0.
  \]
  \item Message: some directions are invisible to data but visible to target properties.
\end{itemize}

\subsection*{Figure 18: Proxy model diagram}

\begin{itemize}[leftmargin=*]
  \item Show two paths:
  \[
    \widetilde{\mathbf d}
    \xrightarrow{\;\mathbf X_c\;}
    \widetilde{\mathbf p}
  \]
  and
  \[
    \widetilde{\mathbf d}
    \longmapsto
    \hat m
    \xrightarrow{\;\Tau\;}
    \widetilde{\mathbf p}.
  \]
  \item Message: SOLA can be interpreted directly or through a proxy model.
\end{itemize}

\subsection*{Figure 19: Resolution operator versus resolving kernel}

\begin{itemize}[leftmargin=*]
  \item Show classical inversion:
  \[
    m\mapsto \mathcal Rm.
  \]
  \item Show SOLA:
  \[
    \mathcal T^{(k)}\mapsto A^{(k)}=\mathcal R^*\mathcal T^{(k)}.
  \]
  \item Message: resolution operators and resolving kernels are related views of the same
  filtering phenomenon.
\end{itemize}

\subsection*{Figure 20: Shifted noise versus likelihood}

\begin{itemize}[leftmargin=*]
  \item Left panel: a Gaussian cloud centred at \(\widetilde{\mathbf d}\).
  \item Right panel: a family of Gaussian clouds centred at \(G(m)\) for different \(m\).
  \item Message: shifted noise is observation-centred, likelihood is model-centred.
\end{itemize}

\subsection*{Figure 21: Generative SOLA}

\begin{itemize}[leftmargin=*]
  \item Show:
  \[
    m
    \mapsto
    L_{\mathrm{prob}}(m,\cdot)
    \xrightarrow{\;\mathbf X_{\#}\;}
    \mathbf X_{\#}L_{\mathrm{prob}}(m,\cdot).
  \]
  \item Also show the set-valued version:
  \[
    m
    \mapsto
    L_{\mathrm{set}}(m)
    \xrightarrow{\;\mathbf X\;}
    \mathbf X(L_{\mathrm{set}}(m)).
  \]
  \item Message: SOLA pushes model-data relations into model-property relations.
\end{itemize}

\subsection*{Figure 22: Final summary graphic}

\begin{itemize}[leftmargin=*]
  \item Three columns:
  \begin{enumerate}[label=(\arabic*)]
    \item Ask with targets:
    \[
      \mathcal T^{(k)}.
    \]
    \item Interpret with resolving kernels:
    \[
      A^{(k)}.
    \]
    \item Add priors when you want truth:
    \[
      \bar m\in B
      \quad\text{or}\quad
      m\sim\mu_{\mathrm{prior}}.
    \]
  \end{enumerate}
  \item Message: the whole document in one picture.
\end{itemize}

\end{document}

\end{document}